\documentclass[11pt,a4paper]{scrreprt}

% ── Encoding & Language ──────────────────────────────────────────────────────
\usepackage[T1]{fontenc}
\usepackage[utf8]{inputenc}
\usepackage[english]{babel}
\usepackage{microtype}

% ── Fonts: Palatino body, Helvetica sans ─────────────────────────────────────
\usepackage{palatino}
\usepackage[scaled=0.90]{helvet}
\usepackage{courier}

% ── Layout ───────────────────────────────────────────────────────────────────
\usepackage[top=2.8cm,bottom=2.8cm,left=3.2cm,right=2.8cm,headheight=32pt]{geometry}
\usepackage{parskip}
\setlength{\parskip}{5pt}
\setlength{\parindent}{0pt}

% ── Color & Graphics ─────────────────────────────────────────────────────────
\usepackage[dvipsnames,svgnames,table]{xcolor}
\usepackage{graphicx}
\usepackage{svg}

% ── Tables ───────────────────────────────────────────────────────────────────
\usepackage{booktabs}
\usepackage{tabularx}
\usepackage{array}

% ── Boxes ────────────────────────────────────────────────────────────────────
\usepackage[most]{tcolorbox}
\tcbuselibrary{breakable,skins}

% ── Lists ────────────────────────────────────────────────────────────────────
\usepackage{enumitem}
\setlist[itemize]{leftmargin=1.5em,itemsep=1pt,topsep=2pt,parsep=0pt}
\setlist[description]{leftmargin=2em,style=nextline,itemsep=6pt}

% ── Headers / Footers ────────────────────────────────────────────────────────
\usepackage{scrlayer-scrpage}

% ── Hyperlinks ───────────────────────────────────────────────────────────────
\usepackage[hidelinks]{hyperref}

% gantt

\usepackage{amssymb}

% Prevent boxes from splitting across two pages
\usepackage{needspace}

% ── Project Colors ───────────────────────────────────────────────────────────
\definecolor{gobcolor}{HTML}{1a4971}
\definecolor{dashcolor}{HTML}{0e7490}
\definecolor{aacolor}{HTML}{2d5a8c}
\definecolor{e360color}{HTML}{1a6b3c}
\definecolor{espcolor}{HTML}{7b5e2a}
\definecolor{fgcolor}{HTML}{7d3c98}
\definecolor{filcolor}{HTML}{b7500a}
\definecolor{comcolor}{HTML}{c0392b}
\definecolor{navycolor}{HTML}{1a2e4a}
\definecolor{bglight}{HTML}{f7f5f0}
\definecolor{bordercolor}{HTML}{e2ddd6}
\definecolor{mutedcolor}{HTML}{6b7280}

\newcommand{\fadelogo}{\includegraphics[height=0.85cm]{marca_ADE_v1_color.png}}

% ── KOMA fonts ───────────────────────────────────────────────────────────────
\addtokomafont{chapter}{\color{navycolor}\Huge\bfseries\sffamily}
\addtokomafont{section}{\color{navycolor}\LARGE\bfseries\sffamily}
\addtokomafont{subsection}{\color{navycolor}\large\bfseries\sffamily}
\setkomafont{chapterprefix}{\small\color{mutedcolor}\normalfont}

% ── Page Style ───────────────────────────────────────────────────────────────
\pagestyle{scrheadings}
\clearpairofpagestyles


\ihead{\small\color{mutedcolor}FADE Strategic Plan 2026--2030}
\ohead{\raisebox{-0.15cm}{\fadelogo}}
\cfoot{\small\color{mutedcolor}\thepage}
% Also apply to "plain" pages (chapter openings, etc.)
\renewcommand*{\chapterpagestyle}{scrheadings}
\KOMAoptions{headsepline=false,footsepline=false}


\RedeclareSectionCommand[
  beforeskip=1.5\baselineskip,
  afterskip=1.2\baselineskip
]{chapter}

\RedeclareSectionCommand[
  beforeskip=1.2\baselineskip,
  afterskip=0.8\baselineskip
]{section}
% ── tcolorbox styles ─────────────────────────────────────────────────────────
\tcbset{
  pabox/.style 2 args={
    breakable, enhanced, skin=enhanced,
    before={\needspace{10\baselineskip}},
    colback=white!98!#1, colframe=#1,
    leftrule=5pt, rightrule=0.5pt, toprule=0.5pt, bottomrule=0.5pt,
    arc=2pt, outer arc=2pt,
    left=10pt, right=10pt, top=8pt, bottom=8pt,
    fonttitle=\bfseries\small\sffamily\raggedright,
    title={#2},
    attach boxed title to top left={yshift=-0.25mm,xshift=4mm},
    % Wrap long titles within the page width, including title padding and offset.
    minipage boxed title={\linewidth-12mm},
    boxed title style={
      colback=#1, colframe=#1,
      arc=2pt, outer arc=2pt,
      left=4pt, right=4pt, top=2pt, bottom=2pt,
    },
    title style={color=white},
  },
  fieldbox/.style={
    breakable, enhanced,
    colback=white, colframe=bordercolor,
    leftrule=0.5pt, rightrule=0.5pt, toprule=0.5pt, bottomrule=0.5pt,
    arc=1pt, left=8pt, right=8pt, top=4pt, bottom=4pt,
  },
  riskbox/.style={
    unbreakable, enhanced,
    colback=red!3!white, colframe=red!40!orange,
    leftrule=4pt, rightrule=0.5pt, toprule=0.5pt, bottomrule=0.5pt,
    arc=2pt, left=8pt, right=8pt, top=5pt, bottom=5pt,
  },
  mitbox/.style={
    unbreakable, enhanced,
    colback=green!3!white, colframe=green!45!black,
    leftrule=4pt, rightrule=0.5pt, toprule=0.5pt, bottomrule=0.5pt,
    arc=2pt, left=8pt, right=8pt, top=5pt, bottom=5pt,
  },
  descbox/.style={
    breakable, enhanced,
    colback=bglight, colframe=bordercolor,
    arc=3pt, left=12pt, right=12pt, top=10pt, bottom=10pt,
  },
  highlightbox/.style={
    breakable, enhanced,
    colback=navycolor!8!white, colframe=navycolor!40!white,
    leftrule=4pt, rightrule=0.5pt, toprule=0.5pt, bottomrule=0.5pt,
    arc=2pt, left=10pt, right=10pt, top=8pt, bottom=8pt,
    fontupper=\itshape,
  },
}

% ── Custom commands ───────────────────────────────────────────────────────────
\newcommand{\fieldlabel}[1]{%
  {\footnotesize\sffamily\bfseries\color{mutedcolor}\MakeUppercase{#1}}\par\nopagebreak[4]
}
\newcommand{\fieldsep}{\vspace{3pt}\noindent\textcolor{bordercolor}{\rule{\linewidth}{0.4pt}}\vspace{3pt}\par}

\begin{document}

% ── COVER ──────────────────────────────────────────────────────────────────
\begin{titlepage}
\pagecolor{navycolor}
\color{white}
\centering
\vspace*{2cm}
\includesvg[width=5cm]{logo}

\vspace{2cm}
{\fontsize{30}{38}\selectfont\bfseries FADE Strategic Plan}

\vspace{0.6cm}
{\fontsize{24}{30}\selectfont 2026--2030}

\vspace{1.2cm}
{\large\itshape UPV School of Management}\\[0.3cm]
{\large Universitat Polit\`ecnica de Val\`encia}\\[0.3cm]
\vspace{1.2cm}
{\large 23/07/2026 - v1.0}

\vspace{3cm}
\begin{tabular}{cccc}
{\Huge\bfseries 8}  & {\Huge\bfseries 40} & {\Huge\bfseries 71}  \\[4pt]
{\small projects}  & {\small objectives}  & {\small action plans} \\
\end{tabular}

\vfill
{\small\color{white!60!navycolor}Institutional reference document \quad\textbullet\quad Universitat Polit\`ecnica de Val\`encia}
\end{titlepage}
\nopagecolor

\tableofcontents
\clearpage

\chapter*{Introduction}
\addcontentsline{toc}{chapter}{Introduction}
Within the framework of the European Higher Education Area, the School has built a diversified academic offering, characterised by the integration of business education with a solid technological foundation, consistent with its polytechnic environment. It currently offers the Bachelor's Degree in Business Administration and Management (ADE), the Bachelor's Degree in Public Administration and Management (GAP) and the Bachelor's Degree in Business and Digital Technology (GETD), as well as a consolidated portfolio of double degrees. Postgraduate education reinforces this orientation with master's programmes in administrative management, taxation, corporate communication, service innovation, and corporate responsibility and sustainability.

In the international and professional sphere, the School maintains an extensive collaboration network with more than 200 exchange agreements with universities worldwide and close ties with companies and institutions, facilitating the connection between academic training and professional reality.

\begin{tcolorbox}[highlightbox]
The Strategic Plan 2026--2030 defines FADE's framework for action. The plan is organised around two strategic lines, Academic Excellence and Impact, and is implemented through 8 strategic projects covering areas from institutional governance to teaching innovation, internationalisation and the student experience.
\end{tcolorbox}

\clearpage
\chapter*{Strategic Framework}
\addcontentsline{toc}{chapter}{Strategic Framework}

\Needspace{14\baselineskip}
\addsec{Mission}
\begin{tcolorbox}[descbox]
\small\itshape The mission of the UPV School of Management is to educate ethical, well-rounded, and highly competent professionals in business and management, generating and transferring relevant knowledge through innovative, practice-oriented education in a technological and interdisciplinary environment, and contributing to sustainable development.
\end{tcolorbox}

\Needspace{14\baselineskip}
\addsec{Vision}
\begin{tcolorbox}[fieldbox]
\small FADE aims to be recognised nationally and internationally for preparing leaders who can drive organisational transformation, adapt to a business environment shaped by continuous technological change, and create value for society.
\end{tcolorbox}

\Needspace{14\baselineskip}
\addsec{Institutional Statements}
\Needspace{10\baselineskip}
\subsection*{Statement 1}
\small FADE serves students, faculty, professional services staff, companies, public administrations, and society at large through high-quality, inclusive, and practice-oriented education with national and international recognition. Within this framework, it promotes talent and prepares its graduates to make a meaningful contribution to sustainable development.

\Needspace{10\baselineskip}
\subsection*{Statement 2}
\small FADE's bachelor's and master's programmes integrate education in business and management with the strong technological and innovative culture that characterises its polytechnic environment. Through practical and experiential learning methodologies, the School prepares students to address the challenges of a dynamic professional environment.

\Needspace{10\baselineskip}
\subsection*{Statement 3}
\small FADE prioritises applied research, complemented by basic research and educational innovation, in order to promote innovation and knowledge transfer. With a focus on technological leadership in education, its research supports advances that have a positive influence on its stakeholders.

\Needspace{14\baselineskip}
\addsec{Values}
\begin{description}
\item[\textbf{Excellence}]
\small We conduct our academic and management activities with rigour, demanding standards and continuous improvement, ensuring high-quality, internationally oriented education and the generation of relevant, validated knowledge.
\item[\textbf{Ethics}]
\small We act with responsibility, transparency, and public commitment, promoting equal opportunities and accountability. We manage resources efficiently and orient our activity toward sustainable development.
\item[\textbf{Innovation}]
\small We promote the integration of technology and the continuous updating of content and methodologies, encouraging new forms of teaching and knowledge production.
\item[\textbf{Professionalism}]
\small We promote practical education aligned with organisational realities, strengthening connections with companies and public administrations to prepare our students for demanding professional environments.
\item[\textbf{Collaboration}]
\small We encourage the active participation of all groups within the School and a multidisciplinary approach, promoting an open and cooperative academic environment enriched by diverse perspectives.
\item[\textbf{Polytechnic Identity}]
\small We are part of UPV and share its technological culture, its orientation toward the application of knowledge, and its commitment to serving society.
\end{description}

\Needspace{14\baselineskip}
\addsec{Strategic Lines}
\Needspace{10\baselineskip}
\subsection*{Academic Excellence}
\small Strengthening the quality of degree programmes, teaching innovation, quality assurance processes and academic monitoring of students. Includes regulatory adaptation, implementation of new degrees, re-accreditation and continuous improvement.

\Needspace{10\baselineskip}
\subsection*{Impact}
\small Institutional outreach, external visibility, connection with the professional and social environment, internationalisation and strengthening of the School's relational ecosystem with companies, administrations, alumni and other collaborating entities.

\Needspace{14\baselineskip}
\addsec{Cross-cutting Action Axes}
\begin{description}
\item[\textbf{Processes}]
\small Structuring and formalisation of institutional procedures, ensuring clear, auditable and sustainable frameworks for action.
\item[\textbf{Data}]
\small Generation of reliable, traceable and comparable information that supports evidence-based academic and institutional decision-making.
\item[\textbf{Communication}]
\small Effective transmission of decisions, results and initiatives to all relevant audiences, reinforcing the identity and visibility of FADE.

\item[\textbf{Artificial Intelligence}]
\small Progressive integration of artificial intelligence as a cross-cutting axis of institutional transformation, consistent with the \textit{Ten Dimension AI Readiness Framework} of the Digital Education Council (2025). The plan addresses AI readiness from multiple dimensions: as a tool supporting institutional analysis and decision-making (Governance, Dashboard), as curricular content and a teaching methodology (Agenda Avanza), as a mechanism for student support (University Experience 360) and as support for relationship and communication management (Impact Lab, COMUNICA). AI is to be incorporated progressively, focusing on specific use cases, under human oversight and in compliance with Regulation (EU) 2024/1689 on Artificial Intelligence, with the aim of moving from an emerging maturity level to an established level over the 2026--2030 period.



\end{description}

\Needspace{14\baselineskip}
\addsec{Logic of Change}

\small
The Strategic Plan 2026--2030 is not a collection of independent projects. It is a commitment to transforming the system that determines FADE's quality and impact: the chain connecting institutional culture with faculty development, teaching, students' competencies, their employability and, ultimately, the School's reputation and demand for its programmes.

This chain has an important characteristic: the effects of decisions taken today are not immediate. An investment in faculty development in 2026 translates into methodological changes in 2027, improvements in students' competencies in 2028 and greater recognition from businesses in 2029--2030. Recognising these time lags is essential to interpreting the plan's progress correctly and avoiding premature conclusions about its effectiveness.

The plan also identifies that the main constraint on transformation is neither technological nor budgetary. It is cultural and organisational: faculty incentives, the flexibility of curricula and the institution's capacity to share and scale up good practice. 

Finally, the plan recognises that organisational systems can generate both virtuous cycles and dynamics that block progress. The cycle we want to activate is: stronger AI competencies $\rightarrow$ better academic and professional outcomes $\rightarrow$ stronger reputation $\rightarrow$ greater capacity 
to attract talent and projects $\rightarrow$ greater investment in people and knowledge. The risk we want to interrupt is the opposite: uncertainty about change $\rightarrow$ resistance $\rightarrow$ superficial adoption $\rightarrow$ poor results $\rightarrow$ greater resistance. The mitigation strategy is not imposition but the creation of visible, early successes that reinforce the positive cycle.

\Needspace{14\baselineskip}
\addsec{Alignment with the AI Maturity Framework (DEC 2025)}

\small
The FADE Strategic Plan 2026--2030 aligns with the \textit{Ten Dimension AI Readiness Framework} of the Digital Education Council (2025), developed with the participation of 27 universities from 17 countries. The following table shows the plan's coverage of the framework's ten dimensions, identifying the projects contributing to each and the estimated maturity level at the start of the plan.

\vspace{6pt}

\begin{tabularx}{\textwidth}{>{\bfseries\small}p{4.2cm} >{\small}X >{\small\centering\arraybackslash}p{2.2cm}}
\toprule
\textbf{DEC Dimension} & \textbf{Coverage in the FADE Plan} & \textbf{Initial Maturity} \\
\midrule
1. Strategic alignment & Strategic framework, lines and cross-cutting axes. All projects. & Established \\
\addlinespace
2. Institutional governance & Governance project (S.O.1--S.O.5). & Developing \\
\addlinespace
3. Stakeholder engagement & University Experience 360, FADE Impact Lab, COMUNICA. & Developing \\
\addlinespace
4. Operational readiness & Academic Dashboard, FADE Spaces. & Emerging \\
\addlinespace
5. AI literacy and ethical use & Agenda Avanza (S.O.4), Governance (S.O.4). & Emerging \\
\addlinespace
6. Accessibility and inclusion & University Experience 360 (A.P.2.2). Partial coverage. & Emerging \\
\addlinespace
7. Faculty professional development & Agenda Avanza (S.O.4, A.P.4.1--A.P.4.2). & Emerging \\
\addlinespace
8. Teaching, learning and assessment & Agenda Avanza (S.O.1--S.O.5). & Developing \\
\addlinespace
9. Curriculum and alignment with the labour market & Agenda Avanza (S.O.2, S.O.5), FADE Impact Lab. & Developing \\
\addlinespace
10. Research and innovation & FADE Impact Lab (S.O.1--S.O.4), Agenda Avanza (S.O.4). & Emerging \\
\bottomrule
\end{tabularx}

\vspace{4pt}
{\footnotesize\color{mutedcolor}\textit{Maturity levels according to the DEC Framework: Emerging - Developing - Established - Mature.}}

\vspace{6pt}

\Needspace{14\baselineskip}
\addsec{Alignment with the UPV Strategy (SIRVE)}

\small
The FADE Strategic Plan 2026--2030 forms part of the institutional strategy of the Universitat Politècnica de València, SIRVE, organised around five objectives: UPV\_Sostenible, UPV\_Internacional, UPV\_Relevante, UPV\_Vital and UPV\_Excelente. The following table presents each strategic project's estimated contribution to each of these objectives, on a scale from 1 (low contribution) to 5 (very high contribution).

\vspace{6pt}

\begin{tabularx}{\textwidth}{>{\bfseries\small}X *{5}{>{\small\centering\arraybackslash}p{1.75cm}}}
\toprule
\textbf{Project} & \textbf{Sustainable} & \textbf{International} & \textbf{Relevant} & \textbf{Vital} & \textbf{Excellent} \\
\midrule
Governance & 2 & 3 & 2 & 2 & \textbf{5} \\
\addlinespace
Academic Dashboard & 3 & 2 & 2 & 4 & \textbf{5} \\
\addlinespace
Agenda Avanza & 1 & 3 & 4 & \textbf{5} & 4 \\
\addlinespace
University Experience 360 & 2 & 3 & 3 & \textbf{5} & 4 \\
\addlinespace
FADE Spaces & 4 & 2 & 3 & 4 & 3 \\
\addlinespace
FADE Global & 2 & \textbf{5} & 3 & 4 & 3 \\
\addlinespace
FADE Impact Lab & 3 & 3 & \textbf{5} & 3 & \textbf{5} \\
\addlinespace
COMUNICA & 2 & 2 & 4 & 2 & 3 \\
\bottomrule
\end{tabularx}

\vspace{4pt}
{\footnotesize\color{mutedcolor}\textit{The same levels are shown in each project's header in the plan's web summaries (Spanish and English).}}

\vspace{6pt}

\clearpage
\chapter[Governance]{\textcolor{gobcolor}{Governance}}
\begin{tcolorbox}[descbox]
\small This project provides FADE with a robust, transparent institutional framework, aligning internal processes, management reports and national and international accreditation requirements. Through the establishment of the Quality Committee and the work of the Academic Committees, it structures decision-making based on reliable data. It integrates processes, communication and artificial intelligence mechanisms to improve monitoring and accountability. Its aim is to make academic excellence manageable and assessable through coherent, auditable governance.
\end{tcolorbox}

\Needspace{18\baselineskip}
\section{Project Description}

The Governance Project aims to provide FADE with a robust, transparent institutional framework focused on continuous improvement, coherently integrating processes, decisions and accountability.   

The project seeks to promote closer alignment between the School's internal processes, Programme Management Reports, the School Management Report and the requirements of quality agencies and national and international accreditations. This is a natural step in a context where national and international accreditation are assessed at School level and the School is facing a period of institutional transformation. It does not replace programmes' academic autonomy but provides a common framework that enables that autonomy to be exercised coherently and in a coordinated manner. 

The Project integrates the four cross-cutting action mechanisms: Data, Processes, Communication and Artificial Intelligence. First, the Data mechanism ensures that governance decisions are based on reliable, traceable and comparable information, through connections to validated official sources, the development of institutional monitoring indicators and the systematic recording of evidence for national and international accreditation processes. Second, the Processes mechanism forms the backbone of the project by structuring and formalising academic governance procedures, from the establishment of the Quality Committee to the updating of the School Regulations, ensuring that decision-making takes place within clear, stable and auditable institutional frameworks. Third, the Communication mechanism acts as a catalyst for the project, establishing channels and protocols that ensure the effective transmission of decisions, the dissemination of accreditation outcomes and informed faculty participation in bachelor's and master's programme coordination meetings, thus turning institutional agreements into a shared academic culture. Artificial Intelligence (AI) joins this ecosystem as a fourth cross-cutting action mechanism, supporting information analysis. While Data ensures information reliability and Processes provides the institutional structure, AI helps improve analysis and monitoring, facilitating the identification of patterns and supporting tasks such as report preparation or academic information analysis. Its incorporation is envisaged as progressive and focused on specific use cases, as a tool supporting decision-making without replacing academic judgement. 

\Needspace{18\baselineskip}
\section{Relationship with the Strategic Lines}

The project's strategic contribution is distributed as follows: 

\begin{itemize}[noitemsep,topsep=2pt]
\item \textbf{Academic Excellence $\rightarrow$ 80\%}
\item \textbf{Impact $\rightarrow$ 20\%}
\end{itemize}

The 80\% contribution to Academic Excellence is based on the project's core purpose of structuring the School's governance and making it fully manageable. Governance directly underpins internal quality assurance processes, continuous programme improvement, the coherence of academic reports and the consolidation of a stable monitoring and accountability system. The 20\% contribution to Impact is justified by its institutional effects in terms of transparency, organisational efficiency and stronger academic governance, as well as its usefulness in external evaluation and accreditation processes. The project is structural in nature, generating lasting institutional capacity. 

The project contributes directly to Academic Excellence by providing a structured governance framework that enables a comprehensive analysis of the School's educational system. 

In particular, the governance project: 

\begin{itemize}[noitemsep,topsep=2pt]
\item Underpins internal quality assurance processes with reliable, comparable information. 

\item Supports continuous programme improvement by providing objective evidence for programme review. 

\item Reinforces the coherence and consistency of Programme and School Management Reports. 

\item Facilitates data-based academic decision-making, reducing information fragmentation. 

\item Integrates AACSB and national accreditation processes within a common quality assurance framework. 
\end{itemize}
The project does not itself measure academic excellence but makes it systematically manageable and assessable through robust, transparent governance mechanisms. 



The project's contribution to the Impact line is institutional and organisational, strengthening the School's capacity for academic governance and accountability. 

Specifically, the project: 

\begin{itemize}[noitemsep,topsep=2pt]
\item Improves academic transparency through structured, validated processes. 

\item Reinforces internal and external accountability, particularly in evaluation and accreditation processes. 

\item Increases institutional efficiency, reducing the administrative workload associated with manual process management. 

\item Strengthens the School's position in relation to AACSB and national accreditation. 
\end{itemize}

\newpage
\Needspace{18\baselineskip}
\section{Objectives, Action Plans and Development}
\Needspace{14\baselineskip}
\addsec{S.O.1. Align internal processes and strengthen the institutional governance of the School}
{\small\color{mutedcolor}\itshape Drive greater alignment between internal processes, Programme Management Reports, the School Management Report and accreditation requirements, ensuring institutional coherence and solid governance.}\par\vspace{6pt}

\begin{tcolorbox}[pabox={gobcolor}{A.P.1.1 \enspace\textbullet\enspace Establishment of the Quality Committee}]

\fieldlabel{Description}
{\small Establishment of the Quality Committee in compliance with the new UPV institutional requirements. It is a structural condition for the development of the improvement monitoring system (A.P.2.2).}\par
\fieldsep
\fieldlabel{Deliverables}
{\small \begin{itemize}[noitemsep,topsep=2pt]
\item Committee founding minutes
\item Annual work plan
\end{itemize}}\par
\fieldsep
\fieldlabel{KPIs}
{\small \begin{itemize}[noitemsep,topsep=2pt]
\item Committee established (Yes/No)
\item No. of meetings held per academic year
\item No. of agreements implemented on time
\end{itemize}}\par
\end{tcolorbox}
\begin{tcolorbox}[riskbox]
{\footnotesize\sffamily\bfseries\color{red!50!black} RISKS}\par\nopagebreak[4]
{\small Delays in formal establishment due to external institutional dependencies.}
\end{tcolorbox}
\begin{tcolorbox}[mitbox]
{\footnotesize\sffamily\bfseries\color{green!40!black} MITIGATION PLAN}\par\nopagebreak[4]
{\small Proactive coordination with the UPV to meet established institutional deadlines.}
\end{tcolorbox}
\vspace{8pt}

\begin{tcolorbox}[pabox={gobcolor}{A.P.1.2 \enspace\textbullet\enspace Review and Update of the School Regulations}]

\fieldlabel{Description}
{\small Review and update of the Regulations to ensure their adaptation to new institutional and academic requirements. It reinforces the regulatory framework for teaching coordination mechanisms (A.P.3.1).}\par
\fieldsep
\fieldlabel{Deliverables}
{\small Updated Regulations document}\par
\fieldsep
\fieldlabel{KPIs}
{\small Regulations approved by the School Council (Yes/No)}\par
\end{tcolorbox}
\begin{tcolorbox}[riskbox]
{\footnotesize\sffamily\bfseries\color{red!50!black} RISKS}\par\nopagebreak[4]
{\small Discrepancies in the interpretation of new institutional regulations.}
\end{tcolorbox}
\begin{tcolorbox}[mitbox]
{\footnotesize\sffamily\bfseries\color{green!40!black} MITIGATION PLAN}\par\nopagebreak[4]
{\small Prior consultation with UPV legal services before beginning drafting.}
\end{tcolorbox}
\vspace{8pt}

\newpage
\Needspace{14\baselineskip}
\addsec{S.O.2. Integrate accreditation and quality assurance processes}
{\small\color{mutedcolor}\itshape Leverage the work already started for AACSB accreditation to strengthen the position for national accreditation, integrating both processes into a common quality assurance framework.}\par\vspace{6pt}

\begin{tcolorbox}[pabox={gobcolor}{A.P.2.1 \enspace\textbullet\enspace Integration of AACSB and national accreditation processes}]

\fieldlabel{Description}
{\small Define the points of convergence between AACSB standards and national accreditation requirements, establishing a common framework of documentation and evidence that avoids duplication.}\par
\fieldsep
\fieldlabel{Deliverables}
{\small \begin{itemize}[noitemsep,topsep=2pt]
\item AACSB / national accreditation convergence map
\item Shared evidence repository
\item Usage guide for academic directors
\end{itemize}}\par
\fieldsep
\fieldlabel{KPIs}
{\small \begin{itemize}[noitemsep,topsep=2pt]
\item \% AACSB standards with identified national correspondence
\item No. of items of evidence shared between both processes
\end{itemize}}\par
\end{tcolorbox}
\begin{tcolorbox}[riskbox]
{\footnotesize\sffamily\bfseries\color{red!50!black} RISKS}\par\nopagebreak[4]
{\small Divergences between international standards and national requirements.}
\end{tcolorbox}
\begin{tcolorbox}[mitbox]
{\footnotesize\sffamily\bfseries\color{green!40!black} MITIGATION PLAN}\par\nopagebreak[4]
{\small Detailed prior analysis of the regulatory frameworks of both accreditations before defining the convergence map.}
\end{tcolorbox}
\vspace{8pt}

\begin{tcolorbox}[pabox={gobcolor}{A.P.2.2 \enspace\textbullet\enspace Monitoring and tracking of improvements by CATs}]

\fieldlabel{Description}
{\small The CATs will take on monitoring and tracking functions for improvements derived from management reports, ensuring the effective implementation of agreed improvement actions.}\par
\fieldsep
\fieldlabel{Deliverables}
{\small \begin{itemize}[noitemsep,topsep=2pt]
\item Periodic monitoring reports by programme
\item Record of implemented improvement actions
\end{itemize}}\par
\fieldsep
\fieldlabel{KPIs}
{\small \begin{itemize}[noitemsep,topsep=2pt]
\item \% improvement actions implemented vs. planned
\item \% actions closed within the deadline
\item No. of CAT meetings with improvement monitoring agenda
\end{itemize}}\par
\end{tcolorbox}
\begin{tcolorbox}[riskbox]
{\footnotesize\sffamily\bfseries\color{red!50!black} RISKS}\par\nopagebreak[4]
{\small Work overload for academic programme directors.}
\end{tcolorbox}
\begin{tcolorbox}[mitbox]
{\footnotesize\sffamily\bfseries\color{green!40!black} MITIGATION PLAN}\par\nopagebreak[4]
{\small Design of agile and standardised monitoring tools that minimise administrative burden.}
\end{tcolorbox}
\vspace{8pt}

\newpage
\Needspace{14\baselineskip}
\addsec{S.O.3. Strengthen teaching coordination and ensure curricular coherence}
{\small\color{mutedcolor}\itshape Maintain and strengthen bachelor's and master's programme coordination meetings as a fundamental forum for communication, curricular coherence and collaborative work among faculty, with the aim of improving coordination between subjects, reducing overlaps and ensuring a more coherent academic experience for students, while also ensuring the correct implementation of any new syllabus requirements.}\par\vspace{6pt}

\begin{tcolorbox}[pabox={gobcolor}{A.P.3.1 \enspace\textbullet\enspace Periodic coordination meetings for bachelor's and master's programmes}]

\fieldlabel{Description}
{\small Systematic planning and holding of coordination meetings, ensuring the participation of all involved faculty and the detailed explanation of curricular and extracurricular improvements introduced.}\par
\fieldsep
\fieldlabel{Deliverables}
{\small \begin{itemize}[noitemsep,topsep=2pt]
\item Meeting minutes with agreements, responsibilities and deadlines
\item Annual teaching coordination report
\end{itemize}}\par
\fieldsep
\fieldlabel{KPIs}
{\small \begin{itemize}[noitemsep,topsep=2pt]
\item No. of coordination meetings per programme and academic year
\item \% faculty attending vs. those invited
\end{itemize}}\par
\end{tcolorbox}
\begin{tcolorbox}[riskbox]
{\footnotesize\sffamily\bfseries\color{red!50!black} RISKS}\par\nopagebreak[4]
{\small \begin{itemize}[noitemsep,topsep=2pt]
\item • Low faculty participation in coordination meetings.
\item • Lack of follow-up on adopted agreements.
\item • That meetings become informational spaces without real impact.
\end{itemize}}
\end{tcolorbox}
\begin{tcolorbox}[mitbox]
{\footnotesize\sffamily\bfseries\color{green!40!black} MITIGATION PLAN}\par\nopagebreak[4]
{\small \begin{itemize}[noitemsep,topsep=2pt]
\item • Integration of meetings into the official academic calendar with formal notice.
\item • Agreement monitoring system with assigned responsibilities and periodic review.
\end{itemize}}
\end{tcolorbox}
\vspace{8pt}

\begin{tcolorbox}[pabox={gobcolor}{A.P.3.2 \enspace\textbullet\enspace Update of Syllabi}]

\fieldlabel{Description}
{\small Incorporation of the new syllabus fields for the following areas: Experiential Learning, SDGs and CSR, Technological Tools, Artificial Intelligence (AI) and AoL-AACSB.}\par
\fieldsep
\fieldlabel{Deliverables}
{\small \begin{itemize}[noitemsep,topsep=2pt]
\item Detailed instructions for faculty on the new fields
\item Compliance report by programme
\end{itemize}}\par
\fieldsep
\fieldlabel{KPIs}
{\small \begin{itemize}[noitemsep,topsep=2pt]
\item \% syllabi with the five new fields correctly completed
\item No. of incidents detected in syllabus review
\item \% incidents resolved before the deadline
\end{itemize}}\par
\end{tcolorbox}
\begin{tcolorbox}[riskbox]
{\footnotesize\sffamily\bfseries\color{red!50!black} RISKS}\par\nopagebreak[4]
{\small \begin{itemize}[noitemsep,topsep=2pt]
\item • Incomplete or inadequate completion of the new fields.
\item • Faculty resistance to new syllabus requirements.
\end{itemize}}
\end{tcolorbox}
\begin{tcolorbox}[mitbox]
{\footnotesize\sffamily\bfseries\color{green!40!black} MITIGATION PLAN}\par\nopagebreak[4]
{\small \begin{itemize}[noitemsep,topsep=2pt]
\item • Clear communication with concrete examples of correct completion.
\item • Informational session per programme with resolution of queries.
\end{itemize}}
\end{tcolorbox}
\vspace{8pt}

\begin{tcolorbox}[pabox={gobcolor}{A.P.3.3 \enspace\textbullet\enspace Effective transmission of agreements concerning institutional information}]

\fieldlabel{Description}
{\small Establish mechanisms to ensure that agreements adopted in the School Council, the Programme Academic Committees (CAT) and coordination meetings are communicated clearly, systematically and traceably to the affected groups: faculty, departments and student representatives. Includes post-meeting communication protocols, standardised formats for disseminating agreements and the active role of CAT representatives as liaisons with their departments.}\par
\fieldsep
\fieldlabel{Deliverables}
{\small \begin{itemize}[noitemsep,topsep=2pt]
\item Post-meeting communication protocol for governing bodies and CATs
\item Standard template for summarising and disseminating agreements
\item Record of communications made per academic year
\end{itemize}}\par
\fieldsep
\fieldlabel{KPIs}
{\small \begin{itemize}[noitemsep,topsep=2pt]
\item \% CAT meetings with documented communication to departments
\item \% School Council agreements with evidence of dissemination to faculty
\item No. of incidents due to lack of information detected in coordination meetings
\end{itemize}}\par
\end{tcolorbox}
\begin{tcolorbox}[riskbox]
{\footnotesize\sffamily\bfseries\color{red!50!black} RISKS}\par\nopagebreak[4]
{\small \begin{itemize}[noitemsep,topsep=2pt]
\item • Resistance or lack of time from representatives to take on this communication role.
\item • Lack of adequate channels to reach all groups.
\end{itemize}}
\end{tcolorbox}
\begin{tcolorbox}[mitbox]
{\footnotesize\sffamily\bfseries\color{green!40!black} MITIGATION PLAN}\par\nopagebreak[4]
{\small \begin{itemize}[noitemsep,topsep=2pt]
\item • Simplify the dissemination format as much as possible to reduce the burden.
\item • Integrate post-meeting communication as a formal step within each body's protocol.
\end{itemize}}
\end{tcolorbox}
\vspace{8pt}

\newpage
\Needspace{14\baselineskip}
\addsec{S.O.4. Improve decision-making through artificial intelligence applied to academic governance}
{\small\color{mutedcolor}\itshape Identify, design and deploy artificial intelligence applications that support the School's academic governance processes, advancing institutional analytical capacity from a descriptive, reactive approach to a predictive, prescriptive approach that enables more agile, well-founded and proactive decision-making. The objective is not limited to identifying opportunities but includes conducting pilots, evaluating their results and producing a map of future opportunities to guide the School's technology investment decisions.}\par\vspace{6pt}

\begin{tcolorbox}[pabox={gobcolor}{A.P.4.1 \enspace\textbullet\enspace Identification and Deployment of AI Applications in Academic Governance}]

\fieldlabel{Description}
{\small Identification and documentation of opportunities for AI use in governance processes (exemption reports, waivers and any process that could be automated using AI). Implementation will take place under human oversight and in compliance with Regulation (EU) 2024/1689 on Artificial Intelligence.}\par
\fieldsep
\fieldlabel{Deliverables}
{\small AI use cases document in academic governance}\par
\fieldsep
\fieldlabel{KPIs}
{\small \begin{itemize}[noitemsep,topsep=2pt]
\item No. of AI use cases identified and documented
\item No. of AI pilots started
\item No. of pilots completed with published results report
\item No. of future opportunities identified from pilot results
\end{itemize}}\par
\fieldsep
\fieldlabel{Required Resources}
{\small \begin{itemize}[noitemsep,topsep=2pt]
\item Head of Studies: 15--20\,h
\item Management / BI specialist: 30--40\,h
\item Vice Dean for Quality: 10--15\,h
\end{itemize}}\par
\end{tcolorbox}
\begin{tcolorbox}[riskbox]
{\footnotesize\sffamily\bfseries\color{red!50!black} RISKS}\par\nopagebreak[4]
{\small \begin{itemize}[noitemsep,topsep=2pt]
\item • Insufficient structured data to feed AI models.
\item • Institutional resistance to the use of advanced analytics tools.
\end{itemize}}
\end{tcolorbox}
\begin{tcolorbox}[mitbox]
{\footnotesize\sffamily\bfseries\color{green!40!black} MITIGATION PLAN}\par\nopagebreak[4]
{\small \begin{itemize}[noitemsep,topsep=2pt]
\item • Prioritise use cases that start from already available and validated data.
\item • Communicate the concrete benefits of each use case before its implementation.
\end{itemize}}
\end{tcolorbox}
\vspace{8pt}

\begin{tcolorbox}[pabox={gobcolor}{A.P.4.2 \enspace\textbullet\enspace Evaluation of Results and AI Roadmap in Governance}]

\fieldlabel{Description}
{\small Evaluation of the results of AI pilots executed, documentation of lessons learned and development of an AI roadmap for the academic governance of the School, prioritising the opportunities of greatest impact and feasibility for the next strategic cycle.}\par
\fieldsep
\fieldlabel{Deliverables}
{\small \begin{itemize}[noitemsep,topsep=2pt]
\item Evaluation report for each pilot (quality, time saved, errors, user satisfaction)
\item AI governance roadmap with prioritisation of opportunities for the next cycle
\end{itemize}}\par
\fieldsep
\fieldlabel{KPIs}
{\small \begin{itemize}[noitemsep,topsep=2pt]
\item No. of pilot evaluation reports produced
\item Roadmap approved by governing body (Yes/No)
\item No. of future opportunities identified and prioritised
\end{itemize}}\par
\fieldsep
\fieldlabel{Required Resources}
{\small \begin{itemize}[noitemsep,topsep=2pt]
\item Head of Studies: 10--15\,h
\item Management / BI specialist: 20--30\,h
\item Vice Dean for Quality: 10\,h
\end{itemize}}\par
\end{tcolorbox}
\begin{tcolorbox}[riskbox]
{\footnotesize\sffamily\bfseries\color{red!50!black} RISKS}\par\nopagebreak[4]
{\small \begin{itemize}[noitemsep,topsep=2pt]
\item • Insufficient pilot results to draw solid conclusions.
\item • Difficulty in comparing results between pilots of different nature.
\end{itemize}}
\end{tcolorbox}
\begin{tcolorbox}[mitbox]
{\footnotesize\sffamily\bfseries\color{green!40!black} MITIGATION PLAN}\par\nopagebreak[4]
{\small \begin{itemize}[noitemsep,topsep=2pt]
\item • Define evaluation criteria before the start of each pilot to ensure comparability of results.
\item • Establish a common evaluation report template for all pilots.
\end{itemize}}
\end{tcolorbox}
\vspace{8pt}

\newpage
\Needspace{14\baselineskip}
\addsec{S.O.5. Establish an adaptive monitoring model for the Strategic Plan}
{\small\color{mutedcolor}\itshape Incorporate formal periodic review mechanisms that allow priorities, projects and indicators to be adjusted without losing institutional coherence or traceability.}\par\vspace{6pt}

\begin{tcolorbox}[pabox={gobcolor}{A.P.5.1 \enspace\textbullet\enspace Periodic Reviews and Change Management of the Strategic Plan}]

\fieldlabel{Description}
{\small Establishment of a formal semester review cycle of the Strategic Plan that allows recording, justifying and communicating changes in projects, indicators and priorities, ensuring the traceability of decisions adopted and the overall coherence of the plan throughout its implementation. The Quality Committee approves the Plan and its changes, and its decisions are submitted to the School Council for ratification. Each academic year, the report on the Plan's level of execution is approved by the Quality Committee, ratified by the School Council and presented to the Advisory Board. The institution-wide risk analysis (academic, financial and operational risks) is reviewed every semester together with the Plan and updated every academic year with the annual report on the level of execution.}\par
\fieldsep
\fieldlabel{Deliverables}
{\small \begin{itemize}[noitemsep,topsep=2pt]
\item Calendar of semester plan reviews
\item Formal record of changes (modified, cancelled or incorporated projects, with justification)
\item Semester Strategic Plan control report presented to the Quality Committee
\item Annual report on the Plan's level of execution, approved by the Quality Committee, ratified by the School Council and presented to the Advisory Board
\item Ratification by the School Council of the Quality Committee's decisions on the Plan
\end{itemize}}\par
\fieldsep
\fieldlabel{KPIs}
{\small \begin{itemize}[noitemsep,topsep=2pt]
\item Semester reviews completed on time (Yes/No)
\item No. of changes formally recorded and justified
\item Semester report presented to the Quality Committee (Yes/No)
\item Annual report on the level of execution approved by the Quality Committee, ratified by the School Council and presented to the Advisory Board (Yes/No)
\item \% of Quality Committee decisions on the Plan ratified by the School Council
\end{itemize}}\par
\end{tcolorbox}
\begin{tcolorbox}[riskbox]
{\footnotesize\sffamily\bfseries\color{red!50!black} RISKS}\par\nopagebreak[4]
{\small \begin{itemize}[noitemsep,topsep=2pt]
\item • That reviews become a formal procedure without real impact on execution.
\item • Accumulation of undocumented changes that hinder subsequent monitoring.
\end{itemize}}
\end{tcolorbox}
\begin{tcolorbox}[mitbox]
{\footnotesize\sffamily\bfseries\color{green!40!black} MITIGATION PLAN}\par\nopagebreak[4]
{\small \begin{itemize}[noitemsep,topsep=2pt]
\item • Structured review format with assigned responsibilities and documented agreements.
\item • Mandatory recording of any scope change, cancellation or incorporation of projects.
\end{itemize}}
\end{tcolorbox}
\vspace{8pt}

\Needspace{18\baselineskip}
\section{Timeline}

\small

\begin{tabularx}{\textwidth}{>{\bfseries}l >{\raggedright\arraybackslash}X *{4}{>{\centering\arraybackslash}p{1.5cm}}}
\toprule
A.P. & Action plan & 2026/27 & 2027/28 & 2028/29 & 2029/30 \\
\midrule
A.P.1.1 & Establishment of the Quality Committee & $\blacksquare$ &  &  &  \\
A.P.1.2 & Review and updating of the Regulations & $\blacksquare$ & $\blacksquare$ &  &  \\
A.P.2.1 & Integration of AACSB and national accreditation & $\blacksquare$ & $\blacksquare$ &  &  \\
A.P.2.2 & Monitoring and tracking of improvements (CATs) &  & $\blacksquare$ & $\blacksquare$ &  \\
A.P.3.1 & Bachelor's and master's programme coordination meetings & $\blacksquare$ & $\blacksquare$ & $\blacksquare$ & $\blacksquare$ \\
A.P.3.2 & Updating of syllabi & $\blacksquare$ &  &  &  \\
A.P.3.3 & Effective transmission of agreements concerning institutional information &  & $\blacksquare$ & $\blacksquare$ & $\blacksquare$ \\
A.P.4.1 & Identification of AI applications in governance &  & $\blacksquare$ &  &  \\
A.P.4.2 & Evaluation of results and AI roadmap &  & $\blacksquare$ & $\blacksquare$ & $\blacksquare$ \\
A.P.5.1 & Periodic reviews and change management of the Plan & $\blacksquare$ & $\blacksquare$ & $\blacksquare$ & $\blacksquare$ \\
\bottomrule
\end{tabularx}




\Needspace{18\baselineskip}
\section{Summary of Objective Achievement}

\begin{tcolorbox}[descbox]
\small For the purposes of monitoring the Strategic Plan, achievement of the project's objectives will be assessed against the following criteria:
\end{tcolorbox}

\vspace{6pt}

\begin{tcolorbox}[pabox={gobcolor}{Objective 1. Strategic alignment and School governance}]
\small It is considered achieved when:
\begin{itemize}[noitemsep,topsep=2pt]
\item The Quality Committee operates regularly, with at least four meetings annually and more than 80 \% of agreements implemented.
\item The updated School Regulations have been formally approved and are in force.
\item There is documented, traceable coordination between internal processes, management reports and accreditation requirements.
\end{itemize}
\end{tcolorbox}

\vspace{6pt}

\begin{tcolorbox}[pabox={gobcolor}{Objective 2. Accreditation and quality assurance}]
\small It is considered achieved when:
\begin{itemize}[noitemsep,topsep=2pt]
\item AACSB and national accreditation processes are integrated within a common framework with an operational shared evidence repository.
\item CATs monitor improvement actions systematically and document that monitoring, with more than 75 \% of actions closed on time.
\item The average time to close improvement actions has decreased compared with the previous cycle.
\end{itemize}
\end{tcolorbox}

\vspace{6pt}

\begin{tcolorbox}[pabox={gobcolor}{Objective 3. Teaching coordination and curricular coherence}]
\small It is considered achieved when:
\begin{itemize}[noitemsep,topsep=2pt]
\item Bachelor's and master's programme coordination meetings are held systematically, with agreements documented and their implementation monitored.
\item All syllabi correctly incorporate the five new fields in Block 3, Contextualisation (New Requirements 2026-27), with a compliance rate above 95 \%.
\end{itemize}
\end{tcolorbox}

\vspace{6pt}

\begin{tcolorbox}[pabox={gobcolor}{Objective 4. Improved decision-making through artificial intelligence applied to academic governance}]
\small It is considered achieved when:
\begin{itemize}[noitemsep,topsep=2pt]
\item At least five AI use cases in academic governance processes have been identified and documented.
\item At least two AI pilots have been launched and all have a published results evaluation report.
\item An AI roadmap with at least three prioritised opportunities for the next strategic cycle has been prepared and approved by the governing body.
\end{itemize}
\end{tcolorbox}

\vspace{6pt}

\begin{tcolorbox}[pabox={gobcolor}{Objective 5. Adaptive monitoring of the Strategic Plan}]
\small It is considered achieved when:
\begin{itemize}[noitemsep,topsep=2pt]
\item Semester reviews of the Strategic Plan are completed on time in every academic year of the 2026--2030 period.
\item The semester Strategic Plan control report is presented to the Quality Committee at each semester review.
\item All changes in scope, project cancellations or additions are formally recorded and justified in the change register.
\item The annual report on the Plan's level of execution is approved by the Quality Committee, ratified by the School Council and presented to the Advisory Board in each of the four academic years of the period.
\item All Quality Committee decisions on the Plan have been ratified by the School Council.
\end{itemize}
\end{tcolorbox}

\clearpage

\chapter[Academic Dashboard]{\textcolor{dashcolor}{Academic Dashboard}}
\begin{tcolorbox}[descbox]
\small Creation of a stable information infrastructure that integrates and visualises key academic performance data for the School and its programmes. It automates the generation of management reports and supports quality assurance processes through documented indicator maps and validated official sources. It facilitates proactive, well-founded decision-making in governing bodies, promoting a culture of data and evidence-based management.
\end{tcolorbox}

\Needspace{18\baselineskip}
\section{Project Description}

\small
The Academic Dashboard of the School of Management (FADE) is a strategic project intended to provide the School with a stable academic information infrastructure that enables the coherent integration, structuring and visualisation of the main data relating to its academic activity.

The project aims to provide a comprehensive, reliable and traceable view of the academic performance of the School and its programmes, supporting decision-making, the preparation of management reports and quality assurance processes. The dashboard is built on previously defined and documented indicator maps for academic areas, and uses business intelligence tools as the final visualisation layer.

The Academic Dashboard of the School of Management is consolidated as a strategic tool through the synergistic integration of the four cross-cutting action mechanisms: Data, Processes, Communication and Artificial Intelligence. First, the Data mechanism ensures the development of an institutional tool through the rigorous definition of indicator maps and exclusive connections to validated official sources, thus ensuring information quality, traceability and historical comparability. Second, the Processes mechanism goes beyond visualisation by structurally integrating the dashboard into academic operations, automating the generation of Programme Management Reports and standardising performance analysis within a common framework. Third, the Communication mechanism acts as a catalyst for the project, establishing clear access protocols and promoting systematic use of the tool by governing bodies, transforming data into shared institutional information and fostering accountability and evidence-based decision-making. Artificial Intelligence (AI) joins this ecosystem as a fourth cross-cutting action mechanism, acting as the project's advanced analytical layer. While Data ensures reliability and Processes provides structure, AI supports information analysis, facilitating pattern identification and supporting tasks such as report preparation or academic performance analysis.

\Needspace{18\baselineskip}
\section{Relationship with the Strategic Lines}

\small
The project's strategic contribution is distributed as follows:
\begin{itemize}[noitemsep,topsep=2pt]
\item \textbf{Academic Excellence $\rightarrow$ 80\%}
\item \textbf{Impact $\rightarrow$ 20\%}
\end{itemize}
The 80\% contribution to Academic Excellence is based on the project's core purpose of structuring, documenting and making the School's academic performance fully manageable. The dashboard directly underpins internal quality assurance processes, continuous programme improvement, the methodological coherence of academic reports and the consolidation of a stable system for monitoring educational performance. The 20\% contribution to Impact is justified by its institutional effects in terms of transparency, accountability, organisational efficiency and stronger academic governance, as well as its usefulness in external evaluation and accreditation processes. The project is structural in nature, generating lasting institutional capacity.



\small
The project contributes directly to Academic Excellence by providing a structured, documented academic information system that enables a comprehensive analysis of the School's educational system.

In particular, the dashboard:
\begin{itemize}[noitemsep,topsep=2pt]
\item Underpins internal quality assurance processes with reliable, comparable information.
\item Supports continuous programme improvement by providing objective evidence for programme review.
\item Reinforces the coherence and consistency of Programme and School Management Reports.
\item Facilitates data-based academic decision-making, reducing information fragmentation.
\end{itemize}

The project does not itself measure academic excellence but makes it systematically manageable and assessable.


\small
The project's contribution to the Impact line is institutional and organisational, strengthening the School's capacity for academic governance and accountability.

Specifically, the project:
\begin{itemize}[noitemsep,topsep=2pt]
\item Improves academic transparency through structured, validated information.
\item Reinforces internal and external accountability, particularly in evaluation and accreditation processes.
\item Increases institutional efficiency, reducing the administrative workload associated with manual report preparation.
\end{itemize}

\newpage
\Needspace{18\baselineskip}
\section{Objectives, Action Plans and Development}
\Needspace{14\baselineskip}
\addsec{S.O.1. Have an integrated view of the School's academic performance}
{\small\color{mutedcolor}\itshape Integrate relevant academic information, currently dispersed, into a single environment to enable a joint, coherent and comparable analysis of the academic functioning of the School and its programmes.}\par\vspace{6pt}

\begin{tcolorbox}[pabox={dashcolor}{A.P.1.1 \enspace\textbullet\enspace Definition of the Academic Framework and Dashboard Areas}]

\fieldlabel{Description}
{\small Definition of the academic areas that structure the dashboard and their relationship with the School's academic processes, with the aim of ensuring a consistent interpretation of academic performance and avoiding inconsistent analyses.}\par
\fieldsep
\fieldlabel{Deliverables}
{\small \begin{itemize}[noitemsep,topsep=2pt]
\item Academic framework document for the dashboard
\item Formal definition of academic areas (list and scope)
\item Conceptual map of areas--academic processes
\end{itemize}}\par
\fieldsep
\fieldlabel{KPIs}
{\small \begin{itemize}[noitemsep,topsep=2pt]
\item Document approved (Yes/No)
\item No. and \% of academic areas defined
\end{itemize}}\par
\end{tcolorbox}
\begin{tcolorbox}[riskbox]
{\footnotesize\sffamily\bfseries\color{red!50!black} RISKS}\par\nopagebreak[4]
{\small Lack of conceptual coherence of the dashboard.}
\end{tcolorbox}
\begin{tcolorbox}[mitbox]
{\footnotesize\sffamily\bfseries\color{green!40!black} MITIGATION PLAN}\par\nopagebreak[4]
{\small Formal validation of the academic framework by a governing body.}
\end{tcolorbox}
\vspace{8pt}

\begin{tcolorbox}[pabox={dashcolor}{A.P.1.2 \enspace\textbullet\enspace Design of the Functional Structure of the Dashboard}]

\fieldlabel{Description}
{\small Design of pages, blocks, filters and levels of analysis of the dashboard to facilitate its use and understanding by different academic user profiles.}\par
\fieldsep
\fieldlabel{Deliverables}
{\small Dashboard architecture (pages, blocks and navigation)}\par
\fieldsep
\fieldlabel{KPIs}
{\small \begin{itemize}[noitemsep,topsep=2pt]
\item No. of views designed and \% according to planning
\item Functional validation by key users (Yes/No)
\end{itemize}}\par
\end{tcolorbox}
\begin{tcolorbox}[riskbox]
{\footnotesize\sffamily\bfseries\color{red!50!black} RISKS}\par\nopagebreak[4]
{\small Low dashboard usability.}
\end{tcolorbox}
\begin{tcolorbox}[mitbox]
{\footnotesize\sffamily\bfseries\color{green!40!black} MITIGATION PLAN}\par\nopagebreak[4]
{\small Iterative design with validation by key users.}
\end{tcolorbox}
\vspace{8pt}

\newpage
\Needspace{14\baselineskip}
\addsec{S.O.2. Define and document academic areas and their indicator maps}
{\small\color{mutedcolor}\itshape Structure the dashboard based on documented and stable academic areas, avoiding the use of isolated indicators and ensuring coherence over time.}\par\vspace{6pt}

\begin{tcolorbox}[pabox={dashcolor}{A.P.2.1 \enspace\textbullet\enspace Definition and Development of Indicator Maps by Academic Areas}]

\fieldlabel{Description}
{\small Definition and documentation of indicator maps by area, establishing what is measured, with which indicators, from which sources and for which analytical and management purpose.}\par
\fieldsep
\fieldlabel{Deliverables}
{\small \begin{itemize}[noitemsep,topsep=2pt]
\item Documented indicator maps by academic area
\item Technical indicator sheets (definition, formula, source, periodicity, level)
\item Traceability matrix: indicator $\rightarrow$ management report section
\end{itemize}}\par
\fieldsep
\fieldlabel{KPIs}
{\small \begin{itemize}[noitemsep,topsep=2pt]
\item No. and \% of areas with complete indicator map
\item No. and \% of dashboard indicators included in each map
\end{itemize}}\par
\end{tcolorbox}
\begin{tcolorbox}[riskbox]
{\footnotesize\sffamily\bfseries\color{red!50!black} RISKS}\par\nopagebreak[4]
{\small Inconsistent definitions between areas or programmes.}
\end{tcolorbox}
\begin{tcolorbox}[mitbox]
{\footnotesize\sffamily\bfseries\color{green!40!black} MITIGATION PLAN}\par\nopagebreak[4]
{\small Use of a common standard template and institutional validation of the maps.}
\end{tcolorbox}
\vspace{8pt}

\begin{tcolorbox}[pabox={dashcolor}{A.P.2.2 \enspace\textbullet\enspace Institutional Validation and Consolidation of Indicator Maps}]

\fieldlabel{Description}
{\small Review, validation and institutional consolidation of the defined indicator maps, ensuring their acceptance as a common reference framework for the dashboard and management reports.}\par
\fieldsep
\fieldlabel{Deliverables}
{\small Document of criteria on the use of indicators}\par
\fieldsep
\fieldlabel{KPIs}
{\small No. and \% of areas with validated maps}\par
\end{tcolorbox}
\begin{tcolorbox}[riskbox]
{\footnotesize\sffamily\bfseries\color{red!50!black} RISKS}\par\nopagebreak[4]
{\small Partial or informal use of the maps.}
\end{tcolorbox}
\begin{tcolorbox}[mitbox]
{\footnotesize\sffamily\bfseries\color{green!40!black} MITIGATION PLAN}\par\nopagebreak[4]
{\small Formal approval and communication of the common framework.}
\end{tcolorbox}
\vspace{8pt}

\newpage
\Needspace{14\baselineskip}
\addsec{S.O.3. Ensure the quality and traceability of academic information}
{\small\color{mutedcolor}\itshape Ensure that dashboard data are reliable, sourced from official sources, correctly documented and comparable across academic years and programmes.}\par\vspace{6pt}

\begin{tcolorbox}[pabox={dashcolor}{A.P.3.1 \enspace\textbullet\enspace Identification and Validation of Official Data Sources}]

\fieldlabel{Description}
{\small Identification of the official data sources feeding the dashboard, assigning responsibilities for their validation, updating and maintenance, and establishing usage criteria.}\par
\fieldsep
\fieldlabel{Deliverables}
{\small Inventory of official sources by area and indicator}\par
\fieldsep
\fieldlabel{KPIs}
{\small \begin{itemize}[noitemsep,topsep=2pt]
\item No. and \% of official sources identified and validated
\item \% of indicators with assigned source
\end{itemize}}\par
\end{tcolorbox}
\begin{tcolorbox}[riskbox]
{\footnotesize\sffamily\bfseries\color{red!50!black} RISKS}\par\nopagebreak[4]
{\small Use of non-official or outdated data.}
\end{tcolorbox}
\begin{tcolorbox}[mitbox]
{\footnotesize\sffamily\bfseries\color{green!40!black} MITIGATION PLAN}\par\nopagebreak[4]
{\small Validated formal inventory of sources and data owners.}
\end{tcolorbox}
\vspace{8pt}

\begin{tcolorbox}[pabox={dashcolor}{A.P.3.2 \enspace\textbullet\enspace Normalisation, Versioning and Control of Academic Indicators}]

\fieldlabel{Description}
{\small Definition of common criteria for the calculation, updating and versioning of academic indicators, as well as quality control mechanisms that ensure temporal comparability.}\par
\fieldsep
\fieldlabel{Deliverables}
{\small Normalised and versioned indicator catalogue}\par
\fieldsep
\fieldlabel{KPIs}
{\small \% of indicators with complete technical sheet}\par
\end{tcolorbox}
\begin{tcolorbox}[riskbox]
{\footnotesize\sffamily\bfseries\color{red!50!black} RISKS}\par\nopagebreak[4]
{\small Loss of comparability between academic years.}
\end{tcolorbox}
\begin{tcolorbox}[mitbox]
{\footnotesize\sffamily\bfseries\color{green!40!black} MITIGATION PLAN}\par\nopagebreak[4]
{\small Mandatory recording of methodological changes.}
\end{tcolorbox}
\vspace{8pt}

\newpage
\Needspace{14\baselineskip}
\addsec{S.O.4. Facilitate the development of Programme and School Management Reports}
{\small\color{mutedcolor}\itshape Use the dashboard as a support tool for the systematic development of reports, ensuring coherence, homogeneity and comparability through a common and stable framework.}\par\vspace{6pt}

\begin{tcolorbox}[pabox={dashcolor}{A.P.4.1 \enspace\textbullet\enspace Alignment of Indicator Maps with the Structure of Management Reports}]

\fieldlabel{Description}
{\small Explicitly define the correspondence between indicator maps and the structure of Programme and School Management Reports, establishing which indicators feed each section of the report.}\par
\fieldsep
\fieldlabel{Deliverables}
{\small \begin{itemize}[noitemsep,topsep=2pt]
\item Mapping document: area $\rightarrow$ map $\rightarrow$ report sections
\item Dashboard usage guide for preparing the management report
\end{itemize}}\par
\fieldsep
\fieldlabel{KPIs}
{\small \% report sections supported by dashboard indicators}\par
\end{tcolorbox}
\begin{tcolorbox}[riskbox]
{\footnotesize\sffamily\bfseries\color{red!50!black} RISKS}\par\nopagebreak[4]
{\small \begin{itemize}[noitemsep,topsep=2pt]
\item • Persistence of heterogeneous reports between programmes.
\item • Partial use of the dashboard as an information source.
\end{itemize}}
\end{tcolorbox}
\begin{tcolorbox}[mitbox]
{\footnotesize\sffamily\bfseries\color{green!40!black} MITIGATION PLAN}\par\nopagebreak[4]
{\small \begin{itemize}[noitemsep,topsep=2pt]
\item • Institutional approval of the common template.
\item • Mandatory use of the dashboard as the base source of the report.
\end{itemize}}
\end{tcolorbox}
\vspace{8pt}

\begin{tcolorbox}[pabox={dashcolor}{A.P.4.2 \enspace\textbullet\enspace Automation of Key Content in Management Reports}]

\fieldlabel{Description}
{\small Automate the generation of tables, graphs and recurring information blocks in management reports, reducing administrative burden and minimising errors arising from manual data manipulation.}\par
\fieldsep
\fieldlabel{Deliverables}
{\small \begin{itemize}[noitemsep,topsep=2pt]
\item Automated tables and visualisations in the dashboard
\item Validation checklist prior to report publication
\end{itemize}}\par
\fieldsep
\fieldlabel{KPIs}
{\small \begin{itemize}[noitemsep,topsep=2pt]
\item \% of report content generated automatically
\item Reduction in average report preparation time (\%)
\end{itemize}}\par
\end{tcolorbox}
\begin{tcolorbox}[riskbox]
{\footnotesize\sffamily\bfseries\color{red!50!black} RISKS}\par\nopagebreak[4]
{\small Excessive dependence on manual processes.}
\end{tcolorbox}
\begin{tcolorbox}[mitbox]
{\footnotesize\sffamily\bfseries\color{green!40!black} MITIGATION PLAN}\par\nopagebreak[4]
{\small Prioritise the automation of recurring and critical content.}
\end{tcolorbox}
\vspace{8pt}

\newpage
\Needspace{14\baselineskip}
\addsec{S.O.5. Promote the systematic use of academic information in School management}
{\small\color{mutedcolor}\itshape Consolidate the academic dashboard as a regularly used tool in academic bodies and in the School's monitoring and decision-making processes, promoting a culture of data-based academic management. This objective has a clearly organisational and cultural component, aimed at ensuring that the dashboard is not an occasional tool but a resource integrated into the School's day-to-day management.}\par\vspace{6pt}

\begin{tcolorbox}[pabox={dashcolor}{A.P.5.1 \enspace\textbullet\enspace Definition of Access Levels and Dashboard Usage Rules}]

\fieldlabel{Description}
{\small Define the access profiles and visibility levels of dashboard information based on different academic and management roles, establishing clear institutional usage rules.}\par
\fieldsep
\fieldlabel{Deliverables}
{\small Institutional dashboard usage protocol}\par
\fieldsep
\fieldlabel{KPIs}
{\small }\par
\end{tcolorbox}
\begin{tcolorbox}[riskbox]
{\footnotesize\sffamily\bfseries\color{red!50!black} RISKS}\par\nopagebreak[4]
{\small \begin{itemize}[noitemsep,topsep=2pt]
\item • Inappropriate access to sensitive information.
\item • Reluctance to use due to lack of clarity.
\end{itemize}}
\end{tcolorbox}
\begin{tcolorbox}[mitbox]
{\footnotesize\sffamily\bfseries\color{green!40!black} MITIGATION PLAN}\par\nopagebreak[4]
{\small \begin{itemize}[noitemsep,topsep=2pt]
\item • Clear and communicated definition of profiles and permissions.
\item • Periodic review of access permissions.
\end{itemize}}
\end{tcolorbox}
\vspace{8pt}

\begin{tcolorbox}[pabox={dashcolor}{A.P.5.2 \enspace\textbullet\enspace Implementation and Institutional Use of the Dashboard in Management Processes}]

\fieldlabel{Description}
{\small Integrate the use of the dashboard into the usual academic monitoring processes and governing bodies, promoting its use as the main source of information for analysis and decision-making.}\par
\fieldsep
\fieldlabel{Deliverables}
{\small \begin{itemize}[noitemsep,topsep=2pt]
\item Integration of the dashboard into the annual management report process
\item Implementation and institutional use report
\end{itemize}}\par
\fieldsep
\fieldlabel{KPIs}
{\small \begin{itemize}[noitemsep,topsep=2pt]
\item No. of academic bodies using the dashboard
\item No. of reports and analyses prepared from the dashboard
\end{itemize}}\par
\end{tcolorbox}
\begin{tcolorbox}[riskbox]
{\footnotesize\sffamily\bfseries\color{red!50!black} RISKS}\par\nopagebreak[4]
{\small \begin{itemize}[noitemsep,topsep=2pt]
\item • Occasional or residual use of the dashboard.
\item • Perception of the dashboard as an extra tool.
\end{itemize}}
\end{tcolorbox}
\begin{tcolorbox}[mitbox]
{\footnotesize\sffamily\bfseries\color{green!40!black} MITIGATION PLAN}\par\nopagebreak[4]
{\small \begin{itemize}[noitemsep,topsep=2pt]
\item • Integration of the dashboard into management routines and calendars.
\item • Explicit use in key meetings and processes.
\end{itemize}}
\end{tcolorbox}
\vspace{8pt}

\Needspace{18\baselineskip}
\section{Timeline}

\small

\begin{tabularx}{\textwidth}{>{\bfseries}l >{\raggedright\arraybackslash}X *{4}{>{\centering\arraybackslash}p{1.5cm}}}
\toprule
A.P. & Action plan & 2026/27 & 2027/28 & 2028/29 & 2029/30 \\
\midrule
A.P.1.1 & Definition of the academic framework and areas & $\blacksquare$ &  &  &  \\
A.P.1.2 & Design of the Functional Structure of the Dashboard & $\blacksquare$ & $\blacksquare$ &  &  \\
A.P.2.1 & Preparation of indicator maps by area & $\blacksquare$ & $\blacksquare$ &  &  \\
A.P.2.2 & Validation and consolidation of the maps &  & $\blacksquare$ &  &  \\
A.P.3.1 & Identification and validation of data sources &  & $\blacksquare$ &  &  \\
A.P.3.2 & Normalisation and versioning of indicators &  & $\blacksquare$ & $\blacksquare$ &  \\
A.P.4.1 & Alignment of maps--management reports &  &  & $\blacksquare$ &  \\
A.P.4.2 & Automation of report content &  &  & $\blacksquare$ & $\blacksquare$ \\
A.P.5.1 & Definition of access rights and usage rules &  &  &  & $\blacksquare$ \\
A.P.5.2 & Implementation and institutional use of the dashboard &  &  &  & $\blacksquare$ \\
\bottomrule
\end{tabularx}


\Needspace{18\baselineskip}
\section{Summary of Objective Achievement}

\small
For the purposes of monitoring the Strategic Plan, achievement of the project's objectives will be assessed against the following criteria:

\vspace{4pt}

\begin{itemize}[leftmargin=1.5em,itemsep=6pt,topsep=4pt]

\item \textbf{Objective 1. Obtain an integrated view of the School's academic performance.}  
It is considered achieved when the School has an operational academic dashboard that coherently integrates key academic information about the School and its programmes, enabling combined analysis and monitoring over time.

\item \textbf{Objective 2. Define and document the academic areas and their indicator maps.}  
It is considered achieved when all defined academic areas have a documented indicator map that has been institutionally validated and is used consistently in the dashboard and management reports.

\item \textbf{Objective 3. Ensure the quality and traceability of academic information.}  
It is considered achieved when dashboard indicators draw exclusively on validated official sources, have standardised/versioned definitions and enable reliable comparisons between academic years.

\item \textbf{Objective 4. Facilitate the preparation of Programme and School Management Reports.}  
It is considered achieved when Programme and School Management Reports are prepared systematically from the academic dashboard, using a common structure and largely automated content.

\item \textbf{Objective 5. Promote systematic use of academic information in School management.}  
It is considered achieved when the academic dashboard is used regularly in academic bodies and in the School's monitoring and decision-making processes.

\end{itemize}

\clearpage

\chapter[Agenda Avanza]{\textcolor{aacolor}{Agenda Avanza}}
\begin{tcolorbox}[descbox]
\small This project ensures the continuous updating of the academic offering and the adaptation of bachelor's and master's programmes to Royal Decree 822/2021 and to the demands of the professional environment. It includes the introduction of new strategic programmes, such as the Bachelor's Degree in Business and Digital Technology, and the progressive integration of artificial intelligence in the classroom. It consolidates an institutional model for curriculum review that strengthens academic quality, graduate employability and the School's capacity to adapt.
\end{tcolorbox}

\Needspace{18\baselineskip}
\section{Project Description}

\small
Agenda Avanza is the strategic project of the UPV School of Management (FADE) aimed at ensuring the continuous, structured and progressive updating of its academic offering, aligning it with emerging needs in the professional, technological and social environment.

The project aims to consolidate an institutional model of programme development that enables the systematic adaptation, introduction and improvement of educational programmes, strengthening academic quality, graduate employability and the School's capacity to adapt.

Agenda Avanza brings together four complementary cross-cutting action mechanisms: Processes, Data, Communication and Artificial Intelligence. First, the academic Processes mechanism ensures that the review, improvement and effective transformation of curricula, subjects and teaching methodologies are integrated into existing institutional structures (CATs and quality systems). Second, the Data mechanism promotes the use of objective institutional information to support and substantiate academic decision-making. Third, the Communication mechanism facilitates the dissemination, transfer and recognition of the innovation and updating activities undertaken in the School. Artificial Intelligence (AI) joins this ecosystem as a fourth cross-cutting action mechanism, acting as the technological core of teaching innovation. While Processes structures curricular updating, AI enhances the academic offering by being progressively integrated as cross-cutting content and a learning methodology in the classroom. Its systematic incorporation, resulting from the transfer of institutional innovation projects, ensures that students acquire advanced technological competencies, consolidating FADE's position as an institution that adapts its teaching to the actual and future demands of the professional environment.

The project integrates the adaptation of existing programmes to Royal Decree 822/2021, the introduction of new strategic programmes, the progressive incorporation of teaching innovation and the updating of subjects and content, establishing a structured model for the continuous updating of the School's academic offering.

\Needspace{18\baselineskip}
\section{Relationship with the Strategic Lines}

\small
The Agenda Avanza project contributes to both strategic lines of the School's Plan 2026--2030, with the following distribution:

\begin{itemize}[noitemsep,topsep=2pt]
\item \textbf{Academic Excellence $\rightarrow$ 75\%}
\item \textbf{Impact $\rightarrow$ 25\%}
\end{itemize}
The project's main contribution is to Academic Excellence, focusing on the updating, design, introduction and continuous improvement of the School's bachelor's and master's programmes, as well as teaching innovation and the integration of new methodologies and technologies into the teaching and learning process.

In particular, the project addresses:
\begin{itemize}[noitemsep,topsep=2pt]
\item The adaptation of the Bachelor's Degree in Business Administration and Management (ADE), as well as the MUDFF, MUGEPS and MUGA master's programmes, to Royal Decree 822/2021.
\item The review of the Bachelor's Degree in Public Administration and Management (GAP), already adapted to regulatory requirements, with the aim of strengthening its technological orientation and its connection with modern public management.
\item The introduction of the new Bachelor's Degree in Business and Digital Technology, whose rollout will begin in the 2026/2027 academic year.
\item The launch of the Master's Degree in Corporate Responsibility and Sustainability.
\item The coordination of educational innovation projects, particularly those aimed at progressively integrating artificial intelligence into teaching, as well as future initiatives to modernise and improve programmes.
\end{itemize}

The project also contributes to Impact (25\%) insofar as these activities improve the fit between the academic offering and the demands of the professional and social environment, strengthen graduate employability and reinforce the School's positioning and reputation.

\newpage

\Needspace{18\baselineskip}
\section{Objectives, Action Plans and Development}
\Needspace{14\baselineskip}
\addsec{S.O.1. Adapt official degree programmes to RD 822/2021}
{\small\color{mutedcolor}\itshape Review, redesign and adapt the ADE bachelor's programme and the MUGA, MUDFF and MUGEPS master's programmes to RD 822/2021, ensuring educational coherence, competency alignment and compliance with quality assurance procedures.}\par\vspace{6pt}

\begin{tcolorbox}[pabox={aacolor}{A.P.1.1 \enspace\textbullet\enspace Diagnosis and Analysis of Adaptation to RD 822/2021}]

\fieldlabel{Description}
{\small Systematic analysis of programmes (ADE, MUGA, MUDFF and MUGEPS), identifying their degree of alignment with the requirements of Royal Decree 822/2021 in terms of structure, competencies, methodologies, assessment and quality assurance. The analysis will follow the criteria and methodological guidelines established by the Quality and Accreditation Area of the UPV Vice-Rectorate for Planning, Studies, Quality and Accreditation.}\par
\fieldsep
\fieldlabel{Deliverables}
{\small \begin{itemize}[noitemsep,topsep=2pt]
\item RD 822/2021 compliance evaluation matrix by programme
\item Diagnostic report by programme
\item Map of external regulatory milestones with realistic approval deadlines (UPV, ANECA, AVAP)
\end{itemize}}\par
\fieldsep
\fieldlabel{KPIs}
{\small \begin{itemize}[noitemsep,topsep=2pt]
\item \% programmes analysed (target: 100\%)
\item No. of reports prepared (target: 4)
\end{itemize}}\par
\end{tcolorbox}
\begin{tcolorbox}[riskbox]
{\footnotesize\sffamily\bfseries\color{red!50!black} RISKS}\par\nopagebreak[4]
{\small \begin{itemize}[noitemsep,topsep=2pt]
\item • Dependence on external UPV criteria and guidelines.
\item • Heterogeneous interpretation of guidelines between programmes.
\item • Uneven depth in the diagnoses.
\item • Overload of academic programme directors.
\item • Difficulties in collecting and organising the necessary information.
\item • External verification deadlines (ANECA/AVAP) that exceed the horizon of the strategic plan.
\end{itemize}}
\end{tcolorbox}
\begin{tcolorbox}[mitbox]
{\footnotesize\sffamily\bfseries\color{green!40!black} MITIGATION PLAN}\par\nopagebreak[4]
{\small \begin{itemize}[noitemsep,topsep=2pt]
\item • Continuous alignment with the Vice-Rectorate for Programme Management.
\item • Internal coordination and comparison sessions between programmes.
\item • Use of common analysis templates.
\item • Planning of work in phases with a defined calendar.
\item • Technical support from the quality area.
\item • Early preparation of the regulatory milestone map and close coordination with the UPV General Secretariat to adjust the timeline to actual external approval deadlines.
\end{itemize}}
\end{tcolorbox}
\vspace{8pt}

\begin{tcolorbox}[pabox={aacolor}{A.P.1.2 \enspace\textbullet\enspace Redesign of Study Plans and Verification Reports}]

\fieldlabel{Description}
{\small Redesign the study plans for the ADE bachelor's programme, MUDFF, MUGEPS and MUGA, adapting the verification reports with the elements required by the regulations and processing them with the agencies until obtaining final verification.}\par
\fieldsep
\fieldlabel{Deliverables}
{\small \begin{itemize}[noitemsep,topsep=2pt]
\item Programme specifications submitted for evaluation
\item ANECA/AVAP evaluation reports
\item Verification resolutions
\end{itemize}}\par
\fieldsep
\fieldlabel{KPIs}
{\small \begin{itemize}[noitemsep,topsep=2pt]
\item \% verification reports submitted (target: 100\%)
\item \% verification reports favourably verified
\item No. of requirements/corrections per verification report
\item Average verification time
\end{itemize}}\par
\end{tcolorbox}
\begin{tcolorbox}[riskbox]
{\footnotesize\sffamily\bfseries\color{red!50!black} RISKS}\par\nopagebreak[4]
{\small \begin{itemize}[noitemsep,topsep=2pt]
\item • Delays in the preparation of the verification reports.
\item • Difficulty in aligning different academic visions.
\item • Inconsistencies between sections of the verification report.
\item • Excessive dependence on a few key people.
\item • Delays arising from the external evaluation process.
\item • That the redesign is limited to a formal adaptation without effective updating of content.
\end{itemize}}
\end{tcolorbox}
\begin{tcolorbox}[mitbox]
{\footnotesize\sffamily\bfseries\color{green!40!black} MITIGATION PLAN}\par\nopagebreak[4]
{\small \begin{itemize}[noitemsep,topsep=2pt]
\item • Establish a detailed timeline with partial milestones and intermediate reviews.
\item • Maintain central methodological coordination.
\item • Use common templates and internal validation procedures.
\item • Distribute tasks among several people and document changes.
\item • Thorough internal review prior to submission.
\item • Incorporate specific review of content and methodologies in the academic redesign.
\end{itemize}}
\end{tcolorbox}
\vspace{8pt}

\newpage
\Needspace{14\baselineskip}
\addsec{S.O.2. Implement and consolidate the new strategic academic offering}
{\small\color{mutedcolor}\itshape Introduce the Bachelor's Degree in Business and Digital Technology and the Master's Degree in Corporate Responsibility and Sustainability in accordance with their approved verification reports, ensuring their successful launch, academic development and integration into the School's coordination, monitoring and quality assurance systems. 

The Bachelor's Degree in Business and Digital Technology will be introduced progressively by academic year, from the 2026/2027 academic year until full implementation. 

The Master's Degree in Corporate Responsibility and Sustainability will be fully introduced from the 2026/2027 academic year, ensuring its appropriate integration into the School's academic and management structure. }\par\vspace{6pt}

\begin{tcolorbox}[pabox={aacolor}{A.P.2.1 \enspace\textbullet\enspace Activation of Academic Management for the Degree in Business and Digital Technology}]

\fieldlabel{Description}
{\small Establishment of the CAT and definition of teaching coordination mechanisms. It will oversee the progressive implementation of the degree and will prepare the annual Programme Management Report.}\par
\fieldsep
\fieldlabel{Deliverables}
{\small \begin{itemize}[noitemsep,topsep=2pt]
\item CAT established and operational
\item Evidence of CAT functioning
\item Teaching coordination mechanisms
\item Annual Programme Management Report
\end{itemize}}\par
\fieldsep
\fieldlabel{KPIs}
{\small \begin{itemize}[noitemsep,topsep=2pt]
\item CAT establishment in first year (Yes/No)
\item No. of CAT meetings per academic year
\item No. of teaching coordination meetings
\item \% subjects with validated syllabus (100\%)
\item Management Report prepared on time (Yes/No)
\item No. of agreed improvement proposals
\end{itemize}}\par
\end{tcolorbox}
\begin{tcolorbox}[riskbox]
{\footnotesize\sffamily\bfseries\color{red!50!black} RISKS}\par\nopagebreak[4]
{\small \begin{itemize}[noitemsep,topsep=2pt]
\item • Misalignments in academic coordination during the first year.
\item • Faculty overload in the initial phase.
\item • Incidents arising from new or interdisciplinary subjects.
\item • Difficulties in adjusting student and faculty expectations during the first years.
\end{itemize}}
\end{tcolorbox}
\begin{tcolorbox}[mitbox]
{\footnotesize\sffamily\bfseries\color{green!40!black} MITIGATION PLAN}\par\nopagebreak[4]
{\small \begin{itemize}[noitemsep,topsep=2pt]
\item • Establish programme coordination from the start, with periodic meetings and continuous monitoring.
\item • Validate syllabi in advance and ensure their overall coherence.
\item • Identify incidents early and activate rapid adjustments.
\item • Maintain close supervision by the Dean's team during the first years of deployment.
\end{itemize}}
\end{tcolorbox}
\vspace{8pt}

\begin{tcolorbox}[pabox={aacolor}{A.P.2.2 \enspace\textbullet\enspace Activation of Academic Management for the Master's in Corporate Responsibility and Sustainability}]

\fieldlabel{Description}
{\small Activate the academic management structure of the Master's Degree in Corporate Responsibility and Sustainability, ensuring its operation in accordance with the approved verification report and its integration into UPV's institutional system. 

This action plan includes the establishment and operation of the Programme Academic Committee, the organisation of teaching coordination, syllabus validation and the launch of academic processes associated with the master's programme, including management of the Master's Thesis (TFM) and monitoring of students' academic records. 

The programme's incorporation into the institutional monitoring system will also be ensured through the preparation of the Programme Management Report, aligning it with UPV's quality mechanisms. }\par
\fieldsep
\fieldlabel{Deliverables}
{\small \begin{itemize}[noitemsep,topsep=2pt]
\item CAT established
\item Validated syllabi
\item TFM management system operational
\item Programme Management Report
\item Coordination meetings held
\end{itemize}}\par
\fieldsep
\fieldlabel{KPIs}
{\small \begin{itemize}[noitemsep,topsep=2pt]
\item CAT establishment first year (Yes/No)
\item \% subjects with validated syllabus before start (100\%)
\item Management Report prepared on time (Yes/No)
\item No. of programme coordination meetings
\end{itemize}}\par
\end{tcolorbox}
\begin{tcolorbox}[riskbox]
{\footnotesize\sffamily\bfseries\color{red!50!black} RISKS}\par\nopagebreak[4]
{\small \begin{itemize}[noitemsep,topsep=2pt]
\item • Delays in establishing the CAT or validating course syllabi before the programme starts.
\item • Coordination difficulties among faculty from different departments in a new, interdisciplinary master's.
\item • Bottlenecks in assigning and supervising master's theses (TFM) for the first cohort.
\item • Enrolment below expectations in the first editions.
\end{itemize}}
\end{tcolorbox}
\begin{tcolorbox}[mitbox]
{\footnotesize\sffamily\bfseries\color{green!40!black} MITIGATION PLAN}\par\nopagebreak[4]
{\small \begin{itemize}[noitemsep,topsep=2pt]
\item • Calendar with a deadline for syllabus validation before enrolment opens, supervised by the Vice Dean for Quality.
\item • Coordination meetings from the first semester, with documented agreements.
\item • TFM procedure (topics, supervisors and deadlines) defined before the second semester.
\item • Targeted promotion of the master's in coordination with COMUNICA and the Advisory Board.
\end{itemize}}
\end{tcolorbox}
\vspace{8pt}

\begin{tcolorbox}[pabox={aacolor}{A.P.2.3 \enspace\textbullet\enspace Integration of New Degrees in Governing Bodies and Academic Planning}]

\fieldlabel{Description}
{\small Incorporate the Degree in Business and Digital Technology and the CSR Master's into the School Council and academic planning processes, ensuring their organisational fit with the overall academic offering.}\par
\fieldsep
\fieldlabel{Deliverables}
{\small \begin{itemize}[noitemsep,topsep=2pt]
\item New degrees presented to School Council
\item Evidence of incorporation into academic planning
\item Record of agreements adopted
\end{itemize}}\par
\fieldsep
\fieldlabel{KPIs}
{\small \begin{itemize}[noitemsep,topsep=2pt]
\item Presentation to School Council (Yes/No)
\item No. of sessions where new degrees are discussed
\item Formal implementation agreements (Yes/No)
\item Integration into academic planning (Yes/No)
\end{itemize}}\par
\end{tcolorbox}
\begin{tcolorbox}[riskbox]
{\footnotesize\sffamily\bfseries\color{red!50!black} RISKS}\par\nopagebreak[4]
{\small \begin{itemize}[noitemsep,topsep=2pt]
\item • Excessive duration and density of School Council sessions.
\item • Low participation or involvement of Council members on matters related to new degrees.
\end{itemize}}
\end{tcolorbox}
\begin{tcolorbox}[mitbox]
{\footnotesize\sffamily\bfseries\color{green!40!black} MITIGATION PLAN}\par\nopagebreak[4]
{\small \begin{itemize}[noitemsep,topsep=2pt]
\item • Plan the inclusion of new degrees on the agenda in a structured manner.
\item • Present information in a concise manner oriented towards decision-making.
\item • Ensure repetition of relevant points in different sessions if necessary.
\end{itemize}}
\end{tcolorbox}
\vspace{8pt}

\newpage
\Needspace{14\baselineskip}
\addsec{S.O.3. Manage accreditation renewal processes}
{\small\color{mutedcolor}\itshape Manage the accreditation renewal processes for the School's official degree programmes, in accordance with the deadlines set by external evaluation agencies, with the aim of securing accreditation renewal under favourable conditions. 

This objective includes preparing the programme self-assessment report, collecting and organising evidence, coordinating the parties involved and managing the external evaluation process, ensuring that the programme is appropriately presented in accordance with the established criteria. 

It will also address responses to provisional reports and the incorporation of improvements arising from the process, ensuring accreditation renewal under favourable conditions. }\par\vspace{6pt}

\begin{tcolorbox}[pabox={aacolor}{A.P.3.1 \enspace\textbullet\enspace Preparation of Accreditation Renewal Processes}]

\fieldlabel{Description}
{\small Prepare the accreditation renewal processes for the School's official degree programmes, in accordance with the schedule established by external evaluation agencies, through the preparation of the programme self-assessment report and the systematic collection of evidence. 

This action plan includes coordinating the Programme Academic Directors, organising the required information (approved verification report, indicators, academic results, evidence of quality and improvement), and preparing the re-accreditation self-assessment report in accordance with the criteria established by ANECA/AVAP. 

Internal review of the self-assessment report and evidence before submission will also be ensured, maintaining coherence between programme design, delivery and results, and facilitating a favourable external evaluation. }\par
\fieldsep
\fieldlabel{Deliverables}
{\small \begin{itemize}[noitemsep,topsep=2pt]
\item Re-accreditation self-assessment report by programme
\item Organised evidence dossier
\item Internally validated documentation
\end{itemize}}\par
\fieldsep
\fieldlabel{KPIs}
{\small \begin{itemize}[noitemsep,topsep=2pt]
\item \% programmes submitting self-assessment report on time (100\%)
\item \% programmes with favourable documentary evaluation
\end{itemize}}\par
\end{tcolorbox}
\begin{tcolorbox}[riskbox]
{\footnotesize\sffamily\bfseries\color{red!50!black} RISKS}\par\nopagebreak[4]
{\small \begin{itemize}[noitemsep,topsep=2pt]
\item • Delays in collecting evidence or in preparing the self-assessment report.
\item • Lack of coherence between the information included in the self-assessment report and the evidence provided.
\item • Dependence on multiple agents.
\item • Insufficient advance preparation of the re-accreditation process.
\end{itemize}}
\end{tcolorbox}
\begin{tcolorbox}[mitbox]
{\footnotesize\sffamily\bfseries\color{green!40!black} MITIGATION PLAN}\par\nopagebreak[4]
{\small \begin{itemize}[noitemsep,topsep=2pt]
\item • Establish a detailed timeline with intermediate milestones and defined responsibilities.
\item • Centralise coordination under the Vice Dean for Quality.
\item • Use structured templates and clear criteria.
\item • Conduct internal reviews prior to submission, verifying consistency between evidence and analysis.
\end{itemize}}
\end{tcolorbox}
\vspace{8pt}

\begin{tcolorbox}[pabox={aacolor}{A.P.3.2 \enspace\textbullet\enspace Management of the External Evaluation Process and Accreditation Renewal}]

\fieldlabel{Description}
{\small Manage the external evaluation process for official degree programmes during accreditation renewal, ensuring the proper organisation of the evaluation committee visit, hearings with the various stakeholder groups and an appropriate response to reports issued by the evaluation agency. 

This action plan includes preparing and coordinating hearings (management team, faculty, students, graduates and employers), making evidence available and presenting it during the visit, analysing the provisional report and preparing representations where appropriate. 

The aim is to ensure that the programme's case is presented effectively during the external evaluation process and to obtain a favourable final accreditation renewal report. }\par
\fieldsep
\fieldlabel{Deliverables}
{\small \begin{itemize}[noitemsep,topsep=2pt]
\item Evaluation committee visit organised
\item Hearings held with the various groups
\item Provisional report analysed
\item Representations prepared if applicable
\item Final renewal report obtained
\end{itemize}}\par
\fieldsep
\fieldlabel{KPIs}
{\small \begin{itemize}[noitemsep,topsep=2pt]
\item Committee visit completed (Yes/No)
\item \% groups in hearings (100\%)
\item No. of critical requirements in provisional report
\item \% representations accepted
\item Final result by programme (Favourable/Not favourable)
\end{itemize}}\par
\end{tcolorbox}
\begin{tcolorbox}[riskbox]
{\footnotesize\sffamily\bfseries\color{red!50!black} RISKS}\par\nopagebreak[4]
{\small \begin{itemize}[noitemsep,topsep=2pt]
\item • An unfavourable evaluation or one with significant requirements arising from weaknesses detected in the programme.
\item • Lack of alignment in the narratives during hearings between the different participating groups.
\item • Difficulties in preparing the evaluation committee visit.
\item • Insufficient or weak response to provisional reports.
\end{itemize}}
\end{tcolorbox}
\begin{tcolorbox}[mitbox]
{\footnotesize\sffamily\bfseries\color{green!40!black} MITIGATION PLAN}\par\nopagebreak[4]
{\small \begin{itemize}[noitemsep,topsep=2pt]
\item • Prepare hearings in advance, aligning key messages between the different participants.
\item • Conduct internal simulations of the external evaluation process to identify possible weaknesses.
\item • Coordinate centrally the organisation of the evaluation committee visit.
\item • Prepare solid and well-founded responses to provisional reports, supported by evidence and concrete improvement actions.
\end{itemize}}
\end{tcolorbox}
\vspace{8pt}

\newpage
\Needspace{14\baselineskip}
\addsec{S.O.4. Systematically integrate teaching innovation and artificial intelligence into programmes}
{\small\color{mutedcolor}\itshape Systematically integrate teaching innovation into the School's official degree programmes through the progressive incorporation into subjects of innovative methodologies, tools and approaches developed in educational innovation projects. 

This objective includes the effective transfer of innovation project results (such as the institutional artificial intelligence projects currently being developed in bachelor's programmes) into regular teaching, ensuring their actual application in subjects, coordination at programme level and consolidation over time. 

Faculty participation in educational innovation initiatives and the development of communities of practice will also be promoted, enabling experiences to be shared, their impact to be evaluated, results to be disseminated and practices that add value to student learning to be extended.  }\par\vspace{6pt}

\begin{tcolorbox}[pabox={aacolor}{A.P.4.1 \enspace\textbullet\enspace Development and Application of Teaching Innovation Projects}]

\fieldlabel{Description}
{\small Develop and implement teaching innovation projects (PIME or others) in the programmes, promoting the continuity of initiatives, their coordination at programme level and the transfer of practices among faculty.}\par
\fieldsep
\fieldlabel{Deliverables}
{\small \begin{itemize}[noitemsep,topsep=2pt]
\item Teaching innovation projects developed
\item Subjects with innovative activities implemented
\item Innovative teaching materials generated

\end{itemize}}\par
\fieldsep
\fieldlabel{KPIs}
{\small \begin{itemize}[noitemsep,topsep=2pt]
\item No. of teaching innovation projects developed
\item No. of subjects with innovation activities
\item No. of participating faculty members
\item \% subjects with documented evidence of innovation applied in the classroom
\end{itemize}}\par
\end{tcolorbox}
\begin{tcolorbox}[riskbox]
{\footnotesize\sffamily\bfseries\color{red!50!black} RISKS}\par\nopagebreak[4]
{\small \begin{itemize}[noitemsep,topsep=2pt]
\item • Limited faculty participation.
\item • Occasional or non-sustained application of innovative initiatives.
\item • Lack of systematic impact evaluation.
\item • Dependence on individual initiatives without sufficient coordination at programme level.
\end{itemize}}
\end{tcolorbox}
\begin{tcolorbox}[mitbox]
{\footnotesize\sffamily\bfseries\color{green!40!black} MITIGATION PLAN}\par\nopagebreak[4]
{\small \begin{itemize}[noitemsep,topsep=2pt]
\item • Foster faculty participation through institutional recognition and alignment with official calls.
\item • Promote innovation projects adapted to the characteristics of each initiative.
\item • Incorporate impact evaluation as a mandatory element.
\item • Promote coordination within each programme.
\end{itemize}}
\end{tcolorbox}
\vspace{8pt}

\begin{tcolorbox}[pabox={aacolor}{A.P.4.2 \enspace\textbullet\enspace Promotion and Dissemination of FADE Teaching Innovation}]

\fieldlabel{Description}
{\small Organise teaching innovation conferences at FADE with scientific standards (scientific committee, proceedings with ISBN from an academic publisher), participate in academic events through institutional parallel sessions and promote faculty participation in external conferences.}\par
\fieldsep
\fieldlabel{Deliverables}
{\small \begin{itemize}[noitemsep,topsep=2pt]
\item FADE teaching innovation conferences (every 2 years) with ISBN
\item Teaching contributions at internal and external conferences
\item Institutional parallel sessions organised at external forums
\end{itemize}}\par
\fieldsep
\fieldlabel{KPIs}
{\small \begin{itemize}[noitemsep,topsep=2pt]
\item No. of teaching innovation conferences organised at FADE
\item No. of parallel sessions organised by the management team
\item No. of contributions presented at FADE conferences
\item No. of FADE faculty participating in conferences
\item No. of faculty participating in parallel sessions
\end{itemize}}\par
\end{tcolorbox}
\begin{tcolorbox}[riskbox]
{\footnotesize\sffamily\bfseries\color{red!50!black} RISKS}\par\nopagebreak[4]
{\small \begin{itemize}[noitemsep,topsep=2pt]
\item • Low faculty participation in conferences, parallel sessions or teaching congresses.
\item • Difficulty in consolidating a stable and periodic dynamic for disseminating teaching innovation.
\item • Insufficient production of teaching contributions with sufficient academic quality.
\item • Excessive dependence on the initiative of the Dean's team.
\end{itemize}}
\end{tcolorbox}
\begin{tcolorbox}[mitbox]
{\footnotesize\sffamily\bfseries\color{green!40!black} MITIGATION PLAN}\par\nopagebreak[4]
{\small \begin{itemize}[noitemsep,topsep=2pt]
\item • Plan conferences and dissemination actions well in advance and integrate them into the academic dynamic of the School.
\item • Incentivise faculty participation through institutional recognition and support for presenting contributions.
\item • Establish a stable organisational structure with a scientific and organising committee.
\item • Promote the generation and maturation of teaching contributions from already developed innovation projects.
\end{itemize}}
\end{tcolorbox}
\vspace{8pt}

\begin{tcolorbox}[pabox={aacolor}{A.P.4.3 \enspace\textbullet\enspace Development of AI Literacy for Faculty}]

\fieldlabel{Description}
{\small Design and implementation of a structured AI competency development programme for FADE faculty, aimed at advancing the AI literacy and professional development dimensions of the \textit{Ten Dimension AI Readiness Framework} (Digital Education Council, 2025). The programme will combine training in the ethical and responsible use of AI, specific pedagogical applications by discipline and forums for exchange among faculty to transfer good practice.}\par
\fieldsep
\fieldlabel{Deliverables}
{\small \begin{itemize}[noitemsep,topsep=2pt]
\item Diagnosis of initial AI competencies of FADE faculty
\item Structured training programme by levels (basic, applied, advanced)
\item Evidence of participation and learning achievement
\item Annual report on faculty AI maturity progress
\end{itemize}}\par
\fieldsep
\fieldlabel{KPIs}
{\small \begin{itemize}[noitemsep,topsep=2pt]
\item \% of faculty completing basic AI training (target: $\geq$70\% in 2028/29)
\item No. of training actions carried out annually
\item Faculty AI maturity level (DEC scale: Emerging $\rightarrow$ Developing $\rightarrow$ Established)
\item No. of good practices of AI use in teaching documented and shared
\end{itemize}}\par
\end{tcolorbox}
\begin{tcolorbox}[riskbox]
{\footnotesize\sffamily\bfseries\color{red!50!black} RISKS}\par\nopagebreak[4]
{\small \begin{itemize}[noitemsep,topsep=2pt]
\item • Low voluntary participation of faculty in training actions.
\item • High heterogeneity in the starting point among faculty.
\item • Risk that training does not translate into real changes in teaching practice.
\end{itemize}}
\end{tcolorbox}
\begin{tcolorbox}[mitbox]
{\footnotesize\sffamily\bfseries\color{green!40!black} MITIGATION PLAN}\par\nopagebreak[4]
{\small \begin{itemize}[noitemsep,topsep=2pt]
\item • Modular and flexible design that allows different levels of entry and dedication.
\item • Linking the programme to existing teaching innovation projects (PIME).
\item • Institutional recognition of participation and good practices generated.
\end{itemize}}
\end{tcolorbox}
\vspace{8pt}



\newpage
\Needspace{14\baselineskip}
\addsec{S.O.5. Drive the permanent updating of the academic offering}
{\small\color{mutedcolor}\itshape Promote the progressive updating of the School's official degree programmes through the review of their content, subjects and academic structure, in response to identified needs and changes in the environment. 

This objective includes specific analyses of different elements of the academic offering, such as elective subjects, compulsory subject content or possible overlaps between subjects, to identify opportunities for improvement. 

The systematic use of artificial intelligence tools to support the continuous review of teaching activity and curricula will also be promoted. These tools will enable the analysis of subject syllabi, competencies, learning outcomes, assessment systems, teaching materials and academic evidence, facilitating the early detection of redundancies, gaps in education, curricular misalignments and opportunities for methodological innovation. AI will act as a mechanism supporting academic decision-making, complementing the expert judgement of faculty and programme coordination bodies.

Proposals arising from these analyses will, where feasible, be incorporated into programme academic coordination processes, fostering continuous updating of the academic offering while respecting existing regulatory and organisational constraints. 

Communication of analysis results and decisions taken by academic bodies will also be promoted, helping faculty to understand, monitor and implement them. }\par\vspace{6pt}

\begin{tcolorbox}[pabox={aacolor}{A.P.5.1 \enspace\textbullet\enspace Review and Proposal for Updating the Academic Offering}]

\fieldlabel{Description}
{\small Periodically analyse elements of the academic offering (elective subjects, content, overlaps) to identify update opportunities and formulate proposals assessed in the CATs, facilitating continuous updating.}\par
\fieldsep
\fieldlabel{Deliverables}
{\small \begin{itemize}[noitemsep,topsep=2pt]
\item Analysis reports of elements of the academic offering
\item Documented identification of review needs
\item Evidence of agreements adopted in the CATs
\item Evidence of academic actions taken
\end{itemize}}\par
\fieldsep
\fieldlabel{KPIs}
{\small \begin{itemize}[noitemsep,topsep=2pt]
\item No. of analyses performed on the academic offering
\item No. of evidence-based academic decisions adopted
\item \% decisions implemented
\item No. of subjects modified after analysis
\item No. of programmes with at least one improvement applied
\end{itemize}}\par
\end{tcolorbox}
\begin{tcolorbox}[riskbox]
{\footnotesize\sffamily\bfseries\color{red!50!black} RISKS}\par\nopagebreak[4]
{\small \begin{itemize}[noitemsep,topsep=2pt]
\item • Identification of a large number of review needs without real viability for action.
\item • Lack of availability or quality of the information used in analyses.
\item • Uneven involvement between programmes or teaching teams.
\item • Identified gaps not translating into real changes due to lack of response from the teaching staff responsible for the course.
\end{itemize}}
\end{tcolorbox}
\begin{tcolorbox}[mitbox]
{\footnotesize\sffamily\bfseries\color{green!40!black} MITIGATION PLAN}\par\nopagebreak[4]
{\small \begin{itemize}[noitemsep,topsep=2pt]
\item • Prioritise viable academic actions within the existing regulatory framework.
\item • Focus analyses on specific and relevant aspects.
\item • Use verified and available information sources.
\item • Integrate this process into the usual dynamic of the Programme Academic Committees.
\item • Establish a clear follow-up protocol in the CAT for identified gaps, with assigned owners and defined response deadlines.
\end{itemize}}
\end{tcolorbox}
\vspace{8pt}

\Needspace{18\baselineskip}
\section{Timeline}

\small

\begin{tabularx}{\textwidth}{>{\bfseries}l >{\raggedright\arraybackslash}X *{4}{>{\centering\arraybackslash}p{1.5cm}}}
\toprule
A.P. & Action plan & 2026/27 & 2027/28 & 2028/29 & 2029/30 \\
\midrule
A.P.1.1 & Diagnosis and Analysis of Adaptation to RD 822/2021 & $\blacksquare$ & $\blacksquare$ & $\blacksquare$ &  \\
A.P.1.2 & Redesign of study plans and verification reports & $\blacksquare$ & $\blacksquare$ & $\blacksquare$ &  \\
A.P.2.1 & Activation of Academic Management for the Degree in Business and Digital Technology & $\blacksquare$ & $\blacksquare$ & $\blacksquare$ & $\blacksquare$ \\
A.P.2.2 & Activation of Academic Management for the Master's in Corporate Responsibility and Sustainability & $\blacksquare$ & $\blacksquare$ & $\blacksquare$ & $\blacksquare$ \\
A.P.2.3 & Integration of new programmes into the School's governing bodies and academic planning & $\blacksquare$ & $\blacksquare$ &  &  \\
A.P.3.1 & Preparation of programme accreditation renewal processes &  & $\blacksquare$ & $\blacksquare$ & $\blacksquare$ \\
A.P.3.2 & Management of the External Evaluation Process and Accreditation Renewal &  & $\blacksquare$ & $\blacksquare$ & $\blacksquare$ \\
A.P.4.1 & Development and application of teaching innovation projects in programmes & $\blacksquare$ & $\blacksquare$ & $\blacksquare$ & $\blacksquare$ \\
A.P.4.2 & Promotion and Dissemination of FADE Teaching Innovation & $\blacksquare$ & $\blacksquare$ & $\blacksquare$ & $\blacksquare$ \\
A.P.4.3 & Development of AI Literacy for Faculty & $\blacksquare$ & $\blacksquare$ & $\blacksquare$ & $\blacksquare$ \\
A.P.5.1 & Review and Proposal for Updating the Academic Offering &  & $\blacksquare$ & $\blacksquare$ & $\blacksquare$ \\
\bottomrule
\end{tabularx}

\Needspace{18\baselineskip}
\section{Summary of Objective Achievement}

\small
Achievement of the project's objectives will be evaluated through the monitoring of indicators defined in each action plan, as well as an overall assessment of the degree of adaptation, implementation, monitoring and improvement of the School's programmes. The approach combines quantitative indicators associated with process execution with a qualitative assessment of the degree to which those processes have become established.

\textbf{S.O.1. Adapt official degree programmes to RD 822/2021.} It will be considered achieved when the School has completed the adaptation of its programmes to the regulatory framework, with updated and approved verification reports aligned with the requirements of Royal Decree 822/2021, and an academic structure that is coherent and validated by the evaluation agencies.

\textbf{S.O.2. Implement and consolidate the new strategic academic offering.} It will be considered achieved when the new programmes are fully implemented and have an operational academic management structure, with active Programme Academic Committees, established teaching coordination mechanisms and monitoring processes integrated into UPV's quality system.

\textbf{S.O.3. Manage programme accreditation renewal processes.} It will be considered achieved when programmes have successfully completed re-accreditation processes in accordance with the planned schedule, obtaining favourable reports and demonstrating coherence between their design, delivery and results.

\textbf{S.O.4. Systematically integrate teaching innovation and artificial intelligence into programmes.} It will be considered consolidated when teaching innovation is continuously integrated into programmes, with projects transferred to the classroom, periodic dissemination events and faculty making sustained progress in their competencies in AI applied to teaching, in line with the maturity objectives defined for faculty.

\textbf{S.O.5. Drive the permanent updating of the academic offering.} It will be considered achieved when the School has a systematic, evidence-based mechanism for reviewing its academic offering, enabling it to identify updating needs and carry out feasible academic actions that improve programme quality and relevance.


\clearpage

\chapter[University Experience 360]{\textcolor{e360color}{University Experience 360}}
\begin{tcolorbox}[descbox]
\small Consolidation of a management model centred on the comprehensive student experience throughout their university journey. The project coherently articulates support, communication, participation and university life, inside and outside the classroom, to reinforce the sense of belonging and offer a more human, approachable, participatory experience connected with the School.
\end{tcolorbox}
\Needspace{18\baselineskip}
\section{Project Description}

\small
The ``University Experience 360'' project aims to consolidate a management model centred on the holistic student experience at the School of Management (FADE), complementing academic education with a more human, participatory and approachable perspective, and providing students with continuous support throughout their university journey.

The project builds on the work already undertaken by the Vice Dean for Student Affairs, particularly through the launch of the PIAE+ programme, and proposes its development towards a more integrated, structured and visible model that avoids fragmented initiatives and brings coherence to the full range of activities directed at students.

Under the ``Experience 360'' concept, the project responds to three identified structural challenges: improving two-way communication between students and the School, which currently lacks fluidity; consolidating support beyond the first year; and strengthening the sense of community and belonging. To this end, it comprises a set of coordinated actions that integrate communication, support and university life as interrelated elements.

Methodologically, the project is based on four cross-cutting action mechanisms: Processes, Data, Communication and Artificial Intelligence. First, the Processes mechanism focuses on improving and structuring student assistance, guidance and monitoring systems, ensuring coherent support. Second, the Data mechanism facilitates the use of institutional information to identify needs and assess the impact of activities. Third, the Communication mechanism acts as a key tool for facilitating participation, accessibility and two-way connection with students. Artificial Intelligence (AI) joins this ecosystem as a fourth cross-cutting action mechanism, supporting information analysis, facilitating the identification of academic risk situations and supporting decision-making in student support. Its use is envisaged as progressive and focused on specific cases, without replacing academic judgement.

The project also takes student diversity into account in the design of activities and recognises the value of co-curricular and extracurricular activities as part of a university experience that extends education beyond the classroom. These activities aim to promote the development of transversal competencies, students' personal and professional growth, their well-being and a more complete experience of their time at the School.

Overall, the aim is to advance towards a more complete, coherent and inclusive university experience that strengthens students' academic and personal development and consolidates a more active, connected university community.

\Needspace{18\baselineskip}
\section{Relationship with the Strategic Lines}

\small
The project's strategic contribution is distributed as follows:

\begin{itemize}[noitemsep,topsep=2pt]
\item \textbf{Academic Excellence $\rightarrow$ 50\%}
\item \textbf{Impact $\rightarrow$ 50\%}
\end{itemize}
This distribution reflects the project's cross-cutting nature, at the intersection of improving students' educational experience and strengthening their connection with the School.

The contribution to Academic Excellence is based on:
\begin{itemize}[noitemsep,topsep=2pt]
\item Strengthening support throughout the student's journey.
\item Improving welcome, guidance, monitoring and support processes to promote a more coherent, structured educational experience.
\item Developing transversal competencies and supporting students' personal and professional growth by integrating and recognising co-curricular and extracurricular activities, extending education beyond the classroom as part of a more complete university experience.
\item Adapting activities to students' diverse profiles and needs.
\end{itemize}

The contribution to the Impact line is reflected in:
\begin{itemize}[noitemsep,topsep=2pt]
\item Improving two-way communication with students.
\item Strengthening the sense of community and belonging.
\item Promoting a more participatory, approachable and visible university experience.
\item Strengthening students' involvement in School life to help build a more active, connected community that identifies with FADE.
\end{itemize}

\newpage
\Needspace{18\baselineskip}
\section{Objectives, Action Plans and Development}
\Needspace{14\baselineskip}
\addsec{S.O.1. Strengthen two-way, transparent and direct communication with students}
{\small\color{mutedcolor}\itshape Enhance transparent, direct, two-way communication between the Dean's team and students through their various representation channels. Improve communication channels and practices with students, facilitating access to information, the collection of feedback and direct interaction between the Dean's team and the various student representation channels. }\par\vspace{6pt}

\begin{tcolorbox}[pabox={e360color}{A.P.1.1 \enspace\textbullet\enspace Organisation of Communication Channels with Students}]

\fieldlabel{Description}
{\small Review and organisation of the communication system, identifying existing channels, defining their purpose and establishing clear usage criteria to improve coherence, avoid duplication, promote the use of formal channels and adapt to student communication habits.}\par
\fieldsep
\fieldlabel{Deliverables}
{\small \begin{itemize}[noitemsep,topsep=2pt]
\item Inventory of communication channels and their purpose
\item Document of channel usage and management criteria
\item Analysis of adaptation to student habits
\end{itemize}}\par
\fieldsep
\fieldlabel{KPIs}
{\small \begin{itemize}[noitemsep,topsep=2pt]
\item No. of channels analysed
\item Criteria document defined and validated (Yes/No)
\item No. of channels adapted to identified habits
\item Average rating (1--5) of perceived clarity (survey/focus group)
\end{itemize}}\par
\end{tcolorbox}
\begin{tcolorbox}[riskbox]
{\footnotesize\sffamily\bfseries\color{red!50!black} RISKS}\par\nopagebreak[4]
{\small \begin{itemize}[noitemsep,topsep=2pt]
\item • Duplication of communication channels, generating confusion among students.
\item • Receipt of the same query through multiple channels and by different persons responsible, making management difficult.
\item • Low effectiveness of channels used to reach students.
\end{itemize}}
\end{tcolorbox}
\begin{tcolorbox}[mitbox]
{\footnotesize\sffamily\bfseries\color{green!40!black} MITIGATION PLAN}\par\nopagebreak[4]
{\small \begin{itemize}[noitemsep,topsep=2pt]
\item • Clear definition of official communication channels, their purpose and their use.
\item • Establishment of criteria on what type of queries are managed in each channel.
\item • Explicit communication to students of the enabled channels and their uses.
\item • Definition of channels based on the needs and proposals raised by students.
\end{itemize}}
\end{tcolorbox}
\vspace{8pt}

\begin{tcolorbox}[pabox={e360color}{A.P.1.2 \enspace\textbullet\enspace Development of Communication Actions at Key Moments of the Student Journey}]

\fieldlabel{Description}
{\small Design and development of structured communication activities at key stages of the student journey, from entry to the School to completion of studies, with the aim of facilitating academic decision-making through clear, timely information and reducing student uncertainty at key moments. Regular forums for dialogue and active listening between the Dean's team and students, particularly through their representatives, will also be promoted to collect concerns, identify needs and strengthen two-way communication.}\par
\fieldsep
\fieldlabel{Deliverables}
{\small \begin{itemize}[noitemsep,topsep=2pt]
\item Informational sessions (welcome, enrolment, pathways, Bachelor's Degree Final Project (TFG)…)
\item Informational materials and unified resource guide
\item Structured calendar of key milestones
\item Integration of the PIAE+ programme in the miUPV app
\item Record of meetings with student representatives
\end{itemize}}\par
\fieldsep
\fieldlabel{KPIs}
{\small \begin{itemize}[noitemsep,topsep=2pt]
\item No. of communication actions carried out
\item No. of material views
\item Milestones calendar defined and implemented (Yes/No)
\item \% students contacting their tutor via app
\end{itemize}}\par
\end{tcolorbox}
\begin{tcolorbox}[riskbox]
{\footnotesize\sffamily\bfseries\color{red!50!black} RISKS}\par\nopagebreak[4]
{\small \begin{itemize}[noitemsep,topsep=2pt]
\item • Information overload or dispersion.
\item • Low utilisation of materials.
\item • Information perceived as not very relevant by students.
\item • Lack of adaptation of communication formats to student habits.
\end{itemize}}
\end{tcolorbox}
\begin{tcolorbox}[mitbox]
{\footnotesize\sffamily\bfseries\color{green!40!black} MITIGATION PLAN}\par\nopagebreak[4]
{\small \begin{itemize}[noitemsep,topsep=2pt]
\item • Advance planning of actions.
\item • Adaptation of content to student needs.
\item • Adaptation of formats to student communication habits.
\item • Recording and subsequent dissemination of sessions.
\end{itemize}}
\end{tcolorbox}
\vspace{8pt}

\newpage
\Needspace{14\baselineskip}
\addsec{S.O.2. Consolidate a model of continuous student support}
{\small\color{mutedcolor}\itshape Maintain, improve and expand welcome, guidance and monitoring programmes, developing PIAE+ towards a more integrated, structured and coherent model that supports students at every stage of their journey through the School. }\par\vspace{6pt}

\begin{tcolorbox}[pabox={e360color}{A.P.2.1 \enspace\textbullet\enspace Organisation and Coordination of the Student Support System (PIAE+)}]

\fieldlabel{Description}
{\small Organisation, coordination and improvement of the student support system through the PIAE+ programme, adapting it to the School's context and ensuring consistent application. This includes coordination between programmes and Schools in the case of double degrees, as well as links with other university services and programmes related to student support, including those addressing student diversity and the appropriate management of the information needed for monitoring.  }\par
\fieldsep
\fieldlabel{Deliverables}
{\small \begin{itemize}[noitemsep,topsep=2pt]
\item PIAE+ organisation and coordination document
\item Common guidelines for PIAE+ tutors
\item Coordination mechanisms with other Schools
\end{itemize}}\par
\fieldsep
\fieldlabel{KPIs}
{\small \begin{itemize}[noitemsep,topsep=2pt]
\item PIAE+ organisation document defined (Yes/No)
\item No. of coordination meetings with PIAE+ tutors
\item No. of coordination actions with other Schools
\end{itemize}}\par
\end{tcolorbox}
\begin{tcolorbox}[riskbox]
{\footnotesize\sffamily\bfseries\color{red!50!black} RISKS}\par\nopagebreak[4]
{\small \begin{itemize}[noitemsep,topsep=2pt]
\item • Lack of effective coordination between the various agents involved.
\item • Differences in programme application between programmes or Schools.
\item • Low involvement of some tutors.
\end{itemize}}
\end{tcolorbox}
\begin{tcolorbox}[mitbox]
{\footnotesize\sffamily\bfseries\color{green!40!black} MITIGATION PLAN}\par\nopagebreak[4]
{\small \begin{itemize}[noitemsep,topsep=2pt]
\item • Establishment of periodic coordination spaces.
\item • Clear definition of roles and responsibilities.
\item • Periodic monitoring of programme functioning.
\end{itemize}}
\end{tcolorbox}
\vspace{8pt}

\begin{tcolorbox}[pabox={e360color}{A.P.2.2 \enspace\textbullet\enspace Monitoring and Intervention with Students Throughout Their Journey}]

\fieldlabel{Description}
{\small Development of student monitoring and intervention activities throughout their academic journey, with the aim of detecting risk situations early, supporting students at critical moments and promoting continuation in their studies. Particular attention is paid to milestones such as the start of studies, academic progress and the development of the Bachelor's Degree Final Project (TFG), incorporating support adapted to students' different needs and profiles. To this end, the plan also provides for the definition and use of information needed to analyse, where appropriate, the possible influence of diversity variables on low performance or academic risk, always within institutional criteria for data access, processing and custody. }\par
\fieldsep
\fieldlabel{Deliverables}
{\small \begin{itemize}[noitemsep,topsep=2pt]
\item Early academic risk identification system
\item Record of monitoring actions via PIAE+
\item Record of students making insufficient progress on their Bachelor's Degree Final Project (TFG)
\item Dropout risk action protocol
\item Analysis of dropout variables and preventive actions
\end{itemize}}\par
\fieldsep
\fieldlabel{KPIs}
{\small \begin{itemize}[noitemsep,topsep=2pt]
\item No. of students identified as at risk
\item Average detection time from course start
\item No. of students contacted via PIAE+
\item No. of students who start/resume the Bachelor's Degree Final Project (TFG) after intervention
\item No. of preventive actions designed from analysis
\item No. of at-risk students who continue in the following year
\item \% low performance situations with identified diversity variables
\end{itemize}}\par
\end{tcolorbox}
\begin{tcolorbox}[riskbox]
{\footnotesize\sffamily\bfseries\color{red!50!black} RISKS}\par\nopagebreak[4]
{\small \begin{itemize}[noitemsep,topsep=2pt]
\item • Late detection of academic risk situations.
\item • Low student response to monitoring actions.
\item • Difficulty in translating data analysis into concrete actions.
\item • Limitations in data availability or quality.
\item • Difficulty in obtaining adequate information on diversity variables.
\end{itemize}}
\end{tcolorbox}
\begin{tcolorbox}[mitbox]
{\footnotesize\sffamily\bfseries\color{green!40!black} MITIGATION PLAN}\par\nopagebreak[4]
{\small \begin{itemize}[noitemsep,topsep=2pt]
\item • Progressive improvement of data access and analysis systems.
\item • Reinforcement of the active role of PIAE+ tutors in intervention.
\item • Design of concrete actions from data analysis.
\item • Adaptation of actions to student characteristics.
\item • Coordination with specialised university services for diversity support.
\item • Prior definition of the necessary information and its access and processing conditions.
\end{itemize}}
\end{tcolorbox}
\vspace{8pt}

\newpage
\Needspace{14\baselineskip}
\addsec{S.O.3. Foster the sense of belonging, participation and identification with the School}
{\small\color{mutedcolor}\itshape Drive active student participation, reinforce their identification with the School and generate initiatives that favour interaction and community building.}\par\vspace{6pt}

\begin{tcolorbox}[pabox={e360color}{A.P.3.1 \enspace\textbullet\enspace Building Community and Sense of Belonging at the School}]

\fieldlabel{Description}
{\small Development of initiatives aimed at generating a stronger sense of student community and belonging to the School, fostering active involvement in university life. The aim is to combine support for initiatives promoted by the Student Union with activities led by the Dean's team that encourage interaction among students, faculty and alumni, creating a more approachable, participatory and connected environment. Particular attention will be paid to significant moments in university life, such as graduation, and to activities that strengthen student involvement and identification with the School. These initiatives may include meetings with graduates, promotion of Alumni UPV mentoring programmes or alumni participation in educational and career guidance activities that encourage interaction between students and graduates. }\par
\fieldsep
\fieldlabel{Deliverables}
{\small \begin{itemize}[noitemsep,topsep=2pt]
\item Annual community-building activities plan
\item Record of activities, events and meetings with institutional presence
\item Annual report of community initiatives
\item Evidence of activity dissemination
\end{itemize}}\par
\fieldsep
\fieldlabel{KPIs}
{\small \begin{itemize}[noitemsep,topsep=2pt]
\item No. of community activities and initiatives carried out
\item No. of activities promoted by students with School support
\item Level of participation (average attendance)
\item No. of initiatives involving student--faculty--graduate interaction
\item \% students rating their satisfaction ≥4/5 (survey)
\end{itemize}}\par
\end{tcolorbox}
\begin{tcolorbox}[riskbox]
{\footnotesize\sffamily\bfseries\color{red!50!black} RISKS}\par\nopagebreak[4]
{\small \begin{itemize}[noitemsep,topsep=2pt]
\item • Low student participation in activities.
\item • Lack of coordination between Student Union and Dean's team.
\item • Concentration of activities with low real impact.
\item • Scarce student involvement in School-driven initiatives.
\end{itemize}}
\end{tcolorbox}
\begin{tcolorbox}[mitbox]
{\footnotesize\sffamily\bfseries\color{green!40!black} MITIGATION PLAN}\par\nopagebreak[4]
{\small \begin{itemize}[noitemsep,topsep=2pt]
\item • Periodic coordination with the Student Union to align initiatives.
\item • Prioritisation of activities with greater student interest.
\item • Reinforcement of activity dissemination through School channels.
\item • Periodic evaluation of participation and adjustment of initiatives.
\end{itemize}}
\end{tcolorbox}
\vspace{8pt}

\newpage
\Needspace{14\baselineskip}
\addsec{S.O.4. Give visibility to the work of the Vice Dean for Student Affairs and student representatives}
{\small\color{mutedcolor}\itshape Give greater visibility to the actions, initiatives and results of the Vice Dean for Student Affairs and student representation, improving their communication and recognition within the School.}\par\vspace{6pt}

\begin{tcolorbox}[pabox={e360color}{A.P.4.1 \enspace\textbullet\enspace Showcasing and Recognising the Work of the Student Area and Students}]

\fieldlabel{Description}
{\small Visibility and recognition of the work of the student affairs area and of students themselves, highlighting their initiatives, achievements and contributions both within the School and across the university and society. The aim is to give greater visibility to student activities, results and journeys, including recognition of academic merit, participation in relevant initiatives and contributions with social impact, thereby strengthening their identification with the School and the value of their university experience. }\par
\fieldsep
\fieldlabel{Deliverables}
{\small \begin{itemize}[noitemsep,topsep=2pt]
\item Communication plan for area initiatives and achievements
\item List of student achievements and recognitions (academic, social)
\item Evidence of dissemination
\end{itemize}}\par
\fieldsep
\fieldlabel{KPIs}
{\small \begin{itemize}[noitemsep,topsep=2pt]
\item No. of initiatives, achievements or activities disseminated
\item No. of recognition actions carried out
\item No. of achievements and recognitions identified and recorded
\end{itemize}}\par
\end{tcolorbox}
\begin{tcolorbox}[riskbox]
{\footnotesize\sffamily\bfseries\color{red!50!black} RISKS}\par\nopagebreak[4]
{\small \begin{itemize}[noitemsep,topsep=2pt]
\item • Low effective visibility of actions despite their dissemination.
\item • Lack of systematic identification or collection of student achievements.
\item • Little student interest in the disseminated content.
\end{itemize}}
\end{tcolorbox}
\begin{tcolorbox}[mitbox]
{\footnotesize\sffamily\bfseries\color{green!40!black} MITIGATION PLAN}\par\nopagebreak[4]
{\small \begin{itemize}[noitemsep,topsep=2pt]
\item • Establishment of mechanisms to identify and collect student achievements.
\item • Prioritisation of relevant and close content for students.
\item • Combined use of different channels and formats to improve dissemination reach.
\end{itemize}}
\end{tcolorbox}
\vspace{8pt}

\newpage
\Needspace{14\baselineskip}
\addsec{S.O.5. Integrate education beyond the classroom as part of the university experience}
{\small\color{mutedcolor}\itshape Promote co-curricular and extracurricular activities that contribute to the development of transversal competencies, personal growth and a more complete university experience, strengthening the educational value of students' time at the School. This objective encompasses initiatives complementing formal education that contribute to students' personal and professional development, such as competency development activities, meetings with professionals and graduates, mentoring programmes, participation in challenges, institutional, cultural or solidarity activities, and other educational experiences promoted by the School or partner organisations.}\par\vspace{6pt}

\begin{tcolorbox}[pabox={e360color}{A.P.5.1 \enspace\textbullet\enspace Promotion of Education Beyond the Classroom}]

\fieldlabel{Description}
{\small Promotion and organisation of co-curricular and extracurricular activities that enrich students' education beyond the classroom, integrating them into a more complete, holistic university experience. The plan seeks to identify, support, manage and bring coherence to initiatives that foster the acquisition of transversal competencies, personal growth and more active student participation in their own education. }\par
\fieldsep
\fieldlabel{Deliverables}
{\small \begin{itemize}[noitemsep,topsep=2pt]
\item Annual plan of competency development actions beyond the classroom
\item Evidence of communication of organised initiatives
\item Evidence of student participation in complementary activities
\end{itemize}}\par
\fieldsep
\fieldlabel{KPIs}
{\small \begin{itemize}[noitemsep,topsep=2pt]
\item No. of initiatives beyond the classroom promoted annually
\item No. of communication actions about these initiatives
\item No. of students participating in complementary educational activities
\item \% of students rating the educational value at 4/5 or above (survey)
\end{itemize}}\par
\end{tcolorbox}
\begin{tcolorbox}[riskbox]
{\footnotesize\sffamily\bfseries\color{red!50!black} RISKS}\par\nopagebreak[4]
{\small \begin{itemize}[noitemsep,topsep=2pt]
\item • Difficulty in clearly identifying which activities genuinely provide educational value.
\item • Low student participation in proposed activities.
\item • Perception that these activities do not provide educational value.
\item • Difficulty in having information on all activities faculty carry out at the School.
\end{itemize}}
\end{tcolorbox}
\begin{tcolorbox}[mitbox]
{\footnotesize\sffamily\bfseries\color{green!40!black} MITIGATION PLAN}\par\nopagebreak[4]
{\small \begin{itemize}[noitemsep,topsep=2pt]
\item • Definition of basic criteria to identify the educational value of activities and prioritisation of those aligned with competency development.
\item • Adaptation of initiatives to student interests and reinforcement of their dissemination.
\item • Clear communication of the educational value of activities and collection of student feedback.
\item • Establishment of simple coordination and information collection mechanisms with faculty and academic managers.
\end{itemize}}
\end{tcolorbox}
\vspace{8pt}

\clearpage
\Needspace{18\baselineskip}
\section{Timeline}

\small

\begin{tabularx}{\textwidth}{>{\bfseries}l >{\raggedright\arraybackslash}X *{8}{>{\centering\arraybackslash}p{1.15cm}}}
\toprule
A.P. & 26/27 S1 & 26/27 S2 & 27/28 S1 & 27/28 S2 & 28/29 S1 & 28/29 S2 & 29/30 S1 & 29/30 S2 \\
\midrule
A.P.1.1  &  & $\blacksquare$ & $\blacksquare$ &  &  &  &  &  \\
A.P.1.2  & $\blacksquare$ & $\blacksquare$ & $\blacksquare$ & $\blacksquare$ & $\blacksquare$ & $\blacksquare$ & $\blacksquare$ & $\blacksquare$ \\
A.P.2.1  &  & $\blacksquare$ & $\blacksquare$ & $\blacksquare$ &  &  &  &  \\
A.P.2.2  &  & $\blacksquare$ &  & $\blacksquare$ &  & $\blacksquare$ &  & $\blacksquare$ \\
A.P.3.1 &  & $\blacksquare$ & $\blacksquare$ & $\blacksquare$ & $\blacksquare$ & $\blacksquare$ & $\blacksquare$ & $\blacksquare$ \\
A.P.4.1  &  &  & $\blacksquare$ & $\blacksquare$ & $\blacksquare$ & $\blacksquare$ & $\blacksquare$ & $\blacksquare$ \\
A.P.5.1 &   &  &  & $\blacksquare$ & $\blacksquare$ & $\blacksquare$ & $\blacksquare$ & $\blacksquare$ \\
\bottomrule
\end{tabularx}

\Needspace{18\baselineskip}
\section{Summary of Objective Achievement}

\small
Achievement of the project's objectives will be evaluated through the monitoring of indicators defined in each action plan, as well as an overall assessment of improvements in the student experience at the School. The project's approach combines quantitative indicators (activity, participation, reach) with qualitative indicators based on student perceptions.

In relation to \textbf{Objective 1}, the aim is to improve the clarity, coherence and effectiveness of communication with students, facilitating both access to information and direct interaction with the Dean's team. Achievement of this objective will be reflected in:
\begin{itemize}[noitemsep,topsep=2pt]
\item Greater use of the defined channels.
\item More positive student assessments of communication.
\end{itemize}

For \textbf{Objective 2}, the impact will focus on consolidating a continuous support model that enables earlier detection of risk situations and a more agile, personalised response. The expected outcomes are:
\begin{itemize}[noitemsep,topsep=2pt]
\item Improved student monitoring.
\item Greater activation of support mechanisms.
\item Positive progress in indicators such as continuation in studies or academic progress.
\end{itemize}

\textbf{Objective 3} aims to strengthen students' sense of belonging and participation in School life. Its achievement will be measured through:
\begin{itemize}[noitemsep,topsep=2pt]
\item Increased participation in activities.
\item Student involvement in their own initiatives.
\item Improved perceptions of university life and the connection with the School.
\end{itemize}

\textbf{Objective 4} seeks to increase the visibility and recognition of the work of the student affairs area and student achievements. The expected outcomes are:
\begin{itemize}[noitemsep,topsep=2pt]
\item Wider dissemination of initiatives and results.
\item Increased student awareness and appreciation of these activities.
\end{itemize}

Finally, \textbf{Objective 5} will advance the integration of education beyond the classroom into students' university experience. The expected outcomes are:
\begin{itemize}[noitemsep,topsep=2pt]
\item Greater identification and recognition of co-curricular and extracurricular activities.
\item A positive perception of their contribution to personal and competency development.
\item A more complete university experience within the School.
\end{itemize}

Overall, the project will help achieve a more complete, coherent and student-centred university experience, strengthening students' academic and personal development and their connection with the School.

\chapter[FADE Spaces]{\textcolor{espcolor}{FADE Spaces}}
\begin{tcolorbox}[descbox]
\small Addresses the comprehensive improvement of School infrastructure. It reorganises classrooms and timetables to optimise their use and adapt them to new active learning methodologies. It transforms shared areas into functional and sustainable environments that favour interaction and collaborative work. In addition, it renews the main hall to create a representative, modern and welcoming space, equipped with a new visual identity and digital displays, projecting an image of innovation and reinforcing the sense of community.
\end{tcolorbox}

\Needspace{18\baselineskip}
\section{Project Description}

\small
The project \textbf{FADE Spaces} aims to achieve a comprehensive, progressive transformation of the physical spaces of the UPV School of Management, adapting them to new academic, organisational and external engagement needs. It is a strategic initiative intended to optimise the use of facilities, improve the university experience and strengthen the School's institutional identity.

The project addresses three integrated areas of action: renovation of the entrance hall as a space for access and institutional representation, academic reorganisation of classrooms and timetables, and improvement of the School's common spaces. This approach enables coherent action across the physical environment, avoiding isolated interventions and ensuring a comprehensive vision aligned with the development of the academic offering and the needs of the university community.

First, renovating the entrance hall transforms the main access space into a functional, representative environment aligned with the institution's identity, improving orientation, communication and the experience of students, faculty and visitors. Second, comprehensive reorganisation of classrooms and timetables optimises facility use, facilitates the introduction of new programmes and teaching methodologies, and improves academic planning. Third, improving common spaces seeks to create more functional, accessible and attractive environments that encourage interaction, collaborative work and a sense of belonging.

The project is supported by the cross-cutting action mechanisms of Processes, Data, Communication and Artificial Intelligence. Processes enables the redefinition of space organisation and use, as well as academic planning; Data facilitates the analysis of flows, needs and utilisation levels; Communication helps improve signage, visual identity and the visibility of School activity; and Artificial Intelligence is incorporated as a tool supporting the analysis and optimisation of space and resource use. All of this is envisaged through a progressive approach focused on improving operational efficiency and the quality of the university experience.

The project also incorporates sustainability criteria in space management, material selection and resource efficiency, contributing to responsible management aligned with the School's institutional values.

\Needspace{18\baselineskip}
\section{Relationship with the Strategic Lines}

\small
The FADE Spaces project's strategic contribution is distributed as follows:

\begin{itemize}[noitemsep,topsep=2pt]
\item \textbf{Academic Excellence → 55\%}
\item \textbf{Impact → 45\%}
\end{itemize}

This distribution reflects the project's nature as an initiative focused primarily on improving the academic environment and the conditions in which teaching takes place, with significant effects also on institutional visibility and organisational efficiency.

The contribution to Academic Excellence is based on improving the physical conditions in which teaching and university life take place, facilitating the implementation of new teaching methodologies, optimising academic planning and creating more functional, accessible environments aligned with the current needs of students and faculty.

The contribution to Impact is justified by the strategic role of spaces in institutional representation, user experience and stakeholder engagement. Improving spaces reinforces the School's image, facilitates the reception of visitors and external evaluators, and helps present a modern, efficient organisation aligned with values such as sustainability, innovation and transparency.
\newpage
\Needspace{18\baselineskip}
\section{Objectives, Action Plans and Development}
\Needspace{14\baselineskip}
\addsec{S.O.1. Reconfigure the hall as a functional and optimised space for daily use}
{\small\color{mutedcolor}\itshape Redesign the layout and organisation of the hall to improve access flows, circulation and functionality, ensuring efficient use adapted to actual needs.}\par\vspace{6pt}

\begin{tcolorbox}[pabox={espcolor}{A.P.1.1 \enspace\textbullet\enspace Analysis of Hall Space Usage}]

\fieldlabel{Description}
{\small Analysis of transit flows, current uses and hall user needs, with the aim of defining an optimal functional model.}\par
\fieldsep
\fieldlabel{Deliverables}
{\small \begin{itemize}[noitemsep,topsep=2pt]
\item Hall usage analysis report
\item Identification of needs by user type
\item Functional zoning proposal
\end{itemize}}\par
\fieldsep
\fieldlabel{KPIs}
{\small \begin{itemize}[noitemsep,topsep=2pt]
\item Analysis document prepared (Yes/No)
\item No. of users/groups consulted
\item Validation of the functional model by Dean's team (Yes/No)
\end{itemize}}\par
\end{tcolorbox}
\begin{tcolorbox}[riskbox]
{\footnotesize\sffamily\bfseries\color{red!50!black} RISKS}\par\nopagebreak[4]
{\small \begin{itemize}[noitemsep,topsep=2pt]
\item • Incomplete or unrepresentative analysis of users.
\item • Lack of consensus in defining the model.
\end{itemize}}
\end{tcolorbox}
\begin{tcolorbox}[mitbox]
{\footnotesize\sffamily\bfseries\color{green!40!black} MITIGATION PLAN}\par\nopagebreak[4]
{\small \begin{itemize}[noitemsep,topsep=2pt]
\item • Include different groups in the analysis.
\item • Formal validation of the functional model before proceeding.
\end{itemize}}
\end{tcolorbox}
\vspace{8pt}

\begin{tcolorbox}[pabox={espcolor}{A.P.1.2 \enspace\textbullet\enspace Design and Implementation of the Hall Layout and Seating Areas}]

\fieldlabel{Description}
{\small Detailed design of the space (layout, zoning, furniture, circulation and seating and informal meeting areas) and implementation of the new hall configuration according to the defined functional model.}\par
\fieldsep
\fieldlabel{Deliverables}
{\small \begin{itemize}[noitemsep,topsep=2pt]
\item Hall layout plan
\item Furniture, seating areas and functional elements proposal
\item Space reorganised and implemented
\end{itemize}}\par
\fieldsep
\fieldlabel{KPIs}
{\small \begin{itemize}[noitemsep,topsep=2pt]
\item Design approved (Yes/No)
\item \% of redesign execution
\item No. of seating areas created
\item User satisfaction level
\end{itemize}}\par
\end{tcolorbox}
\begin{tcolorbox}[riskbox]
{\footnotesize\sffamily\bfseries\color{red!50!black} RISKS}\par\nopagebreak[4]
{\small \begin{itemize}[noitemsep,topsep=2pt]
\item • Physical space limitations.
\item • Deviations in cost or execution times.
\item • Low use of the seating areas created.
\item • Inappropriate use of the space.
\end{itemize}}
\end{tcolorbox}
\begin{tcolorbox}[mitbox]
{\footnotesize\sffamily\bfseries\color{green!40!black} MITIGATION PLAN}\par\nopagebreak[4]
{\small \begin{itemize}[noitemsep,topsep=2pt]
\item • Prior technical validation of the design.
\item • Execution in phases if necessary.
\item • Design based on detected actual needs.
\item • Signage and basic usage rules.
\end{itemize}}
\end{tcolorbox}
\vspace{8pt}

\newpage
\Needspace{14\baselineskip}
\addsec{S.O.2. Set up a work room in the hall}
{\small\color{mutedcolor}\itshape Refurbish the hall's computer room, currently underused, adapting its furniture and technological infrastructure to students' new ways of working.}\par\vspace{6pt}

\begin{tcolorbox}[pabox={espcolor}{A.P.2.1 \enspace\textbullet\enspace Refurbishment of the Hall's Work Room}]

\fieldlabel{Description}
{\small Changes in students' working patterns and use of facilities mean that the hall's computer room has a low level of use and effectiveness. In order to encourage greater use and utility, a change in furniture and technological infrastructure is proposed for the room.}\par
\fieldsep
\fieldlabel{Deliverables}
{\small \begin{itemize}[noitemsep,topsep=2pt]
\item Initial study of room usage and needs
\item Work room refurbished and implemented
\end{itemize}}\par
\fieldsep
\fieldlabel{KPIs}
{\small \begin{itemize}[noitemsep,topsep=2pt]
\item \% of execution
\item Student satisfaction level
\end{itemize}}\par
\end{tcolorbox}
\begin{tcolorbox}[riskbox]
{\footnotesize\sffamily\bfseries\color{red!50!black} RISKS}\par\nopagebreak[4]
{\small Persistent low use despite refurbishment.}
\end{tcolorbox}
\begin{tcolorbox}[mitbox]
{\footnotesize\sffamily\bfseries\color{green!40!black} MITIGATION PLAN}\par\nopagebreak[4]
{\small Initial study of actual needs and promotion of the new space among students.}
\end{tcolorbox}
\vspace{8pt}

\newpage
\Needspace{14\baselineskip}
\addsec{S.O.3. Reinforce the institutional image and visual identity of FADE in the hall}
{\small\color{mutedcolor}\itshape Provide the hall with a coherent, modern aesthetic aligned with institutional values, and integrate communication media that disseminate academic activity.}\par\vspace{6pt}

\begin{tcolorbox}[pabox={espcolor}{A.P.3.1 \enspace\textbullet\enspace Implementation of Identity and Signage}]

\fieldlabel{Description}
{\small Design and implementation of a signage system coherent with the institutional identity of FADE and the UPV, including branding and visual communication elements.}\par
\fieldsep
\fieldlabel{Deliverables}
{\small \begin{itemize}[noitemsep,topsep=2pt]
\item Signage system implemented
\item Branding elements installed
\item Visually coherent space
\end{itemize}}\par
\fieldsep
\fieldlabel{KPIs}
{\small \begin{itemize}[noitemsep,topsep=2pt]
\item \% of elements implemented
\item Improvement in image perception (survey)
\end{itemize}}\par
\end{tcolorbox}
\begin{tcolorbox}[riskbox]
{\footnotesize\sffamily\bfseries\color{red!50!black} RISKS}\par\nopagebreak[4]
{\small \begin{itemize}[noitemsep,topsep=2pt]
\item • Visual saturation.
\item • Lack of coherence between elements.
\end{itemize}}
\end{tcolorbox}
\begin{tcolorbox}[mitbox]
{\footnotesize\sffamily\bfseries\color{green!40!black} MITIGATION PLAN}\par\nopagebreak[4]
{\small \begin{itemize}[noitemsep,topsep=2pt]
\item • Application of defined visual guide.
\item • Supervision of the final ensemble.
\end{itemize}}
\end{tcolorbox}
\vspace{8pt}

\begin{tcolorbox}[pabox={espcolor}{A.P.3.2 \enspace\textbullet\enspace Digital Communication Displays}]

\fieldlabel{Description}
{\small Installation and management of screens and digital displays in the hall to disseminate the academic, research and institutional activity of FADE.}\par
\fieldsep
\fieldlabel{Deliverables}
{\small \begin{itemize}[noitemsep,topsep=2pt]
\item Screens installed and operational
\item Content management system defined
\item First content published
\end{itemize}}\par
\fieldsep
\fieldlabel{KPIs}
{\small \begin{itemize}[noitemsep,topsep=2pt]
\item No. of displays installed
\item No. of content items disseminated monthly
\item Update frequency
\end{itemize}}\par
\end{tcolorbox}
\begin{tcolorbox}[riskbox]
{\footnotesize\sffamily\bfseries\color{red!50!black} RISKS}\par\nopagebreak[4]
{\small \begin{itemize}[noitemsep,topsep=2pt]
\item • Lack of content updating.
\item • Technical problems.
\end{itemize}}
\end{tcolorbox}
\begin{tcolorbox}[mitbox]
{\footnotesize\sffamily\bfseries\color{green!40!black} MITIGATION PLAN}\par\nopagebreak[4]
{\small \begin{itemize}[noitemsep,topsep=2pt]
\item • Assignment of content responsibilities.
\item • Periodic technical maintenance.
\end{itemize}}
\end{tcolorbox}
\vspace{8pt}

\newpage
\Needspace{14\baselineskip}
\addsec{S.O.4. Improve the assignment and use of classrooms and teaching spaces}
{\small\color{mutedcolor}\itshape Ensure that teaching facilities are used efficiently, reducing overlaps, under-utilisation and mismatches between the type of classroom and pedagogical needs.}\par\vspace{6pt}

\begin{tcolorbox}[pabox={espcolor}{A.P.4.1 \enspace\textbullet\enspace Diagnosis and Analysis of Space and Timetable Usage}]

\fieldlabel{Description}
{\small Analysis of classroom occupancy flows, current uses and needs of degree programmes, identifying opportunities for improvement in the assignment and use of teaching spaces.}\par
\fieldsep
\fieldlabel{Deliverables}
{\small \begin{itemize}[noitemsep,topsep=2pt]
\item Classroom occupancy diagnostic report
\item Space and time slot utilisation map
\item Identification of improvement opportunities
\end{itemize}}\par
\fieldsep
\fieldlabel{KPIs}
{\small \begin{itemize}[noitemsep,topsep=2pt]
\item Average classroom occupancy rate
\item \% of under-utilised classrooms
\item No. of time slots with deficit in current planning
\end{itemize}}\par
\end{tcolorbox}
\begin{tcolorbox}[riskbox]
{\footnotesize\sffamily\bfseries\color{red!50!black} RISKS}\par\nopagebreak[4]
{\small \begin{itemize}[noitemsep,topsep=2pt]
\item • Lack of complete data.
\item • Limitations of current systems for extracting detailed information.
\item • Coordination between programmes.
\end{itemize}}
\end{tcolorbox}
\begin{tcolorbox}[mitbox]
{\footnotesize\sffamily\bfseries\color{green!40!black} MITIGATION PLAN}\par\nopagebreak[4]
{\small Collaboration with programme managers and professional services staff (PTGAS) to ensure information quality.}
\end{tcolorbox}
\vspace{8pt}

\newpage
\Needspace{14\baselineskip}
\addsec{S.O.5. Design and implement a new model of academic planning}
{\small\color{mutedcolor}\itshape Develop a more flexible and efficient planning model that optimises the assignment of timetables and classrooms, facilitating coordination between programmes and the incorporation of new methodologies.}\par\vspace{6pt}

\begin{tcolorbox}[pabox={espcolor}{A.P.5.1 \enspace\textbullet\enspace Design of the New Academic Planning Model}]

\fieldlabel{Description}
{\small Design and implementation of a more efficient, flexible teaching planning model aligned with the evolution of the School's academic offering, including the operational protocol for its annual application.}\par
\fieldsep
\fieldlabel{Deliverables}
{\small \begin{itemize}[noitemsep,topsep=2pt]
\item Approved teaching planning model
\item Operational protocol for annual timetable planning
\end{itemize}}\par
\fieldsep
\fieldlabel{KPIs}
{\small \begin{itemize}[noitemsep,topsep=2pt]
\item Reduction in the No. of timetable conflicts
\item Improvement in classroom occupancy index
\item No. of requests processed and successfully resolved
\end{itemize}}\par
\end{tcolorbox}
\begin{tcolorbox}[riskbox]
{\footnotesize\sffamily\bfseries\color{red!50!black} RISKS}\par\nopagebreak[4]
{\small \begin{itemize}[noitemsep,topsep=2pt]
\item • Organisational complexity in coordination between programmes and departments.
\item • Limitations of existing computer systems.
\item • Changes in the academic offering or in the number of students that alter initial assumptions.
\item • Resistance to change.
\end{itemize}}
\end{tcolorbox}
\begin{tcolorbox}[mitbox]
{\footnotesize\sffamily\bfseries\color{green!40!black} MITIGATION PLAN}\par\nopagebreak[4]
{\small \begin{itemize}[noitemsep,topsep=2pt]
\item • Participatory design process involving managers and administrative staff.
\item • Development of an iterative timetable adjustment process.
\item • Inform managers in advance about planning and changes made.
\end{itemize}}
\end{tcolorbox}
\vspace{8pt}

\newpage
\Needspace{14\baselineskip}
\addsec{S.O.6. Progressively adapt teaching spaces to pedagogical needs}
{\small\color{mutedcolor}\itshape Ensure that available spaces support the methodologies of new degrees (flexible classrooms, collaborative work, technological equipment).}\par\vspace{6pt}

\begin{tcolorbox}[pabox={espcolor}{A.P.6.1 \enspace\textbullet\enspace Functional Adaptation of Teaching Spaces}]

\fieldlabel{Description}
{\small Diagnosis of the pedagogical adequacy of classrooms and a prioritised action plan to adapt spaces to active methodologies, collaborative work and the technological needs of new programmes.}\par
\fieldsep
\fieldlabel{Deliverables}
{\small \begin{itemize}[noitemsep,topsep=2pt]
\item Diagnosis of pedagogical adequacy of classrooms
\item Classification of space typologies
\item Functional adaptation plan with prioritised actions
\end{itemize}}\par
\fieldsep
\fieldlabel{KPIs}
{\small \begin{itemize}[noitemsep,topsep=2pt]
\item No. of classrooms adapted or improved
\item \% spaces aligned with methodological needs
\item No. of incidents due to space inadequacy
\end{itemize}}\par
\end{tcolorbox}
\begin{tcolorbox}[riskbox]
{\footnotesize\sffamily\bfseries\color{red!50!black} RISKS}\par\nopagebreak[4]
{\small \begin{itemize}[noitemsep,topsep=2pt]
\item • Budget limitations for carrying out space adaptations.
\item • Physical or structural restrictions of existing facilities.
\item • Interference with teaching activity during adaptation works.
\end{itemize}}
\end{tcolorbox}
\begin{tcolorbox}[mitbox]
{\footnotesize\sffamily\bfseries\color{green!40!black} MITIGATION PLAN}\par\nopagebreak[4]
{\small \begin{itemize}[noitemsep,topsep=2pt]
\item • Prioritisation of actions according to academic impact and technical feasibility.
\item • Staggered planning of adaptations to minimise interference with teaching activity.
\item • Coordination with university infrastructure services to assess the technical feasibility of interventions.
\end{itemize}}
\end{tcolorbox}
\vspace{8pt}

\newpage
\Needspace{14\baselineskip}
\addsec{S.O.7. Optimise the use and management of common areas}
{\small\color{mutedcolor}\itshape Improve the organisation, assignment and use of common areas, ensuring efficient, coordinated use adapted to the actual needs of the various groups.}\par\vspace{6pt}

\begin{tcolorbox}[pabox={espcolor}{A.P.7.1 \enspace\textbullet\enspace Diagnosis and Analysis of Common Area Usage}]

\fieldlabel{Description}
{\small Analysis of the usage patterns, occupancy and needs of the various School common areas (circulation areas, seating areas, signage) to identify improvement opportunities.}\par
\fieldsep
\fieldlabel{Deliverables}
{\small \begin{itemize}[noitemsep,topsep=2pt]
\item Common area usage and occupancy report
\item Needs map by group
\end{itemize}}\par
\fieldsep
\fieldlabel{KPIs}
{\small \begin{itemize}[noitemsep,topsep=2pt]
\item Report prepared (Yes/No)
\item No. of spaces analysed
\end{itemize}}\par
\end{tcolorbox}
\begin{tcolorbox}[riskbox]
{\footnotesize\sffamily\bfseries\color{red!50!black} RISKS}\par\nopagebreak[4]
{\small Incomplete data.}
\end{tcolorbox}
\begin{tcolorbox}[mitbox]
{\footnotesize\sffamily\bfseries\color{green!40!black} MITIGATION PLAN}\par\nopagebreak[4]
{\small Include multiple sources and groups.}
\end{tcolorbox}
\vspace{8pt}

\begin{tcolorbox}[pabox={espcolor}{A.P.7.2 \enspace\textbullet\enspace Implementation of a Management and Optimisation System for Common Areas}]

\fieldlabel{Description}
{\small Design and implementation of a coordinated management model for common areas, including usage regulations, sustainability criteria and an occupancy monitoring system.}\par
\fieldsep
\fieldlabel{Deliverables}
{\small \begin{itemize}[noitemsep,topsep=2pt]
\item Common area management model defined
\item Usage regulations approved
\item Occupancy monitoring system
\end{itemize}}\par
\fieldsep
\fieldlabel{KPIs}
{\small \begin{itemize}[noitemsep,topsep=2pt]
\item Model approved (Yes/No)
\item \% improvement in perceived occupancy
\item No. of incidents resolved
\end{itemize}}\par
\end{tcolorbox}
\begin{tcolorbox}[riskbox]
{\footnotesize\sffamily\bfseries\color{red!50!black} RISKS}\par\nopagebreak[4]
{\small Resistance to change.}
\end{tcolorbox}
\begin{tcolorbox}[mitbox]
{\footnotesize\sffamily\bfseries\color{green!40!black} MITIGATION PLAN}\par\nopagebreak[4]
{\small Communication and training.}
\end{tcolorbox}
\vspace{8pt}

\newpage
\Needspace{14\baselineskip}
\addsec{S.O.8. Improve the user experience and foster interaction in common areas}
{\small\color{mutedcolor}\itshape Transform common areas into accessible, comfortable and attractive environments that favour interaction, collaborative work and sense of belonging, incorporating sustainability criteria.}\par\vspace{6pt}

\begin{tcolorbox}[pabox={espcolor}{A.P.8.1 \enspace\textbullet\enspace Design and Creation of Interaction and Collaborative Work Zones}]

\fieldlabel{Description}
{\small Design of seating areas, collaborative work and informal interaction zones in common areas, fostering university life and reinforcing the institutional identity of the School.}\par
\fieldsep
\fieldlabel{Deliverables}
{\small \begin{itemize}[noitemsep,topsep=2pt]
\item Design of interaction zones and work spaces
\item Spaces enabled (meeting points, collaborative areas)
\item Sustainability criteria incorporated in materials and furniture
\end{itemize}}\par
\fieldsep
\fieldlabel{KPIs}
{\small \begin{itemize}[noitemsep,topsep=2pt]
\item No. of spaces/zones created or renovated
\item Level of use (periodic observation)
\item User satisfaction level (\%)
\end{itemize}}\par
\end{tcolorbox}
\begin{tcolorbox}[riskbox]
{\footnotesize\sffamily\bfseries\color{red!50!black} RISKS}\par\nopagebreak[4]
{\small Low use.}
\end{tcolorbox}
\begin{tcolorbox}[mitbox]
{\footnotesize\sffamily\bfseries\color{green!40!black} MITIGATION PLAN}\par\nopagebreak[4]
{\small Design based on needs.}
\end{tcolorbox}
\vspace{8pt}


\Needspace{18\baselineskip}
\section{Timeline}

\small

\begin{tabularx}{\textwidth}{>{\bfseries}l *{8}{>{\centering\arraybackslash}p{1.15cm}}}
\toprule
A.P. & 26/27 S1 & 26/27 S2 & 27/28 S1 & 27/28 S2 & 28/29 S1 & 28/29 S2 & 29/30 S1 & 29/30 S2 \\
\midrule
A.P.1.1 & $\blacksquare$ & $\blacksquare$ &  &  &  &  &  &  \\
A.P.1.2 &  & $\blacksquare$ & $\blacksquare$ & $\blacksquare$ &  &  &  &  \\
A.P.2.1 & $\blacksquare$ & $\blacksquare$ & $\blacksquare$ &  &  &  &  &  \\
A.P.3.1 & $\blacksquare$ &  &  &  &  &  &  &  \\
A.P.3.2 &  &  & $\blacksquare$ &  &  &  &  &  \\
A.P.4.1 & $\blacksquare$ & $\blacksquare$ &  &  &  &  &  &  \\
A.P.5.1 &  & $\blacksquare$ & $\blacksquare$ &  &  &  &  &  \\
A.P.6.1 &  &  & $\blacksquare$ & $\blacksquare$ & $\blacksquare$ & $\blacksquare$ & $\blacksquare$ & $\blacksquare$ \\
A.P.7.1 & $\blacksquare$ & $\blacksquare$ &  &  &  &  &  &  \\
A.P.7.2 &  & $\blacksquare$ & $\blacksquare$ & $\blacksquare$ &  &  &  &  \\
A.P.8.1 &  &  & $\blacksquare$ & $\blacksquare$ &  &  &  &  \\
\bottomrule
\end{tabularx}

\Needspace{18\baselineskip}
\section{Summary of Objective Achievement}

\small
For the purposes of monitoring the Strategic Plan, achievement of the project's objectives will be assessed against the following criteria:

\begin{itemize}

\item \textbf{S.O.1. Reconfigure the hall as a functional and optimised space for daily use.}  
It is considered achieved when the entrance hall has a clearly defined functional layout, with optimised access and circulation flows, areas differentiated by use (circulation, seating, information) and a space arrangement that enables efficient, orderly use adapted to the needs of the School's various groups.

\item \textbf{S.O.2. Set up a work room in the hall.}
It will be considered achieved when the work room in the entrance hall has been renovated and is in effective use, with a high level of student satisfaction.

\item \textbf{S.O.3. Strengthen FADE's institutional image and visual identity (entrance hall).}  
It is considered achieved when the entrance hall presents a coherent, modern image aligned with institutional values, with a clearly recognisable visual identity, consistent signage and an improved perception of the School both internally and among external visitors, and facilitates the visibility of relevant content disseminated through physical and digital media.

\item \textbf{S.O.4. Improve the assignment and use of classrooms and teaching spaces.}  
It is considered achieved when teaching spaces are used highly efficiently and in a balanced way, based on precise, systematic knowledge of occupancy patterns, reducing underutilisation and mismatches between capacity and demand.

\item \textbf{S.O.5. Design and implement a new model of academic planning.}  
It is considered achieved when there is a fully implemented, structured and shared planning model that enables more agile, coordinated and efficient timetabling, reducing conflicts in classroom allocation and improving the fit between spaces and teaching needs.

\item \textbf{S.O.6. Progressively adapt teaching spaces to pedagogical needs.}  
It is considered achieved when a significant proportion of classrooms have been transformed or equipped to facilitate active methodologies, collaborative work and technology use, improving the learning experience and aligning infrastructure with the academic offering.

\item \textbf{S.O.7. Optimise the use and management of common areas.}  
It is considered achieved when the School has a structured space management system, based on usage and needs data, enabling efficient allocation and reducing conflicts over use.

\item \textbf{S.O.8. Improve the user experience and foster interaction in common areas.}  
It is considered achieved when common spaces enable flexible, multipurpose use, facilitating collaborative work, active learning and high utilisation in line with the needs of the university community.

\end{itemize}
\clearpage

\chapter[FADE Global]{\textcolor{fgcolor}{FADE Global}}
\begin{tcolorbox}[descbox]
\small Improves access to the academic exchange offering through interactive digital tools, such as a destinations map and artificial intelligence assistance for resolving frequently asked questions. It systematises recognition regulations and centralises student support through Policonsulta, reducing administrative barriers. It fosters balanced mobility by adjusting the offering to demand. In addition, it promotes double degree agreements and staff mobility to reinforce the international reputation of the School.
\end{tcolorbox}

\Needspace{18\baselineskip}
\section{Project Description}

\small
The project \textbf{FADE Global} is a strategic initiative aimed at providing students with clear, up-to-date and understandable access to the School's international academic exchange opportunities. The project is structured around three complementary areas: access to information, improvement of student services and the strategic development of internationalisation. Its purpose is to reduce the information and administrative barriers that hinder destination choice and understanding of academic recognition criteria, through the development of digital tools, the structuring of management processes and the improvement of communication channels.

The project is based on the synergistic integration of four cross-cutting action mechanisms: Processes, Data, Communication and Artificial Intelligence. First, the Processes mechanism ensures the systematisation of recognition regulations and the structuring of mobility procedures, guaranteeing an efficient and consistent service. Second, the Data mechanism provides traceability and quality in the information available on destinations and agreements, facilitating informed academic decisions. Third, the Communication mechanism centralises and clarifies information, improving academic guidance and transparency for students. Artificial Intelligence (AI) joins this ecosystem as a fourth cross-cutting action mechanism, acting as the project's interactive assistance layer. While Processes formalises rules and Data provides information, AI acts as a support tool through conversational assistants that answer frequently asked questions swiftly, improving information accessibility and reducing the administrative workload.

Through this initiative, FADE strengthens its international positioning, increases the visibility of its academic offering and projects an image of institutional modernisation, while contributing to academic excellence by improving student guidance, planning and academic recognition processes. The digital infrastructure and methodologies adopted will provide lasting institutional capacity to manage and communicate international mobility effectively, ensuring the sustainability and impact of activities in the medium and long term.

\Needspace{18\baselineskip}
\section{Relationship with the Strategic Lines}

\small
The project \textbf{FADE Global: Improving International Mobility} distributes its strategic contribution as follows:

\begin{itemize}[noitemsep,topsep=2pt]
\item \textbf{Academic Excellence → 60\%}
\item \textbf{Impact → 40\%}
\end{itemize}

The 60\% contribution to Academic Excellence is based on the project's direct improvement of student guidance and support processes, its facilitation of informed academic decisions in planning international pathways, and its assurance of clear, consistent academic recognition. The implementation of interactive digital tools and AI-based support systems strengthens the quality of the academic experience and mobility management, ensuring that learning objectives and programme coherence are achieved efficiently and traceably.

The 40\% contribution to Impact is justified by the international visibility generated by the initiative. The visibility of mobility opportunities, the consolidation of strategic agreements with partner universities and the integration of double degree programmes help strengthen FADE's global reputation, increase its competitiveness relative to other institutions and project the School's academic excellence internationally. This institutional effect expands FADE's presence in global higher education networks and strengthens its positioning as an innovative, inclusive academic institution serving students and society.



\small
The project contributes to Academic Excellence through:

\begin{itemize}[noitemsep,topsep=2pt]
\item Improved academic guidance and support for students participating in international mobility.
\item Facilitating understanding of academic recognition procedures and assessment criteria.
\item Implementation of digital tools for structured access to information.
\item Traceability and continuous updating of international academic content.
\end{itemize}



\small
The contribution to the Impact line is reflected in:

\begin{itemize}[noitemsep,topsep=2pt]
\item FADE's international visibility through strategic agreements and double degree programmes.
\item Increased visibility and reputation of the School globally.
\item Strengthening students' international experience and employability.
\item Institutional dissemination of academic excellence through digital channels and communication tools.
\end{itemize}
\newpage
\Needspace{18\baselineskip}
\section{Objectives, Action Plans and Development}
\Needspace{14\baselineskip}
\addsec{S.O.1. Improve access, understanding and usability of information on academic exchange}
{\small\color{mutedcolor}\itshape Provide students with clear, up-to-date and understandable access to FADE's international mobility opportunities, reducing information barriers in destination selection and understanding of academic recognition conditions, through interactive digital tools and AI-based support systems. This activity reinforces the processes + analysis + communication mechanisms, ensuring that international mobility is manageable and accessible in a structured manner. }\par\vspace{6pt}

\begin{tcolorbox}[pabox={fgcolor}{A.P.1.1 \enspace\textbullet\enspace Development of an Interactive Map with Available Destinations}]

\fieldlabel{Description}
{\small Design and implementation of an interactive digital map that allows visualising all academic exchange destinations, incorporating information on partner universities, places, language, requirements and basic characteristics of each destination.}\par
\fieldsep
\fieldlabel{Deliverables}
{\small \begin{itemize}[noitemsep,topsep=2pt]
\item FADE interactive digital map of international destinations
\item Basic information sheets per destination
\item Integration into institutional web environment
\end{itemize}}\par
\fieldsep
\fieldlabel{KPIs}
{\small \begin{itemize}[noitemsep,topsep=2pt]
\item Tool published and operational (Yes/No)
\item \% active destinations incorporated into the map
\end{itemize}}\par
\end{tcolorbox}
\begin{tcolorbox}[riskbox]
{\footnotesize\sffamily\bfseries\color{red!50!black} RISKS}\par\nopagebreak[4]
{\small \begin{itemize}[noitemsep,topsep=2pt]
\item • Incomplete or outdated information.
\item • Low initial utilisation.
\end{itemize}}
\end{tcolorbox}
\begin{tcolorbox}[mitbox]
{\footnotesize\sffamily\bfseries\color{green!40!black} MITIGATION PLAN}\par\nopagebreak[4]
{\small \begin{itemize}[noitemsep,topsep=2pt]
\item • Annual content review.
\item • Active dissemination at FADE conferences and channels.
\end{itemize}}
\end{tcolorbox}
\vspace{8pt}

\begin{tcolorbox}[pabox={fgcolor}{A.P.1.2 \enspace\textbullet\enspace Development of an AI Chatbot for Students and Proposal of Recognition Rules}]

\fieldlabel{Description}
{\small Development of an AI-based conversational assistant that answers frequently asked questions about mobility, administrative procedures and general academic recognition criteria, with a validated knowledge base, initially limited to frequently asked questions and with a knowledge base validated by academic leads. It will never replace personalised academic assistance but will act as a complement. }\par
\fieldsep
\fieldlabel{Deliverables}
{\small \begin{itemize}[noitemsep,topsep=2pt]
\item Operational AI chatbot integrated in FADE environment
\item Periodic content update system
\end{itemize}}\par
\fieldsep
\fieldlabel{KPIs}
{\small \begin{itemize}[noitemsep,topsep=2pt]
\item Chatbot implemented and operational (Yes/No)
\item No. of queries handled automatically
\end{itemize}}\par
\end{tcolorbox}
\begin{tcolorbox}[riskbox]
{\footnotesize\sffamily\bfseries\color{red!50!black} RISKS}\par\nopagebreak[4]
{\small \begin{itemize}[noitemsep,topsep=2pt]
\item • Incorrect or outdated responses.
\item • Low initial student trust.
\end{itemize}}
\end{tcolorbox}
\begin{tcolorbox}[mitbox]
{\footnotesize\sffamily\bfseries\color{green!40!black} MITIGATION PLAN}\par\nopagebreak[4]
{\small \begin{itemize}[noitemsep,topsep=2pt]
\item • Validation by academic managers.
\item • Periodic updating.
\item • Initially limited to frequent queries.
\end{itemize}}
\end{tcolorbox}
\vspace{8pt}

\newpage
\Needspace{14\baselineskip}
\addsec{S.O.2. Improve support and accompaniment for outgoing students}
{\small\color{mutedcolor}\itshape Strengthen the quality, traceability and consistency of assistance to outgoing students through structured digital communication tools and updated regulations, facilitating understanding of academic recognition procedures and reducing students' administrative uncertainty. }\par\vspace{6pt}

\begin{tcolorbox}[pabox={fgcolor}{A.P.2.1 \enspace\textbullet\enspace Implementation of Policonsulta as a Communication Tool}]

\fieldlabel{Description}
{\small Centralise the management of outgoing student enquiries through the Policonsulta system, ensuring traceability and efficiency, avoiding the dispersion of enquiries across multiple channels and improving information traceability. }\par
\fieldsep
\fieldlabel{Deliverables}
{\small \begin{itemize}[noitemsep,topsep=2pt]
\item Policonsulta channel configured
\item Usage protocol
\item Student manual
\item Frequently asked questions record
\end{itemize}}\par
\fieldsep
\fieldlabel{KPIs}
{\small \begin{itemize}[noitemsep,topsep=2pt]
\item Channel operational (Yes/No)
\item \% queries handled
\item Average response time
\end{itemize}}\par
\end{tcolorbox}
\begin{tcolorbox}[riskbox]
{\footnotesize\sffamily\bfseries\color{red!50!black} RISKS}\par\nopagebreak[4]
{\small \begin{itemize}[noitemsep,topsep=2pt]
\item • Initial student resistance to changing channels.
\item • Duplication of queries due to simultaneous use of email and Policonsulta.
\end{itemize}}
\end{tcolorbox}
\begin{tcolorbox}[mitbox]
{\footnotesize\sffamily\bfseries\color{green!40!black} MITIGATION PLAN}\par\nopagebreak[4]
{\small \begin{itemize}[noitemsep,topsep=2pt]
\item • Clear communication of the single official channel.
\item • Initial transition period with automatic reminders.
\end{itemize}}
\end{tcolorbox}
\vspace{8pt}

\begin{tcolorbox}[pabox={fgcolor}{A.P.2.2 \enspace\textbullet\enspace Regulatory Update on Recognition and Collection}]

\fieldlabel{Description}
{\small Systematise the academic recognition regulations applicable to international mobility in a single operational document, simplified and digitally accessible.}\par
\fieldsep
\fieldlabel{Deliverables}
{\small \begin{itemize}[noitemsep,topsep=2pt]
\item Updated single document of recognition regulations
\item Digital version accessible in institutional environment
\end{itemize}}\par
\fieldsep
\fieldlabel{KPIs}
{\small \begin{itemize}[noitemsep,topsep=2pt]
\item Document approved and published (Yes/No)
\item Reduction of interpretative incidents in recognition
\end{itemize}}\par
\end{tcolorbox}
\begin{tcolorbox}[riskbox]
{\footnotesize\sffamily\bfseries\color{red!50!black} RISKS}\par\nopagebreak[4]
{\small \begin{itemize}[noitemsep,topsep=2pt]
\item • Difficulty in homogenising criteria between programmes.
\item • Subsequent changes in UPV regulations.
\end{itemize}}
\end{tcolorbox}
\begin{tcolorbox}[mitbox]
{\footnotesize\sffamily\bfseries\color{green!40!black} MITIGATION PLAN}\par\nopagebreak[4]
{\small \begin{itemize}[noitemsep,topsep=2pt]
\item • Prior academic validation.
\item • Annual document review.
\end{itemize}}
\end{tcolorbox}
\vspace{8pt}

\newpage
\Needspace{14\baselineskip}
\addsec{S.O.3. Facilitate the academic integration of incoming international students}
{\small\color{mutedcolor}\itshape Promote the effective integration of international students through accessible academic organisation consistent with their educational needs, strengthening their academic and social integration and improving the perceived quality of FADE's international offering. }\par\vspace{6pt}

\begin{tcolorbox}[pabox={fgcolor}{A.P.3.1 \enspace\textbullet\enspace Assign a Coordinated Timetable for Electives in English}]

\fieldlabel{Description}
{\small Coordinate the timetable offering of elective subjects in English so that incoming students can enrol in pathways without overlaps, with advance publication of the offering.}\par
\fieldsep
\fieldlabel{Deliverables}
{\small \begin{itemize}[noitemsep,topsep=2pt]
\item Proposal for coordinated timetable of electives in English
\item Advance publication of the academic offering for incoming students
\end{itemize}}\par
\fieldsep
\fieldlabel{KPIs}
{\small \begin{itemize}[noitemsep,topsep=2pt]
\item Coordinated timetable approved (Yes/No)
\item No. of mutually compatible subjects in English
\end{itemize}}\par
\end{tcolorbox}
\begin{tcolorbox}[riskbox]
{\footnotesize\sffamily\bfseries\color{red!50!black} RISKS}\par\nopagebreak[4]
{\small \begin{itemize}[noitemsep,topsep=2pt]
\item • Timetable rigidity due to general School restrictions.
\item • Coordination difficulty between programmes.
\end{itemize}}
\end{tcolorbox}
\begin{tcolorbox}[mitbox]
{\footnotesize\sffamily\bfseries\color{green!40!black} MITIGATION PLAN}\par\nopagebreak[4]
{\small \begin{itemize}[noitemsep,topsep=2pt]
\item • Prioritise minimum compatibilities in strategic subjects.
\item • Early coordination during the teaching planning phase.
\end{itemize}}
\end{tcolorbox}
\vspace{8pt}

\begin{tcolorbox}[pabox={fgcolor}{A.P.3.2 \enspace\textbullet\enspace Creation of FADE International Seminars}]

\fieldlabel{Description}
{\small Design a modular cross-cutting subject for incoming students, with sessions taught by faculty from other Schools on business, innovation, sustainability, culture and the international environment, strengthening the academic value and international distinctiveness of FADE's offering. }\par
\fieldsep
\fieldlabel{Deliverables}
{\small \begin{itemize}[noitemsep,topsep=2pt]
\item Academic design of the International Seminars subject
\item Annual programme and session calendar
\item Initial catalogue of participating faculty
\end{itemize}}\par
\fieldsep
\fieldlabel{KPIs}
{\small \begin{itemize}[noitemsep,topsep=2pt]
\item Subject designed and approved (Yes/No)
\item No. of incoming students enrolled
\item No. of participating faculty
\end{itemize}}\par
\end{tcolorbox}
\begin{tcolorbox}[riskbox]
{\footnotesize\sffamily\bfseries\color{red!50!black} RISKS}\par\nopagebreak[4]
{\small \begin{itemize}[noitemsep,topsep=2pt]
\item • Low initial enrolment in the subject.
\item • Difficulty in consolidating a stable calendar of participating faculty.
\item • Overlaps with other teaching activities.
\end{itemize}}
\end{tcolorbox}
\begin{tcolorbox}[mitbox]
{\footnotesize\sffamily\bfseries\color{green!40!black} MITIGATION PLAN}\par\nopagebreak[4]
{\small \begin{itemize}[noitemsep,topsep=2pt]
\item • Flexible initial implementation with a reduced number of seminars.
\item • Progressive recruitment of volunteer faculty.
\item • Advance coordination in teaching planning.
\end{itemize}}
\end{tcolorbox}
\vspace{8pt}

\newpage
\Needspace{14\baselineskip}
\addsec{S.O.4. Foster the international exchange of FADE students}
{\small\color{mutedcolor}\itshape Promote more balanced, accessible and sustainable outgoing mobility by adjusting academic destinations to actual student demand and encouraging additional financial incentives, particularly in international double degree programmes and excellence schemes such as PROMOE, thereby facilitating participation in experiences of high academic value and reducing the concentration of demand on particular destinations. }\par\vspace{6pt}

\begin{tcolorbox}[pabox={fgcolor}{A.P.4.1 \enspace\textbullet\enspace Adjust Supply and Demand between Destinations and Student Preferences}]

\fieldlabel{Description}
{\small Periodic analysis of the correspondence between the destination offering and student preferences to optimise the distribution of places and reinforce strategic agreements.}\par
\fieldsep
\fieldlabel{Deliverables}
{\small \begin{itemize}[noitemsep,topsep=2pt]
\item Annual supply-demand report on destinations
\item Proposal for adjustment or reorientation of agreements
\end{itemize}}\par
\fieldsep
\fieldlabel{KPIs}
{\small \begin{itemize}[noitemsep,topsep=2pt]
\item Report prepared annually (Yes/No)
\item \% occupancy of offered places
\end{itemize}}\par
\end{tcolorbox}
\begin{tcolorbox}[riskbox]
{\footnotesize\sffamily\bfseries\color{red!50!black} RISKS}\par\nopagebreak[4]
{\small \begin{itemize}[noitemsep,topsep=2pt]
\item • Persistence of demand concentration in few destinations.
\item • Institutional limitations for modifying agreements.
\end{itemize}}
\end{tcolorbox}
\begin{tcolorbox}[mitbox]
{\footnotesize\sffamily\bfseries\color{green!40!black} MITIGATION PLAN}\par\nopagebreak[4]
{\small \begin{itemize}[noitemsep,topsep=2pt]
\item • Reinforcement of information on alternative destinations.
\item • Progressive prioritisation of strategic destinations.
\end{itemize}}
\end{tcolorbox}
\vspace{8pt}

\begin{tcolorbox}[pabox={fgcolor}{A.P.4.2 \enspace\textbullet\enspace Mobility Incentives through Additional Fundraising}]

\fieldlabel{Description}
{\small Fundraising of complementary financing to reinforce mobility, especially in international double degrees and PROMOE programmes, reducing economic barriers.}\par
\fieldsep
\fieldlabel{Deliverables}
{\small \begin{itemize}[noitemsep,topsep=2pt]
\item Identification of complementary funding sources
\item Specific incentives proposal
\end{itemize}}\par
\fieldsep
\fieldlabel{KPIs}
{\small \begin{itemize}[noitemsep,topsep=2pt]
\item No. of applications submitted
\item Additional funds raised (€)
\item No. of students benefited
\end{itemize}}\par
\end{tcolorbox}
\begin{tcolorbox}[riskbox]
{\footnotesize\sffamily\bfseries\color{red!50!black} RISKS}\par\nopagebreak[4]
{\small \begin{itemize}[noitemsep,topsep=2pt]
\item • Failure to obtain external funding.
\item • Dependence on competitive calls.
\end{itemize}}
\end{tcolorbox}
\begin{tcolorbox}[mitbox]
{\footnotesize\sffamily\bfseries\color{green!40!black} MITIGATION PLAN}\par\nopagebreak[4]
{\small \begin{itemize}[noitemsep,topsep=2pt]
\item • Diversify funding sources.
\item • Prioritise calls compatible with international mobility.
\end{itemize}}
\end{tcolorbox}
\vspace{8pt}

\newpage
\Needspace{14\baselineskip}
\addsec{S.O.5. Foster the international reputation of the School}
{\small\color{mutedcolor}\itshape Reinforce the international positioning of FADE through strategic academic alliances with prestigious institutions and the promotion of faculty and staff mobility.}\par\vspace{6pt}

\begin{tcolorbox}[pabox={fgcolor}{A.P.5.1 \enspace\textbullet\enspace Promote International Double Degree Agreements with Business Schools of Recognised Standing}]

\fieldlabel{Description}
{\small Promote international double degree agreements with business schools of recognised standing, prioritising institutions with academic affinity and strategic complementarity.}\par
\fieldsep
\fieldlabel{Deliverables}
{\small \begin{itemize}[noitemsep,topsep=2pt]
\item Identification of target universities
\item Institutional negotiation roadmap
\item Formal academic agreement proposals
\end{itemize}}\par
\fieldsep
\fieldlabel{KPIs}
{\small \begin{itemize}[noitemsep,topsep=2pt]
\item No. of institutions analysed
\item No. of negotiations started
\item No. of agreements formalised or in advanced phase
\item No. of active international double degrees
\end{itemize}}\par
\end{tcolorbox}
\begin{tcolorbox}[riskbox]
{\footnotesize\sffamily\bfseries\color{red!50!black} RISKS}\par\nopagebreak[4]
{\small \begin{itemize}[noitemsep,topsep=2pt]
\item • Long negotiation processes.
\item • Complex curricular differences.
\item • High mutual demands.
\end{itemize}}
\end{tcolorbox}
\begin{tcolorbox}[mitbox]
{\footnotesize\sffamily\bfseries\color{green!40!black} MITIGATION PLAN}\par\nopagebreak[4]
{\small \begin{itemize}[noitemsep,topsep=2pt]
\item • Prioritisation of compatible strategic partners.
\item • Start with academically close agreements.
\end{itemize}}
\end{tcolorbox}
\vspace{8pt}

\begin{tcolorbox}[pabox={fgcolor}{A.P.5.2 \enspace\textbullet\enspace Foster Professional Services Staff (PTGAS) and Faculty (PDI) Mobility}]

\fieldlabel{Description}
{\small Drive the participation of professional services staff (PTGAS) and faculty (PDI) in international mobility programmes, academic visits and institutional cooperation, strengthening academic networks and good practices.}\par
\fieldsep
\fieldlabel{Deliverables}
{\small \begin{itemize}[noitemsep,topsep=2pt]
\item Specific dissemination of mobility calls
\item Annual record of mobility periods carried out
\end{itemize}}\par
\fieldsep
\fieldlabel{KPIs}
{\small \begin{itemize}[noitemsep,topsep=2pt]
\item No. of professional services staff (PTGAS) mobility periods carried out
\item No. of faculty (PDI) mobility periods carried out
\end{itemize}}\par
\end{tcolorbox}
\begin{tcolorbox}[riskbox]
{\footnotesize\sffamily\bfseries\color{red!50!black} RISKS}\par\nopagebreak[4]
{\small \begin{itemize}[noitemsep,topsep=2pt]
\item • Low participation due to teaching or administrative workload.
\item • Limitation of available funding.
\end{itemize}}
\end{tcolorbox}
\begin{tcolorbox}[mitbox]
{\footnotesize\sffamily\bfseries\color{green!40!black} MITIGATION PLAN}\par\nopagebreak[4]
{\small \begin{itemize}[noitemsep,topsep=2pt]
\item • Advance dissemination.
\item • Prioritisation of strategic mobilities.
\item • Utilisation of existing active calls.
\end{itemize}}
\end{tcolorbox}
\vspace{8pt}

\Needspace{18\baselineskip}
\section{Timeline}

\small

\begin{tabularx}{\textwidth}{>{\bfseries}l >{\raggedright\arraybackslash}X *{4}{>{\centering\arraybackslash}p{1.5cm}}}
\toprule
A.P. & Action plan & 2026/27 & 2027/28 & 2028/29 & 2029/30 \\
\midrule
A.P.1.1 & Interactive destination map & $\blacksquare$ &  &  &  \\
A.P.1.2 & AI chatbot and recognition &  & $\blacksquare$ & $\blacksquare$ &  \\
A.P.2.1 & Policonsulta implementation & $\blacksquare$ &  &  &  \\
A.P.2.2 & Updating of recognition regulations & $\blacksquare$ &  &  &  \\
A.P.3.1 & Coordinated timetable for electives in English & $\blacksquare$ & $\blacksquare$ &  &  \\
A.P.3.2 & FADE International Seminars &  & $\blacksquare$ & $\blacksquare$ &  \\
A.P.4.1 & Adjustment of destination supply/demand &  & $\blacksquare$ & $\blacksquare$ & $\blacksquare$ \\
A.P.4.2 & Mobility incentives / PROMOE funds &  & $\blacksquare$ & $\blacksquare$ & $\blacksquare$ \\
A.P.5.1 & International double degree agreements &  &  & $\blacksquare$ & $\blacksquare$ \\
A.P.5.2 & Professional services staff and faculty mobility &  & $\blacksquare$ & $\blacksquare$ & $\blacksquare$ \\
\bottomrule
\end{tabularx}

\Needspace{18\baselineskip}
\section{Summary of Objective Achievement}

\small
Achievement of the project's objectives will be evaluated through the monitoring of indicators defined in each action plan, as well as an overall assessment of the impact on students' international mobility and the School's international positioning.

\textbf{S.O.1. Improve access to, understanding and usability of information on international academic exchange.} It is considered achieved when students have an updated interactive map and an operational AI chatbot that enable clear, structured access to international mobility opportunities, understanding of academic recognition conditions and a reduction in repetitive enquiries, ensuring a transparent and manageable mobility experience.

\textbf{S.O.2. Improve support and accompaniment for outgoing students.} It is considered achieved when all outgoing students use Policonsulta as their main communication channel, and academic recognition regulations are consolidated in a single accessible digital document, enabling consistent monitoring, fewer incidents and high-quality academic support.

\textbf{S.O.3. Facilitate the academic integration of incoming international students.} It is considered achieved when international students can enrol in academic pathways without timetable conflicts thanks to the coordination of subjects taught in English, and actively participate in FADE International Seminars, receiving comprehensive academic and cultural integration that enhances the perceived quality of the international offering.

\textbf{S.O.4. Foster the international exchange of FADE students.} It is considered achieved when destination opportunities match student preferences, balancing the allocation of places, and financial incentives are implemented through the acquisition of additional funds (including PROMOE programmes and international double degrees), facilitating participation in outgoing mobility of high academic value and reducing financial barriers.

\textbf{S.O.5. Foster the international reputation of the School.} It is considered achieved when FADE has formalised strategic double degree agreements with business schools of recognised standing and has promoted professional services staff and faculty mobility, generating international visibility, strengthening academic networks and consolidating the School's reputation as a global benchmark in education in organisational management and leadership.
\clearpage

\chapter[FADE Impact Lab]{\textcolor{filcolor}{FADE Impact Lab}}
\begin{tcolorbox}[descbox]
\small Structures and empowers the School's relationship ecosystem, connecting its academic activity with the professional and social environment through co-created challenges or missions. It integrates employability, entrepreneurship, internship and collaboration initiatives with companies, institutions and alumni, facilitating their coordination and monitoring. It uses data and artificial intelligence to identify opportunities and improve relational management. Its purpose is to convert FADE's knowledge into real projects with impact, enriching the educational experience of students.
\end{tcolorbox}

\Needspace{18\baselineskip}
\section{Project Description}

\small

\small
FADE Impact Lab is a strategic project built on a structural competitive advantage: FADE's membership of the Universitat Politècnica de València. In an environment where the speed of technological change is redefining the model of traditional business schools, having a leading technological university such as UPV represents a barrier to entry that private competitors would find difficult to replicate. This symbiosis between management and leadership education and frontier technological and research capacity is the central pillar of FADE's value proposition to the business community and society, and the foundation on which this project is built.

On this basis, FADE Impact Lab aims to structure, coordinate and strengthen the School's ecosystem of relationships, activities and initiatives, coherently connecting its academic activity with its professional and social environment. Its objective is to transform activity that is currently valuable but dispersed into an organised system that enables more effective activation, monitoring and communication of FADE's interactions with businesses, institutions, alumni and other actors, organising them around shared challenges or missions (FADE Challenges) that guide collaboration and generate impact.

The project integrates within a common framework the full range of FADE activities and initiatives related to teaching, applied research, employability, entrepreneurship, internships, Business Chair and master's programme activities, participatory activities and engagement with alumni and other external actors. Its purpose is to organise and activate these connections, facilitating their coordination, monitoring and recognition, and strengthening collaboration with businesses, professional bodies, associations, public administrations and other organisations. This approach enables different forms of interaction (internships, Bachelor's Degree Final Projects/Master's Theses in organisations, joint projects, activities with organisations or knowledge transfer initiatives) to be coherently organised within the FADE ecosystem, avoiding fragmentation and improving their academic, professional and social impact.

The project relies on four cross-cutting action mechanisms: Processes, to systematise and coordinate activities, relationships and initiatives in the FADE ecosystem; Data, to record and monitor activities, participating actors and results; Communication, to make opportunities, connections and achievements generated by the School more visible and accessible; and Artificial Intelligence (AI), as a tool supporting analysis and management of the relationship ecosystem. While Processes coordinates and Data records, AI facilitates the identification of potential connections between challenges, activities and actors, supporting tasks such as relationship analysis, information organisation or the generation of collaboration proposals. Its incorporation is envisaged as progressive and focused on specific use cases, with the aim of improving efficiency in ecosystem management without increasing the administrative workload. It will never replace academic or institutional judgement but will act as a complementary tool supporting decision-making.

\Needspace{18\baselineskip}
\section{Relationship with the Strategic Lines}

\small
The FADE Impact Lab project's strategic contribution is distributed as follows:

\begin{itemize}[noitemsep,topsep=2pt]
\item \textbf{Academic Excellence → 35\%}
\item \textbf{Impact → 65\%}
\end{itemize}

The main weighting towards Impact is justified because the project aims to structure, activate, monitor and make visible FADE's ecosystem of relationships with businesses, professional bodies, associations, public administrations, partner organisations and alumni, as well as to integrate and communicate activities connecting the School with its social and professional environment.

At the same time, the project contributes to Academic Excellence by improving students' connections with employability, entrepreneurship and applied learning opportunities, fostering a university experience more closely connected to professional reality, and providing useful information to organise, monitor and improve initiatives with educational value.

The 35\% / 65\% distribution reflects the project's nature: an initiative clearly oriented towards external impact, but with direct effects on the quality and added value of the academic experience.



\small
The project contributes to Academic Excellence through:

\begin{itemize}[noitemsep,topsep=2pt]
\item Improved connections between students and employability, entrepreneurship and internship opportunities.
\item Strengthening applied learning by incorporating activities, challenges and interactions with the external environment into subjects, Bachelor's Degree Final Projects/Master's Theses and other educational experiences.
\item Integration of complementary activities of high educational value (talks, lectures, participatory activities).
\item Development of a more active university experience connected to professional reality.
\item Use of structured information to guide academic decision-making.
\end{itemize}



\small
The contribution to the Impact line is reflected in:

\begin{itemize}[noitemsep,topsep=2pt]
\item Structuring and coherence of FADE's ecosystem of relationships with its environment.
\item Integration and visibility of activities and interactions with social, professional and institutional value.
\item Strengthening knowledge transfer and external engagement through internships, Bachelor's Degree Final Projects/Master's Theses, joint projects and external collaborations.
\item Improved traceability of impact through monitoring, data and communication systems.
\item Increased external visibility and institutional capacity to manage and demonstrate impact.
\end{itemize}
\newpage
\Needspace{18\baselineskip}
\section{Objectives, Action Plans and Development}
\Needspace{14\baselineskip}
\addsec{S.O.1. Structure and integrate the FADE ecosystem through co-created challenges or missions}
{\small\color{mutedcolor}\itshape Develop a common framework enabling FADE to identify, organise and coordinate initiatives, activities and connections related to employability, entrepreneurship, internships, applied teaching, Business Chair and master's programme activities, actions promoted by services and teams, participatory activities and stakeholder collaboration, paying particular attention to areas where FADE's academic capacity matches real external needs, such as sustainability, technological innovation or corporate governance. This objective places FADE challenges or co-created missions at the centre of external engagement and collaboration. FADE Challenges are real missions or problems, defined collaboratively with external partners, that bring together the School's academic activity, educational opportunities and relationships, connecting FADE's knowledge with specific professional and social needs. They are therefore the means of turning our knowledge into real projects with impact. }\par\vspace{6pt}

\begin{tcolorbox}[pabox={filcolor}{A.P.1.1 \enspace\textbullet\enspace Mapping, Inventory and Classification of the FADE Ecosystem}]

\fieldlabel{Description}
{\small Identification, collection, organisation and classification of the set of activities, agents and relationships existing in FADE linked to employability, entrepreneurship, internships, applied teaching, alumni and stakeholder collaboration.}\par
\fieldsep
\fieldlabel{Deliverables}
{\small \begin{itemize}[noitemsep,topsep=2pt]
\item FADE ecosystem framework document
\item Structured inventory of activities, agents and relationships
\item Basic taxonomy for ecosystem classification
\item Basic information structure for management and monitoring
\end{itemize}}\par
\fieldsep
\fieldlabel{KPIs}
{\small \begin{itemize}[noitemsep,topsep=2pt]
\item Framework document prepared and validated (Yes/No)
\item \% FADE areas/units incorporated into the inventory
\item No. of activities and relationships identified and classified
\item Basic information structure defined (Yes/No)
\end{itemize}}\par
\end{tcolorbox}
\begin{tcolorbox}[riskbox]
{\footnotesize\sffamily\bfseries\color{red!50!black} RISKS}\par\nopagebreak[4]
{\small \begin{itemize}[noitemsep,topsep=2pt]
\item • Incomplete identification of existing activities or relationships.
\item • Lack of homogeneity in information provided by different units.
\item • Generation of an inventory of little use for subsequent management.
\end{itemize}}
\end{tcolorbox}
\begin{tcolorbox}[mitbox]
{\footnotesize\sffamily\bfseries\color{green!40!black} MITIGATION PLAN}\par\nopagebreak[4]
{\small \begin{itemize}[noitemsep,topsep=2pt]
\item • Use a common information collection template and a shared basic classification.
\item • Carry out a brief validation with managers of the main units involved.
\item • Design the inventory with a practical approach, oriented towards subsequent coordination and monitoring.
\end{itemize}}
\end{tcolorbox}
\vspace{8pt}

\begin{tcolorbox}[pabox={filcolor}{A.P.1.2 \enspace\textbullet\enspace Definition of the Annual Co-creation and Prioritisation Process for FADE Challenges}]

\fieldlabel{Description}
{\small Definition of an annual process for identifying, co-creating, prioritising and validating the FADE Challenges that will guide activity in each academic year. This action plan will establish a shared process, undertaken at the end of each academic year or the beginning of the next, to specify the challenges or missions that will form the backbone of activities, opportunities and connections in the FADE ecosystem related to employability, entrepreneurship, internships, applied teaching, stakeholder collaboration and other School initiatives. FADE Challenges will be defined consistently with the areas in which the School has the greatest academic capacity and specialist knowledge, drawing on Dashboard information to identify areas where FADE can provide greater value to its environment. Thus, the challenges not only respond to external needs but are grounded in the School's own knowledge and experience, strengthening its positioning and thought leadership in those areas. }\par
\fieldsep
\fieldlabel{Deliverables}
{\small \begin{itemize}[noitemsep,topsep=2pt]
\item Annual co-creation process document
\item Challenge formulation and selection criteria
\item Process calendar
\item Initial catalogue of prioritised FADE Challenges
\end{itemize}}\par
\fieldsep
\fieldlabel{KPIs}
{\small \begin{itemize}[noitemsep,topsep=2pt]
\item Annual process defined and validated (Yes/No)
\item Formulation and selection criteria defined (Yes/No)
\item No. of FADE Challenges prioritised for the year
\item Initial catalogue prepared (Yes/No)
\end{itemize}}\par
\end{tcolorbox}
\begin{tcolorbox}[riskbox]
{\footnotesize\sffamily\bfseries\color{red!50!black} RISKS}\par\nopagebreak[4]
{\small \begin{itemize}[noitemsep,topsep=2pt]
\item • Defining challenges that are too generic or of little use for guiding the year's activity.
\item • Lack of alignment between the prioritised challenges and FADE's real capabilities or initiatives.
\item • Scarce involvement of key units in the co-creation and prioritisation process.
\end{itemize}}
\end{tcolorbox}
\begin{tcolorbox}[mitbox]
{\footnotesize\sffamily\bfseries\color{green!40!black} MITIGATION PLAN}\par\nopagebreak[4]
{\small \begin{itemize}[noitemsep,topsep=2pt]
\item • Establish clear formulation and selection criteria, oriented towards relevance, feasibility and mobilisation capacity.
\item • Link the definition of challenges to the prior map of activities, agents and opportunities of the FADE ecosystem.
\item • Carry out a final validation with academic managers and key units before approving each year's challenges.
\end{itemize}}
\end{tcolorbox}
\vspace{8pt}

\begin{tcolorbox}[pabox={filcolor}{A.P.1.3 \enspace\textbullet\enspace Definition of the Governance and Coordination Model of the FADE Ecosystem}]

\fieldlabel{Description}
{\small Definition of a governance and coordination model for the FADE ecosystem that enables the organisation of responsibilities, information flows, internal interaction mechanisms and common criteria for bringing together activities, opportunities and external relationships. This action plan will establish a shared organisational basis for ecosystem management, including coordination between units, validation of FADE Challenges and definition of the functional support needed to record and manage project-related information. }\par
\fieldsep
\fieldlabel{Deliverables}
{\small \begin{itemize}[noitemsep,topsep=2pt]
\item Ecosystem governance and coordination model
\item Roles, responsibilities and coordination mechanisms
\item Functional proposal for a support tool
\item Basic information management protocol
\end{itemize}}\par
\fieldsep
\fieldlabel{KPIs}
{\small \begin{itemize}[noitemsep,topsep=2pt]
\item Governance model validated (Yes/No)
\item Roles and mechanisms formalised (Yes/No)
\item Functional tool proposal defined (Yes/No)
\item Basic protocol approved (Yes/No)
\end{itemize}}\par
\end{tcolorbox}
\begin{tcolorbox}[riskbox]
{\footnotesize\sffamily\bfseries\color{red!50!black} RISKS}\par\nopagebreak[4]
{\small \begin{itemize}[noitemsep,topsep=2pt]
\item • Definition of a governance model that is unclear or difficult to sustain in practice.
\item • Excessive organisational complexity or coordination burden between units.
\item • Mismatch between ecosystem needs and the functional solution planned for its management.
\end{itemize}}
\end{tcolorbox}
\begin{tcolorbox}[mitbox]
{\footnotesize\sffamily\bfseries\color{green!40!black} MITIGATION PLAN}\par\nopagebreak[4]
{\small \begin{itemize}[noitemsep,topsep=2pt]
\item • Design a simple, realistic model supported by already existing structures in FADE wherever possible.
\item • Limit the number of roles and coordination circuits to those strictly necessary.
\item • Define the functional support solution based on actual needs previously identified in A.P.1.1 and aligned with the planned uses in the project.
\end{itemize}}
\end{tcolorbox}
\vspace{8pt}

\newpage
\Needspace{14\baselineskip}
\addsec{S.O.2. Improve student connection with employability, entrepreneurship, internship and active participation opportunities}
{\small\color{mutedcolor}\itshape Provide students with more structured, visible and coherent access to the full range of activities, experiences and opportunities offered by FADE in career guidance, employability, entrepreneurship, internships, participation in challenges, collaboration with organisations and other initiatives of high educational value and value for personal and professional development.  }\par\vspace{6pt}

\begin{tcolorbox}[pabox={filcolor}{A.P.2.1 \enspace\textbullet\enspace Integrated Organisation and Visibility of the Opportunities Offering for Students}]

\fieldlabel{Description}
{\small Organisation, structuring and integrated visibility of the full range of activities, experiences and opportunities offered by FADE to students in employability, entrepreneurship, internships, participation in FADE Challenges, collaboration with organisations and other initiatives of high educational value and value for personal and professional development. This action plan will enable the offering to be presented more clearly, accessibly and coherently, helping students become more aware of available opportunities and more easily identify participation pathways aligned with their interests. }\par
\fieldsep
\fieldlabel{Deliverables}
{\small \begin{itemize}[noitemsep,topsep=2pt]
\item Offering organisation framework document
\item Structured opportunities catalogue
\item Classification and presentation criteria
\item Basic prototype of integrated channel
\end{itemize}}\par
\fieldsep
\fieldlabel{KPIs}
{\small \begin{itemize}[noitemsep,topsep=2pt]
\item Framework document validated (Yes/No)
\item Structured catalogue prepared (Yes/No)
\item No. of opportunity typologies classified
\item Basic prototype of integrated channel defined (Yes/No)
\end{itemize}}\par
\end{tcolorbox}
\begin{tcolorbox}[riskbox]
{\footnotesize\sffamily\bfseries\color{red!50!black} RISKS}\par\nopagebreak[4]
{\small \begin{itemize}[noitemsep,topsep=2pt]
\item • Persistence of a fragmented or difficult-to-homogenise offering.
\item • Difficulty in presenting information clearly and usefully for diverse student profiles.
\item • Lack of updating or continuity in the organisation of the offering.
\end{itemize}}
\end{tcolorbox}
\begin{tcolorbox}[mitbox]
{\footnotesize\sffamily\bfseries\color{green!40!black} MITIGATION PLAN}\par\nopagebreak[4]
{\small \begin{itemize}[noitemsep,topsep=2pt]
\item • Establish a common classification and homogeneous criteria for presenting opportunities.
\item • Validate the organisation and visibility proposal with a small sample of students and involved units.
\item • Design the solution with a simple, practical and sustainable approach over time.
\end{itemize}}
\end{tcolorbox}
\vspace{8pt}

\begin{tcolorbox}[pabox={filcolor}{A.P.2.2 \enspace\textbullet\enspace Promotion of Student Participation and Connection with the FADE Ecosystem and its Environment}]

\fieldlabel{Description}
{\small Design and promotion of actions that encourage greater student participation in activities, experiences and opportunities linked to the FADE ecosystem, while strengthening their connections with businesses, institutions, alumni and other stakeholders. This action plan will promote more active student involvement in employability, entrepreneurship and internship activities, participation in FADE Challenges, participatory activities and collaboration with organisations, contributing to an educational experience more closely connected to the professional and social environment. }\par
\fieldsep
\fieldlabel{Deliverables}
{\small \begin{itemize}[noitemsep,topsep=2pt]
\item Approach document for student activation
\item Proposal of activation actions/formats
\item Pilot programme of activities
\item Basic participation monitoring criteria
\end{itemize}}\par
\fieldsep
\fieldlabel{KPIs}
{\small \begin{itemize}[noitemsep,topsep=2pt]
\item Approach document validated (Yes/No)
\item No. of activation actions/formats defined
\item Pilot programme defined (Yes/No)
\item Monitoring criteria established (Yes/No)
\end{itemize}}\par
\end{tcolorbox}
\begin{tcolorbox}[riskbox]
{\footnotesize\sffamily\bfseries\color{red!50!black} RISKS}\par\nopagebreak[4]
{\small \begin{itemize}[noitemsep,topsep=2pt]
\item • Low initial student participation in proposed activities.
\item • Scarce connection between the promoted actions and students' real interests.
\item • Difficulty in sustaining participation and involvement of entities or alumni over time.
\end{itemize}}
\end{tcolorbox}
\begin{tcolorbox}[mitbox]
{\footnotesize\sffamily\bfseries\color{green!40!black} MITIGATION PLAN}\par\nopagebreak[4]
{\small \begin{itemize}[noitemsep,topsep=2pt]
\item • Design participation actions using attractive, useful formats aligned with real student needs.
\item • Previously validate the proposal with a small sample of students and agents involved.
\item • Prioritise a limited number of viable actions with capacity for continuity, in collaboration with units and actors already active in FADE.
\end{itemize}}
\end{tcolorbox}
\vspace{8pt}

\newpage
\Needspace{14\baselineskip}
\addsec{S.O.3. Systematise the monitoring of FADE's interactions and collaborations with its stakeholders}
{\small\color{mutedcolor}\itshape Implement a monitoring system enabling the consistent recording, classification and analysis of the School's relationships and interactions with businesses, professional bodies, associations, public administrations, partner organisations and alumni, including dissemination activities, business internships, Bachelor's Degree Final Projects (TFG) and Master's Theses (TFM) linked to organisations, joint projects, patronage and other forms of institutional, professional or academic collaboration.  }\par\vspace{6pt}

\begin{tcolorbox}[pabox={filcolor}{A.P.3.1 \enspace\textbullet\enspace Implementation of the Interaction and Collaboration Recording and Monitoring System}]

\fieldlabel{Description}
{\small Implementation of a consistent system for recording, updating and monitoring FADE's stakeholder interactions and collaborations, in accordance with the previously defined information structure and governance model. This action plan will enable the orderly recording of activities, actors, relationship types and results associated with interaction with businesses, professional bodies, associations, public administrations, partner organisations and alumni, generating a common operational basis for institutional monitoring of the School's relationship ecosystem. }\par
\fieldsep
\fieldlabel{Deliverables}
{\small \begin{itemize}[noitemsep,topsep=2pt]
\item Basic recording and monitoring system implemented
\item Operational protocol for use, updating and validation
\item Initial set of interactions recorded
\item Basic usage guide
\end{itemize}}\par
\fieldsep
\fieldlabel{KPIs}
{\small \begin{itemize}[noitemsep,topsep=2pt]
\item Basic system implemented (Yes/No)
\item Operational protocol approved (Yes/No)
\item No. of interactions recorded in initial phase
\item No. of units incorporated into system use
\end{itemize}}\par
\end{tcolorbox}
\begin{tcolorbox}[riskbox]
{\footnotesize\sffamily\bfseries\color{red!50!black} RISKS}\par\nopagebreak[4]
{\small \begin{itemize}[noitemsep,topsep=2pt]
\item • Low system adoption by the units or people involved.
\item • Incomplete or irregular recording of information.
\item • Excessive complexity in the tool or in the usage protocol.
\end{itemize}}
\end{tcolorbox}
\begin{tcolorbox}[mitbox]
{\footnotesize\sffamily\bfseries\color{green!40!black} MITIGATION PLAN}\par\nopagebreak[4]
{\small \begin{itemize}[noitemsep,topsep=2pt]
\item • Implement a simple solution, adjusted to actual needs and easy to use.
\item • Define minimum responsibilities and a clear operational protocol for recording and updating information.
\item • Carry out an initial piloting or gradual deployment phase with a limited number of units before progressive extension.
\end{itemize}}
\end{tcolorbox}
\vspace{8pt}

\begin{tcolorbox}[pabox={filcolor}{A.P.3.2 \enspace\textbullet\enspace Analysis and Use of Information to Improve FADE's Relational Management}]

\fieldlabel{Description}
{\small Development of basic mechanisms for analysing, querying and using the information recorded on FADE's stakeholder interactions and collaborations, to improve institutional understanding of its relationship ecosystem and support decision-making. This action plan will identify relationship patterns, areas of activity, types of collaboration, the most relevant actors and opportunities for improvement, generating useful information for the project's planning, coordination and strategic monitoring. }\par
\fieldsep
\fieldlabel{Deliverables}
{\small \begin{itemize}[noitemsep,topsep=2pt]
\item Basic analysis criteria and indicators
\item Basic information analysis report/dashboard
\item Periodic reporting scheme
\item Initial improvement recommendations
\end{itemize}}\par
\fieldsep
\fieldlabel{KPIs}
{\small \begin{itemize}[noitemsep,topsep=2pt]
\item Basic criteria and indicators defined (Yes/No)
\item Basic information analysis report prepared (Yes/No)
\item Reporting scheme defined (Yes/No)
\item No. of improvement recommendations formulated
\end{itemize}}\par
\end{tcolorbox}
\begin{tcolorbox}[riskbox]
{\footnotesize\sffamily\bfseries\color{red!50!black} RISKS}\par\nopagebreak[4]
{\small \begin{itemize}[noitemsep,topsep=2pt]
\item • Insufficient availability or limited quality of recorded information for carrying out useful analyses.
\item • Generation of reports poorly oriented towards decision-making.
\item • Excessive complexity in the indicators or in the analysis and use of information.
\end{itemize}}
\end{tcolorbox}
\begin{tcolorbox}[mitbox]
{\footnotesize\sffamily\bfseries\color{green!40!black} MITIGATION PLAN}\par\nopagebreak[4]
{\small \begin{itemize}[noitemsep,topsep=2pt]
\item • Prioritise a reduced set of relevant and viable indicators from the information effectively available.
\item • Orient the analysis towards concrete and useful management questions for FADE.
\item • Develop a basic progressive reporting scheme, expandable as the monitoring system matures.
\end{itemize}}
\end{tcolorbox}
\vspace{8pt}

\newpage
\Needspace{14\baselineskip}
\addsec{S.O.4. Reinforce the visibility, communication and measurement of the impact generated by FADE}
{\small\color{mutedcolor}\itshape Identify, structure and highlight the activities, connections and results of the FADE ecosystem, generating evidence of impact enabling the School to demonstrate its contribution to society, organisations and the responsible development of its environment. Dissemination of this evidence is coordinated with the COMUNICA project, which acts as the institutional channel for external communication. }\par\vspace{6pt}

\begin{tcolorbox}[pabox={filcolor}{A.P.4.1 \enspace\textbullet\enspace Structuring, Visibility and Highlighting the Value of the Impact of the FADE Ecosystem}]

\fieldlabel{Description}
{\small Identification, structuring and recognition of the FADE ecosystem's content, results and evidence that need to be made visible and understandable to students, faculty, alumni and partner organisations, consistently with the procedures, channels and criteria defined in the COMUNICA Project. This action plan will clearly organise which project activities, opportunities, connections and results should be communicated, with what approach and to which audiences, strengthening FADE's capacity to make visible and communicate the academic, professional and social impact generated by its relationship ecosystem. Public dissemination of this evidence will be coordinated with the COMUNICA Project, avoiding duplication and ensuring a coherent institutional narrative. 
 }\par
\fieldsep
\fieldlabel{Deliverables}
{\small \begin{itemize}[noitemsep,topsep=2pt]
\item Content identification and structuring framework document
\item Selection criteria and criteria for highlighting value
\item Prioritised content, evidence and results
\item Visibility scheme aligned with COMUNICA
\end{itemize}}\par
\fieldsep
\fieldlabel{KPIs}
{\small \begin{itemize}[noitemsep,topsep=2pt]
\item Framework document validated (Yes/No)
\item Selection criteria and criteria for highlighting value defined (Yes/No)
\item No. of content items, evidence or results prioritised
\item Visibility scheme aligned with COMUNICA (Yes/No)
\end{itemize}}\par
\end{tcolorbox}
\begin{tcolorbox}[riskbox]
{\footnotesize\sffamily\bfseries\color{red!50!black} RISKS}\par\nopagebreak[4]
{\small \begin{itemize}[noitemsep,topsep=2pt]
\item • Overlap or lack of coordination with the COMUNICA Project.
\item • Unclear or insufficiently strategic selection of content and impacts to make visible.
\item • Difficulty in translating project information into understandable and relevant messages for different audiences.
\end{itemize}}
\end{tcolorbox}
\begin{tcolorbox}[mitbox]
{\footnotesize\sffamily\bfseries\color{green!40!black} MITIGATION PLAN}\par\nopagebreak[4]
{\small \begin{itemize}[noitemsep,topsep=2pt]
\item • Align this action plan from the start with the criteria, procedures and channels defined in the COMUNICA Project.
\item • Establish clear selection and prioritisation criteria for content, oriented towards relevance, utility and institutional value.
\item • Prioritise a simple, comprehensible visibility proposal adapted to specific audiences.
\end{itemize}}
\end{tcolorbox}
\vspace{8pt}

\Needspace{18\baselineskip}
\section{Timeline}

\small

\begin{tabularx}{\textwidth}{>{\bfseries}l *{8}{>{\centering\arraybackslash}p{1.3cm}}}
\toprule
A.P. & 26/27 S1 & 26/27 S2 & 27/28 S1 & 27/28 S2 & 28/29 S1 & 28/29 S2 & 29/30 S1 & 29/30 S2 \\
\midrule
A.P.1.1 & $\blacksquare$ & $\blacksquare$ &  &  &  &  &  &  \\
A.P.1.2 &  & $\blacksquare$ & $\blacksquare$ & $\blacksquare$ & $\blacksquare$ & $\blacksquare$ & $\blacksquare$ & $\blacksquare$ \\
A.P.1.3 &  & $\blacksquare$ & $\blacksquare$ &  &  &  &  &  \\
A.P.2.1 &  & $\blacksquare$ & $\blacksquare$ & $\blacksquare$ &  &  &  &  \\
A.P.2.2 &  &  & $\blacksquare$ & $\blacksquare$ & $\blacksquare$ & $\blacksquare$ & $\blacksquare$ & $\blacksquare$ \\
A.P.3.1 &  &  & $\blacksquare$ & $\blacksquare$ & $\blacksquare$ &  &  &  \\
A.P.3.2 &  &  &  & $\blacksquare$ & $\blacksquare$ & $\blacksquare$ & $\blacksquare$ & $\blacksquare$ \\
A.P.4.1 &  &  &  & $\blacksquare$ & $\blacksquare$ & $\blacksquare$ & $\blacksquare$ & $\blacksquare$ \\
\bottomrule
\end{tabularx}



\Needspace{18\baselineskip}
\section{Summary of Objective Achievement}

\small
Achievement of the project's objectives will be evaluated through the monitoring of indicators defined in each action plan, as well as an overall assessment of the degree of structuring, activation and visibility of the School's ecosystem of relationships, activities and initiatives.

\textbf{S.O.1. Structure and integrate the FADE ecosystem of activities, opportunities and external relationships through co-created challenges or missions.} It will be considered achieved when FADE has a common framework for coherently identifying, organising and coordinating the activities, initiatives, actors and relationships forming its ecosystem, as well as a defined annual process for co-creating and prioritising FADE Challenges and a basic governance model enabling sustained coordination and development over time.

\textbf{S.O.2. Improve student connection with employability, entrepreneurship, internship and active participation opportunities.} It will be considered achieved when students have a more structured, visible and accessible range of opportunities, and FADE has introduced mechanisms and actions that effectively encourage participation in activities, challenges, internships, entrepreneurship initiatives and other experiences connected with organisations, alumni and the professional and social environment.

\textbf{S.O.3. Systematise the monitoring of FADE's interactions and collaborations with its stakeholders.} It will be considered achieved when the School has an operational system for recording and monitoring its stakeholder interactions and collaborations, used consistently by the units involved, and when the information generated enables basic analysis of FADE ecosystem relationships and supports decision-making to improve relationship management.

\textbf{S.O.4. Reinforce the visibility, communication and measurement of the impact generated by FADE.} It will be considered achieved when FADE has a structured proposal for ecosystem content, results and evidence that makes its impact more visible, understandable and valuable to students, faculty, alumni and partner organisations, consistently with the criteria, procedures and channels defined in the COMUNICA Project.


\clearpage

\chapter[COMUNICA]{\textcolor{comcolor}{COMUNICA}}
\begin{tcolorbox}[descbox]
\small Strengthens the internal and external communication system of FADE through the definition of common processes, flows, criteria and channels to provide coherence and traceability to all its activity. It organises the collection, validation and publication of content generated by programmes, chairs, services and other agents. Supported by data and artificial intelligence, it improves the visibility of teaching, research and outreach. It reinforces institutional visibility and systematically disseminates the intellectual leadership of the FADE community.
\end{tcolorbox}
\Needspace{18\baselineskip}
\section{Project Description}

\small
COMUNICA is a strategic project aimed at structuring and strengthening FADE's internal and external communication system, based on a clear premise: institutional communication is not merely a dissemination function but a strategic tool for building reputation, generating trust and consolidating a distinctive positioning among the School's various stakeholder groups. With this orientation, the project develops a common framework for collecting, validating, prioritising, scheduling and publishing content generated by the various internal and external actors linked to FADE, including bachelor's and master's programmes, Business Chairs, projects, services, teams, faculty, students, alumni and partner organisations, with the aim of creating a coherent, recognisable institutional narrative that conveys the FADE community's thought leadership and impact.

The project is structured around four cross-cutting action mechanisms: Processes, Data, Communication and Artificial Intelligence. First, the Processes mechanism defines publication procedures, workflows and criteria, ensuring an orderly, efficient working model. Second, the Data mechanism enables FADE's communication activity to be recorded, archived and made traceable, building a robust institutional memory. Third, the Communication mechanism adapts content and channels to different audiences, strengthening the School's institutional visibility. Artificial Intelligence (AI) joins this ecosystem as a fourth cross-cutting action mechanism, acting as the project's operational optimisation layer. While Processes structures and Data archives, Artificial Intelligence (AI) is incorporated as a tool supporting content management and processing, facilitating tasks such as information classification, adaptation of formats to different channels or generation of summaries. Its use is envisaged as progressive and focused on specific cases, with the aim of improving operational efficiency without replacing the responsible team's editorial judgement, contributing to more agile, structured dissemination of FADE's thought leadership.

\Needspace{18\baselineskip}
\section{Relationship with the Strategic Lines}

\small
The COMUNICA project's strategic contribution is distributed as follows:

\begin{itemize}[noitemsep,topsep=2pt]
\item \textbf{Academic Excellence → 35\%}
\item \textbf{Impact → 65\%}
\end{itemize}

The project contributes predominantly to Impact by strengthening FADE's external and institutional visibility through the structured, coherent presentation of its academic, research and outreach activity, consolidating its capacity to disseminate its thought leadership systematically and make the FADE community's output accessible and understandable.

It also gives prioritised, organised attention to areas with strong potential for visibility, such as the Open Agenda, Science Communication, Research Projects and dissemination of Bachelor's Degree Final Projects/Master's Theses, improving institutional capacity to communicate publications, projects, events, results and initiatives coherently to different audiences.

In this way, it strengthens the School's identity and positioning while providing cross-cutting support to the other Strategic Plan projects, facilitating their visibility, coordination and recognition within a common communication framework.

At the same time, COMUNICA contributes significantly to Academic Excellence by organising and improving internal and external communication of relevant activities, opportunities, teaching content and academic results, and by incorporating an approach to recording and traceability that strengthens coordination, institutional coherence and management based on structured information.

The 35\% / 65\% distribution reflects the project's nature: a cross-cutting communication infrastructure clearly oriented towards institutional and social impact, but with significant effects also on the quality, coordination and accessibility of the School's academic life.


\newpage

\Needspace{18\baselineskip}
\section{Objectives, Action Plans and Development}
\Needspace{14\baselineskip}
\addsec{S.O.1. Structure FADE's communication system through common processes, flows and criteria}
{\small\color{mutedcolor}\itshape Develop a common framework to organise the collection, validation, prioritisation, scheduling and publication of content generated by the various agents linked to FADE.}\par\vspace{6pt}

\begin{tcolorbox}[pabox={comcolor}{A.P.1.1 \enspace\textbullet\enspace Definition of the Process, Flow and Criteria Model for the Communication System}]

\fieldlabel{Description}
{\small Definition of a common framework to organise FADE's communication system: processes, work flows, editorial criteria, content typologies and channel adaptation guidelines.}\par
\fieldsep
\fieldlabel{Deliverables}
{\small \begin{itemize}[noitemsep,topsep=2pt]
\item FADE communication system framework document
\item Basic process and flow map
\item Editorial and publication criteria
\item Classification of content typologies and channels
\end{itemize}}\par
\fieldsep
\fieldlabel{KPIs}
{\small \begin{itemize}[noitemsep,topsep=2pt]
\item Framework document prepared and validated (Yes/No)
\item Process and flow map defined (Yes/No)
\item Editorial criteria defined (Yes/No)
\item Classification of typologies and channels defined (Yes/No)
\end{itemize}}\par
\end{tcolorbox}
\begin{tcolorbox}[riskbox]
{\footnotesize\sffamily\bfseries\color{red!50!black} RISKS}\par\nopagebreak[4]
{\small \begin{itemize}[noitemsep,topsep=2pt]
\item • Design of an excessively complex or insufficiently operational model.
\item • Lack of adjustment to FADE's reality.
\item • Classification that is unclear or difficult to apply.
\end{itemize}}
\end{tcolorbox}
\begin{tcolorbox}[mitbox]
{\footnotesize\sffamily\bfseries\color{green!40!black} MITIGATION PLAN}\par\nopagebreak[4]
{\small \begin{itemize}[noitemsep,topsep=2pt]
\item • Design a simple, realistic model adjusted to FADE's management capacities.
\item • Discuss it with the main units involved.
\item • Prioritise a clear, functional and applicable classification.
\end{itemize}}
\end{tcolorbox}
\vspace{8pt}

\begin{tcolorbox}[pabox={comcolor}{A.P.1.2 \enspace\textbullet\enspace Implementation of the Operational Content Management and Publication Workflow}]

\fieldlabel{Description}
{\small Put into operation a consistent workflow for the reception, validation, prioritisation, scheduling and publication of content generated by the various agents linked to the School.}\par
\fieldsep
\fieldlabel{Deliverables}
{\small \begin{itemize}[noitemsep,topsep=2pt]
\item Operational workflow implemented
\item Basic operating and usage protocol
\item Initial set of managed content
\item Basic usage guide
\end{itemize}}\par
\fieldsep
\fieldlabel{KPIs}
{\small \begin{itemize}[noitemsep,topsep=2pt]
\item Operational workflow implemented (Yes/No)
\item Basic protocol approved (Yes/No)
\item No. of content items managed in initial phase
\item No. of units/persons responsible incorporated
\end{itemize}}\par
\end{tcolorbox}
\begin{tcolorbox}[riskbox]
{\footnotesize\sffamily\bfseries\color{red!50!black} RISKS}\par\nopagebreak[4]
{\small \begin{itemize}[noitemsep,topsep=2pt]
\item • Low adoption of the workflow by the units or people involved.
\item • Irregular or incomplete initial recording of content.
\item • Excessive operational complexity in workflow operation.
\end{itemize}}
\end{tcolorbox}
\begin{tcolorbox}[mitbox]
{\footnotesize\sffamily\bfseries\color{green!40!black} MITIGATION PLAN}\par\nopagebreak[4]
{\small \begin{itemize}[noitemsep,topsep=2pt]
\item • Implement a simple, clear workflow adjusted to FADE's actual needs.
\item • Define minimum responsibilities and easy-to-apply operational guidelines.
\item • Carry out an initial piloting or gradual deployment phase with a limited number of units before progressive extension.
\end{itemize}}
\end{tcolorbox}
\vspace{8pt}

\newpage
\Needspace{14\baselineskip}
\addsec{S.O.2. Organise, make visible and highlight the value of strategic content and the activity generated by FADE}
{\small\color{mutedcolor}\itshape Reinforce FADE's capacity to communicate its activity, giving priority attention to the Open Agenda, Science Communication, Research Projects and the Dissemination of Bachelor's Degree Final Projects/Master's Theses.}\par\vspace{6pt}

\begin{tcolorbox}[pabox={comcolor}{A.P.2.1 \enspace\textbullet\enspace Structuring and Visibility of FADE's Priority Areas and Strategic Content}]

\fieldlabel{Description}
{\small Identification, organisation and definition of visibility criteria for the priority areas (Open Agenda, Science Communication, Research Projects, Bachelor's Degree Final Projects/Master's Theses). Includes specific channels for making FADE's Thought Leadership visible.}\par
\fieldsep
\fieldlabel{Deliverables}
{\small \begin{itemize}[noitemsep,topsep=2pt]
\item Structuring and visibility framework document
\item Classification of priority areas and content typologies
\item Basic visibility criteria
\item Content--audience--channel linkage framework
\end{itemize}}\par
\fieldsep
\fieldlabel{KPIs}
{\small \begin{itemize}[noitemsep,topsep=2pt]
\item Framework document validated (Yes/No)
\item Areas and typologies defined (Yes/No)
\item Visibility criteria defined (Yes/No)
\item Linkage scheme defined (Yes/No)
\end{itemize}}\par
\end{tcolorbox}
\begin{tcolorbox}[riskbox]
{\footnotesize\sffamily\bfseries\color{red!50!black} RISKS}\par\nopagebreak[4]
{\small \begin{itemize}[noitemsep,topsep=2pt]
\item • Excessively broad or insufficiently operational definition of priority areas and strategic content.
\item • Overlap between content typologies or lack of clarity in their delimitation.
\item • Difficulty in reaching consensus on priorities and visibility criteria between different units or agents.
\end{itemize}}
\end{tcolorbox}
\begin{tcolorbox}[mitbox]
{\footnotesize\sffamily\bfseries\color{green!40!black} MITIGATION PLAN}\par\nopagebreak[4]
{\small \begin{itemize}[noitemsep,topsep=2pt]
\item • Prioritise a clear, limited and functional proposal of areas, content and visibility criteria.
\item • Define simple classification and communicative treatment criteria.
\item • Discuss the proposal with the main units involved before final validation.
\end{itemize}}
\end{tcolorbox}
\vspace{8pt}

\newpage
\Needspace{14\baselineskip}
\addsec{S.O.3. Reinforce the traceability, archiving and monitoring of FADE's communication activity}
{\small\color{mutedcolor}\itshape Implement a system to record, classify, archive and consult publications and communication actions, reinforcing traceability and institutional memory.}\par\vspace{6pt}

\begin{tcolorbox}[pabox={comcolor}{A.P.3.1 \enspace\textbullet\enspace Implementation of the Recording, Archiving and Monitoring System for Publications and Communication Actions}]

\fieldlabel{Description}
{\small Structured system for recording and archiving publications and communication actions, incorporating partial structural alignment with the FADE Impact Lab agent categories to favour institutional coherence.}\par
\fieldsep
\fieldlabel{Deliverables}
{\small \begin{itemize}[noitemsep,topsep=2pt]
\item Basic recording and monitoring system implemented
\item Basic recording and consultation protocol
\item Category and field structure defined (with references to FADE Impact Lab)
\item Initial set of publications recorded
\end{itemize}}\par
\fieldsep
\fieldlabel{KPIs}
{\small \begin{itemize}[noitemsep,topsep=2pt]
\item Basic system implemented (Yes/No)
\item Basic protocol approved (Yes/No)
\item Category structure defined (Yes/No)
\item No. of publications and actions recorded in initial phase
\end{itemize}}\par
\end{tcolorbox}
\begin{tcolorbox}[riskbox]
{\footnotesize\sffamily\bfseries\color{red!50!black} RISKS}\par\nopagebreak[4]
{\small \begin{itemize}[noitemsep,topsep=2pt]
\item • Design of an excessively complex or difficult-to-maintain recording system.
\item • Lack of homogeneity in the classification and archiving of publications and communication actions.
\item • Mismatch between COMUNICA's own categories and the references shared with other strategic projects, especially FADE Impact Lab.
\end{itemize}}
\end{tcolorbox}
\begin{tcolorbox}[mitbox]
{\footnotesize\sffamily\bfseries\color{green!40!black} MITIGATION PLAN}\par\nopagebreak[4]
{\small \begin{itemize}[noitemsep,topsep=2pt]
\item • Design a simple, functional system adjusted to FADE's real management and consultation needs.
\item • Define clear classification, archiving and recording criteria from the start.
\item • Incorporate only partial structural alignment with other projects, limited to common categories and references that add real value and avoid unnecessary duplication or rigidities.
\end{itemize}}
\end{tcolorbox}
\vspace{8pt}

\Needspace{18\baselineskip}
\section{Timeline}

\small

\begin{tabularx}{\textwidth}{>{\bfseries}l *{8}{>{\centering\arraybackslash}p{1.3cm}}}
\toprule
A.P. & 26/27 S1 & 26/27 S2 & 27/28 S1 & 27/28 S2 & 28/29 S1 & 28/29 S2 & 29/30 S1 & 29/30 S2 \\
\midrule
A.P.1.1 & $\blacksquare$ & $\blacksquare$ &  &  &  &  &  &  \\
A.P.1.2 &  & $\blacksquare$ & $\blacksquare$ & $\blacksquare$ &  &  &  &  \\
A.P.2.1 &  & $\blacksquare$ & $\blacksquare$ & $\blacksquare$ & $\blacksquare$ & $\blacksquare$ & $\blacksquare$ & $\blacksquare$ \\
A.P.3.1 &  &  & $\blacksquare$ & $\blacksquare$ & $\blacksquare$ & $\blacksquare$ & $\blacksquare$ & $\blacksquare$ \\
\bottomrule
\end{tabularx}

\Needspace{18\baselineskip}
\section{Summary of Objective Achievement}

\small
Achievement of the project's objectives will be evaluated through the monitoring of indicators defined in the action plans, as well as an overall assessment of the degree of structuring, visibility and traceability of the School's communication system.

\textbf{S.O.1. Structure FADE's communication system through common processes, flows and criteria.} It will be considered achieved when FADE has a defined, operational communication system, with clearly established procedures, workflows, editorial criteria and guidelines for adapting content to channels, enabling consistent management of the collection, initial recording, validation, prioritisation, scheduling and publication of content generated by the various actors linked to the School.

\textbf{S.O.2. Organise, make visible and highlight the value of strategic content and the activity generated by FADE.} It will be considered achieved when FADE has a structured framework for coherently identifying, classifying and making visible its strategic content, giving priority attention to the Open Agenda, Science Communication, Research Projects and dissemination of Bachelor's Degree Final Projects/Master's Theses, as well as communication support for the other Strategic Plan projects, so that its academic, research, outreach and institutional activity becomes more accessible, understandable and suitable for wider dissemination.

\textbf{S.O.3. Reinforce the traceability, archiving and monitoring of FADE's communication activity.} It will be considered achieved when the School has a structured system for recording, archiving and consulting publications and communication activities, enabling the disseminated activity to be traced, strengthening institutional memory and facilitating monitoring and continuous improvement of the communication system, incorporating, where relevant, common references with other strategic projects, particularly the ecosystem of actors defined in FADE Impact Lab.



\end{document}